\documentclass[11pt,a4paper]{article}
\usepackage[utf8]{inputenc}
\usepackage[T1]{fontenc}
\usepackage[french]{babel}
\usepackage{lmodern}
\usepackage{amsmath,amssymb,amsthm,mathtools}
\usepackage{physics}
\usepackage{siunitx}
\usepackage{bm}
\usepackage{geometry}
\usepackage{booktabs}
\usepackage{array}
\usepackage{enumitem}
\usepackage{hyperref}
\usepackage{microtype}
\geometry{margin=2.2cm}
\hypersetup{colorlinks=true,linkcolor=blue,citecolor=blue,urlcolor=blue}
\setlength{\parskip}{0.5em}
\setlength{\parindent}{0pt}

\title{Méthodes SRC/MPCD quasi-incompressibles :\\ projection Q6/Q9, resampling pondéré, conditions limites virtuelles et fermeture virielle}
\author{}
\date{\today}

\begin{document}
\maketitle



\begin{abstract}
Ce rapport rassemble l'état de développement de la méthode SRC/MPCD quasi-incompressible en distinguant explicitement deux générations de code. La première, historique, reposait sur une redistribution particulaire dure vers une cible d'occupation et sur une fermeture liquide virielle. Elle reste utile pour l'interprétation physique, mais ses résultats quantitatifs doivent être relus avec prudence car une partie des validations a été effectuée avant la correction du thermostat, avant la stabilisation des conditions limites par particules virtuelles et avant les premières vectorisations MATLAB. La seconde génération, désormais considérée comme la chaîne de référence, est la méthode Q6/Q9 : collision SRC/MPCD classique, projection de vitesse incompressible Q6, correction basse fréquence du flux de masse Q9, thermostat cellulaire corrigé et diagnostics conservatifs associés.

La version actuelle, sur la branche \texttt{q9-piston-case-virial-kick}, valide quatre points principaux. Premièrement, les cas Taylor--Green et Poiseuille montrent que Q6/Q9 réduit les modes compressifs de grande longueur d'onde sans détruire les structures hydrodynamiques moyennes. Deuxièmement, les campagnes Poiseuille multi-résolutions ont mis en évidence une dépendance apparente de la viscosité identifiée au nombre de mailles transverse ; une correction géométrique en \(N_y^2\) ramène les viscosités corrigées vers la référence bulk Taylor--Green, ce qui impose de distinguer viscosité brute de fit et viscosité corrigée par la hauteur discrète du canal. Troisièmement, le piston staircase avec fermeture virielle bulk fournit une équation d'état contrôlée, avec \(P_{\mathrm{vir}}=K_{\mathrm{virial}}(\rho-\rho_{\mathrm{EOSRef}})\), \(P_{\mathrm{tot}}=P_{\mathrm{kin}}+P_{\mathrm{vir}}\), \(P_{\mathrm{drive}}=P_{\mathrm{tot}}\) et des fits staircase de pente quasi parfaite sur les grandeurs bulk. Quatrièmement, le diagnostic mécanique de paroi avec particules virtuelles \emph{wallVP} est maintenant séparé en impulsions signées, pressions positives de compression et moyennes cumulées ; le test d'équilibre gaz entre deux murs valide l'ordre de grandeur de la pression mécanique cumulée, tandis que le staircase piston demande encore des fenêtres longues pour réduire la variance de l'estimateur impulsionnel.



Enfin, le chantier CUDA-TOPO ajoute une brique de pénalisation Darcy--Brinkman pilotée par un champ externe $\chi$. Ce champ permet de représenter des géométries complexes et de construire des proxies de coût, de portance et de traînée pour des canaux, obstacles et profils NACA. Le point d'étape validé est toutefois nuancé : Darcy--Brinkman est robuste comme pénalisation poreuse volumique, mais ne reproduit pas à lui seul une frontière solide impenetrable. La comparaison Von Karman montre qu'un cylindre Darcy peut annuler la vitesse dans la zone pénalisée sans produire le déficit de vitesse et le sillage du solide immergé codé ; un futur chantier devra donc dériver une vraie condition solide à partir de $\chi$.

Les développements 0418--0426 révisent ce statut. Le champ $\chi$ n'est plus seulement utilisé comme frein Darcy pur : il pilote désormais une chaîne plus complète combinant désactivation initiale des particules dans les zones $\chi\simeq0$, kick moyen Darcy--Brinkman, bain de réémission orienté par la normale $\nabla\chi/\|\nabla\chi\|$, et contribution de particules virtuelles de collision dérivées de $\chi$. Cette variante, notée dans les scripts \texttt{mean\_outward\_bath+chiVP}, permet d'utiliser $\chi$ comme représentation géométrique souple de formes solides quelconques, tout en conservant la distinction conceptuelle entre pénalisation poreuse et condition pariétale exacte. La correction 0426 des \emph{fastflags} CUDA résidents est capitale : elle fait repasser le cas Darcy/chiVP segmenté de l'ordre de $0.19$--$0.23\,\si{\second}$ par pas à environ $0.056\,\si{\second}$ par pas sur le backward step $960\times480$, soit un coût voisin du chemin solide spécialisé. Le goulet mesuré ne venait donc pas de Darcy/chiVP mais d'un aiguillage CUDA incomplet dans les scripts de démonstration.


Depuis les développements C++/OpenMP réalisés ensuite sur \texttt{SRC\_incompressible\_openMP}, une nouvelle étape peut être isolée. Les conditions limites inlet/outlet Q6/Q9 ont été combinées avec des solides immergés fixes à l'aide d'un masque cellule/face : les cellules fluides adjacentes au solide restent actives, les faces fluide--solide sont fermées, le halo Q9 autour du solide est désactivé comme stratégie nominale, et les flux normaux Q6/Q9 au travers du solide sont imposés nuls après fermeture géométrique. Les validations périodiques et inlet/outlet avec cylindre fixe montrent que cette formulation solide immergé est suffisamment robuste pour être considérée comme validée. En revanche, les tests ultérieurs sur poches de cellules pauvres ont mis en évidence une limite indépendante des solides : Q9, ses sécurités low-mass et les tentatives de garde virielle locale peuvent transformer une cellule pauvre en défaut eulérien actif, alors que le transport SRC classique tend plutôt à la convecter et la diffuser.

La réponse méthodologique récente est une troisième génération de prototype MATLAB, séparée de la branche solides : une formulation SRC/MPCD pondérée avec resampling local conservatif. Les particules portent désormais une masse variable $m_i$ ; les champs de cellule sont reconstruits par dépôt massique ; un pool préalloué, des rôles de particules \emph{fluide}/\emph{latente}/\emph{inactive}, un masque de cellules mouillées et une boucle extraction--insertion--remap maintiennent une population locale suffisante sans transformer les cellules pauvres en puits de pression. La masse de cellule, le moment et, lorsque nécessaire, l'énergie thermique locale sont restaurés par des corrections conservatrices. Les premiers tests démontrent la fermeture de la boucle sur poches vide/surpeuplée, sur hétérogénéités distribuées, sur remplissage d'un domaine initialement latent et, surtout, sur Taylor--Green forcé et Poiseuille avec particules virtuelles de paroi. Le cas Poiseuille wallVP indique que Q6+resampling conserve une viscosité effective proche de la référence SRC wallVP pure tout en contrôlant la masse cellulaire au niveau du roundoff.


Depuis le portage C++/OpenMP réalisé dans \texttt{SRC\_incompressible\_openMP}, cette stratégie n'est plus seulement un prototype MATLAB. L'infrastructure particulaire pondérée, les rôles \texttt{Fluid}/\texttt{Inactive}/\texttt{Latent}, les fichiers d'état \texttt{.smpcd} enrichis, le dépôt pondéré, la collision SRC pondérée, Q6 sur champs pondérés, le pool de recyclage, le remap local, la renormalisation thermique, le \texttt{mass guard} et, surtout, le garde de population $N_{\min}/N_{\mathrm{target}}/N_{\max}$ ont été traduits dans le code OpenMP. Les validations récentes montrent que la chaîne OpenMP conserve la dynamique Taylor--Green et Poiseuille, impose la masse cellulaire au roundoff dans les cas resamplés, accélère fortement le remplissage d'un domaine initialement inactif et contrôle le support particulaire dans les cellules mouillées. La différence essentielle par rapport au premier portage OpenMP est que le resampling est maintenant redevenu \emph{population-driven}, comme dans MATLAB : on corrige d'abord le nombre de particules actives par cellule, puis seulement ensuite la masse et l'énergie locales.

Le document corrige donc le rapport précédent en reclassant les résultats historiques, en documentant le thermostat corrigé, les conditions limites virtuelles, le solveur \texttt{general\_bc} avec fixation du mode constant, les diagnostics de conservation du moment, les corrections de viscosité liées à la discrétisation transverse, les premières stratégies de parallélisation/vectorisation MATLAB et la nouvelle stratégie de resampling pondéré. Le portage C++/OpenMP reste une étape naturelle, mais le socle MATLAB fournit maintenant un protocole de validation plus fiable : Taylor--Green pour la préservation des structures et la viscosité bulk de référence, Poiseuille pour la viscosité effective corrigée et les parois, piston staircase pour l'EOS bulk, gaz à l'équilibre entre murs pour la pression mécanique wallVP, et tests de population contrôlée pour la robustesse locale du support particulaire.


Le dernier cycle de développement OpenMP ajoute quatre éléments qui modifient fortement l'architecture pratique du code. Premièrement, les conditions limites ouvertes ne sont plus seulement définies par face entière : des segments \texttt{inlet}/\texttt{outlet} multiples, en coordonnées relatives, peuvent être placés sur une même face solide, les portions non couvertes restant solides. Deuxièmement, l'ancien aiguillage \texttt{method=classic/q6} a été supprimé au profit de deux interrupteurs explicites, \texttt{projectionEnable} et \texttt{resamplingEnable}, qui définissent les quatre cas d'ablation \texttt{classic}, \texttt{classic\_ resampling}, \texttt{q6} et \texttt{q6\_ resampling}. Troisièmement, un mécanisme de réponse de capacité fermée a été introduit : lorsque la masse contenue dans un domaine fluide fermé dépasse la masse nominale, l'intensité effective de Q6 est diminuée par modulation de \texttt{q6ProjectionStrength}, le module viriel effectif augmente, et le remap de masse suit une cible cellulaire effective qui conserve la surmasse globale au lieu de l'effacer. Quatrièmement, la pression bulk reconstruite, $P_{\mathrm{kin}}+P_{\mathrm{vir}}$, est maintenant projetée en charge normale sur les portions solides des parois ; le test statique de boîte uniformément surcompressée vérifie que les quatre faces reçoivent la même pression et que la force pariétale nette est nulle à la précision numérique.

Enfin, la branche \texttt{SRC\_GPU} ajoute un jalon d'architecture parallèle : le SRC classic, défini ici comme advection/streaming, décalage de grille, collision/rotation SRC et thermostat, dispose désormais d'un chemin CUDA résident validé sur les conditions limites périodiques, wall-simple, solide rectangle, piston/mobile wall, inlet/outlet full-face et inlet/outlet segmenté. La fermeture liquide Q6+resampling+viriel reste un module séparé, conservé pour une migration CUDA ultérieure.



Le cycle CUDA 0294--0303 complète ce statut : un module de resampling CUDA post-SRC, non destructif et conservatif, a été ajouté au chemin résident. Il comprend un survey passif du support particulaire, un reconditionnement conservatif des masses, un garde de population local par split/merge de particules représentatives, une restauration de l'énergie cinétique relative et un filtrage boundary-aware. Les validations longues sur marche arrière montrent que le réglage nominal provisoire $N_{\min}=12$, $N_{\mathrm{target}}=20$, $N_{\mathrm{max}}=32$ supprime les cellules humides vides derrière la marche pour $U_{\mathrm{in}}=0.30$, $0.45$ et $0.60$, tout en conservant masse, impulsion et énergie relative au bruit numérique.



Le chantier 0490a--0490p ajoute enfin un traitement multi-espèces au resampling pondéré. Le champ \texttt{type} devient une identité physique transportée, distincte du \texttt{role} algorithmique \texttt{Fluid}/\texttt{Inactive}/\texttt{Latent}. Les dépôts par cellule et par espèce, la fermeture de masse, le garde de population, le remplissage des cellules temporairement vides et les transferts donneur--receveur sont désormais conservatifs par espèce. Le chemin validé 0490p garde ces opérations sur l'état CUDA partagé, y compris la politique wet/poor/rich et le pool de slots, et termine l'audit avec \texttt{remaining\_cpu\_scope=none} dans le sous-ensemble périodique actuellement qualifié.

Le cycle 0491 complète cette infrastructure par un Q6 multi-espèces CUDA résident. Une seule projection elliptique reste appliquée au mouvement barycentrique ; sa correction finale est ensuite distribuée selon les fractions massiques et le coefficient \texttt{q6StrengthDeclared} de chaque espèce, avec une identité de recomposition qui conserve exactement la correction barycentrique. Le mode commun reproduit le Q6 historique, tandis que le mode pondéré rend les vitesses relatives sensibles au \texttt{type} sans modifier masse, population ou identité particulaire. Les validations bulk sont encourageantes ; le resampling demeure cependant optionnel en raison de son coût et des questions encore ouvertes sur l'autodiffusion.

\end{abstract}


\tableofcontents


\section*{Note de version et statut des résultats}
\addcontentsline{toc}{section}{Note de version et statut des résultats}

Ce document complète une version antérieure du rapport. Les sections consacrées à la redistribution dure, au zonage et aux premières validations Q6/Q9 sont conservées parce qu'elles expliquent la construction progressive de la méthode. Elles ne doivent toutefois plus être lues comme des validations quantitatives définitives lorsque les résultats provenaient des scripts historiques : ces scripts ne contenaient pas encore le thermostat cellulaire corrigé, n'utilisaient pas les conditions limites par particules virtuelles dans les canaux, et n'étaient pas encore structurés pour les premiers gains de vectorisation/parallélisation MATLAB.



Depuis les patchs post-0143, le code OpenMP utilise en outre une configuration plus explicite : \texttt{projectionEnable} contrôle seul Q6, \texttt{resamplingEnable} contrôle seul le garde de population et le remap pondéré, et les ouvertures partielles sont décrites par des segments relatifs. Les développements de capacité fermée, de viriel faiblement compressible et de charge pariétale ne remplacent pas les validations précédentes ; ils constituent un enrichissement activé par paramètres, conçu pour les domaines fermés sur-remplis et pour le futur couplage à des parois solides déformables.


Depuis le cycle 0294--0303, la branche \texttt{SRC\_GPU} inclut en outre un resampling CUDA post-SRC destiné au contrôle du support particulaire. Ce module est volontairement séparé de Q6 CUDA : Q6 CPU/OpenMP reste réactivable dans l'architecture, mais la validation initiale du resampling CUDA concerne le chemin SRC classic résident suivi d'une opération de re-échantillonnage conservatif. Les listings de paramètres 0292 joints au présent rapport restent la référence pour le socle CUDA classic, inlet/outlet, thermostat et fermeture liquide OpenMP ; les options 0295--0303 sont documentées ci-dessous comme complément spécifique au resampling CUDA.


Le cycle 0490a--0490p étend ce resampling aux mélanges de plusieurs espèces enregistrées. La conservation n'est plus contrôlée seulement sur la masse totale : les opérations discrètes, le remplissage, la fermeture de masse et les transferts sont audités séparément pour chaque valeur physique de \texttt{type}. Le checkpoint de référence de cette section est le commit \texttt{36abd23} de la branche \texttt{surf}. Les inventaires CSV joints ont été reconstruits à partir des appels \texttt{getenv} et du parseur \texttt{params.kv} de cet état du code ; ils remplacent les inventaires consolidés 0436b pour les options ajoutées ou omises depuis cette date.

Le cycle 0491 rend ensuite l'application Q6 sensible aux espèces sans créer de solveur elliptique par phase. La section 33 distingue désormais explicitement la formulation 0490 du resampling et la formulation 0491 de Q6, documente leurs paramètres, leur ordre d'exécution, leurs invariants et leur niveau de qualification.

La chaîne de référence actuelle est celle de la branche \texttt{q9-piston-case-virial-kick}. Sauf indication contraire, on appelle désormais \textbf{Q6/Q9 complet} la séquence suivante : étape SRC/MPCD classique, projection de vitesse Q6, projection basse fréquence du flux de masse Q9, thermostat corrigé, correction globale exacte du moment pour les kicks additionnels, puis diagnostics bulk/paroi. Les résultats historiques restent utiles pour comparer les idées, mais les chiffres de calibration doivent être repris avec les scripts corrigés.

\section{Architecture générale de la méthode}

Le code simule un fluide particulaire en deux dimensions dans un domaine rectangulaire de taille $L_x \times L_y$, discrétisé en un maillage cartésien régulier de $N_x \times N_y$ cellules de taille
\begin{equation}
 a_0 = \Delta x = \frac{L_x}{N_x},
 \qquad
 \Delta y = \frac{L_y}{N_y},
 \qquad
 V_c = \Delta x\,\Delta y,
\end{equation}
où $V_c$ désigne l'aire d'une cellule dans la simulation bidimensionnelle.

Les variables particulaires sont la position $\bm{x}_i \in \mathbb{R}^2$ et la vitesse $\bm{v}_i \in \mathbb{R}^2$ de la particule $i$, pour $i=1,\dots,N_p$, où $N_p$ est le nombre total de particules. Le pas de temps est noté $\Delta t$.

La séquence temporelle de référence pour les validations récentes n'est plus la chaîne historique avec redistribution dure. Elle s'écrit maintenant, au niveau d'un pas de temps,
\begin{enumerate}[label=\arabic*)]
  \item advection particulaire avec les forces imposées et application des conditions limites géométriques ;
  \item collision SRC/MPCD avec décalage aléatoire de grille, rotation SRD et thermostat cellulaire corrigé ;
  \item reconstruction particules--grille des champs d'occupation, de vitesse moyenne, de température et de pression cinétique ;
  \item projection de vitesse Q6, c'est-à-dire correction du champ moyen pour réduire \(\nabla_h\cdot u\) ;
  \item projection basse fréquence du flux de masse Q9, c'est-à-dire relaxation de \(\nabla_h\cdot(Nu)\) sur les modes cohérents de grande longueur d'onde ;
  \item interpolation des corrections grille vers les particules avec correction globale exacte de la quantité de mouvement lorsque des kicks non strictement conservatifs sont appliqués ;
  \item diagnostics de densité, divergence, flux de masse, température, pression bulk, pression mécanique de paroi et, selon le cas, vorticité ou profils moyens.
\end{enumerate}

La chaîne historique, fondée sur redistribution particulaire puis fermeture liquide par pression virielle, est décrite plus loin parce qu'elle a joué un rôle important dans la construction de la méthode. Elle n'est toutefois plus le chemin principal pour les validations Q6/Q9 actuelles. Dans les cas piston récents, le terme viriel n'est pas utilisé pour déplacer durement les particules vers une cible de densité ; il sert à construire une fermeture bulk explicite,
\begin{equation}
  P_{\mathrm{vir}} = K_{\mathrm{virial}}(\rho-\rho_{\mathrm{EOSRef}}),
  \qquad
  P_{\mathrm{tot}}=P_{\mathrm{kin}}+P_{\mathrm{vir}},
  \qquad
  P_{\mathrm{drive}}=P_{\mathrm{tot}}.
\end{equation}

On distinguera donc dans la suite :
\begin{itemize}
  \item l'\textbf{état MPCD classique}, obtenu après transport, conditions limites, collision et thermostat ;
  \item l'\textbf{état Q6}, obtenu après projection de vitesse ;
  \item l'\textbf{état Q9}, obtenu après relaxation basse fréquence du flux de masse ;
  \item l'\textbf{état piston/viriel}, dans lequel une pression bulk additionnelle est reconstruite pour tester une réponse EOS quasi incompressible ;
  \item les \textbf{diagnostics de paroi}, qui mesurent un flux d'impulsion et ne doivent pas être confondus avec la pression thermodynamique reconstruite dans le bulk.
\end{itemize}

\section{Méthode MPCD : principes, advection, rotation, décalage de grille}

\subsection{Principe général}

La méthode MPCD représente un fluide par un ensemble de particules ponctuelles qui alternent deux sous-étapes : une sous-étape de transport balistique (ou advection libre) et une sous-étape de collision collective, dans laquelle les vitesses sont mélangées cellule par cellule.

La variable conservative fondamentale au niveau microscopique est la quantité de mouvement particulaire. La dissipation, la diffusion de quantité de mouvement et la température émergent de l'alternance entre transport et collisions.

\subsection{Advection avec forçage volumique}

À l'instant $t^n = n\Delta t$, la vitesse de la particule $i$ est d'abord modifiée par les forces volumiques imposées. Le code utilise actuellement une force volumique horizontale $f_x$ et une accélération verticale uniforme $g$. L'étape d'advection s'écrit donc
\begin{equation}
 v_{i,x}^{\star} = v_{i,x}^{n} + f_x\,\Delta t,
 \qquad
 v_{i,y}^{\star} = v_{i,y}^{n} + g\,\Delta t,
\end{equation}
puis
\begin{equation}
 \bm{x}_i^{\star} = \bm{x}_i^{n} + \Delta t\,\bm{v}_i^{\star}.
\end{equation}

Ici $v_{i,x}$ et $v_{i,y}$ sont les composantes de la vitesse $\bm{v}_i$, et l'exposant $\star$ désigne l'état pré-collision après transport et forçage.

\subsection{Découpage en cellules et décalage aléatoire de grille}

L'étape de collision est effectuée dans des cellules cartésiennes de taille $a_0$. Pour limiter les artefacts liés à la grille et restaurer l'invariance galiléenne de la méthode, le code applique à chaque pas un décalage aléatoire uniforme
\begin{equation}
 \bm{s} \in \left[-\frac{a_0}{2},\frac{a_0}{2}\right]^2,
\end{equation}
puis assigne les particules aux cellules de la grille décalée. Si la cellule d'indice $c$ contient $N_c$ particules, la vitesse moyenne cellulaire est
\begin{equation}
 \bm{u}_c = \frac{1}{N_c}\sum_{i\in c} \bm{v}_i.
\end{equation}

Le code distingue proprement les axes périodiques et non périodiques au moment de cette affectation, ce qui est essentiel lorsque les bords physiques ou un piston mobile rendent le domaine non périodique dans une direction.

\subsection{Collision SRD par rotation des vitesses relatives}

Dans la collision SRD, on décompose la vitesse particulaire sous la forme
\begin{equation}
 \bm{v}_i = \bm{u}_c + \bm{w}_i,
 \qquad i\in c,
\end{equation}
où $\bm{w}_i$ est la fluctuation de vitesse de la particule $i$ relativement à la vitesse moyenne de sa cellule.

Ces fluctuations sont ensuite tournées d'un angle $\alpha$ ou $-\alpha$ aléatoire, identique pour toutes les particules d'une même cellule. Si $\mathbf{R}(\theta_c)$ désigne la matrice de rotation d'angle $\theta_c\in\{+\alpha,-\alpha\}$ dans la cellule $c$, la collision s'écrit
\begin{equation}
 \bm{w}_i' = \mathbf{R}(\theta_c)\,\bm{w}_i,
 \qquad
 \bm{v}_i' = \bm{u}_c + \bm{w}_i'.
\end{equation}

Cette étape conserve exactement la quantité de mouvement de chaque cellule, puisque
\begin{equation}
 \sum_{i\in c} \bm{w}_i = \bm{0}
 \quad \Longrightarrow \quad
 \sum_{i\in c} \bm{v}_i' = \sum_{i\in c} \bm{v}_i.
\end{equation}

\section{Liens entre unités MPCD et grandeurs physiques}

\subsection{Choix d'unités réduites}

Le code est écrit en unités réduites MPCD. Les grandeurs indépendantes choisies comme unités de base sont généralement :
\begin{itemize}
  \item la longueur de maille $a_0$ ;
  \item le pas de temps $\Delta t$ ;
  \item la masse particulaire $m$.
\end{itemize}

À partir de ces trois échelles, on déduit :
\begin{equation}
 u_0 = \frac{a_0}{\Delta t}
\end{equation}
comme échelle de vitesse,
\begin{equation}
 \varepsilon_0 = m u_0^2 = m\frac{a_0^2}{\Delta t^2}
\end{equation}
comme échelle d'énergie, et
\begin{equation}
 \nu_0 = \frac{a_0^2}{\Delta t}
\end{equation}
comme échelle de viscosité cinématique.

Dans le code actuel, les diagnostics et la fermeture liquide utilisent la variable nommée \verb|rho| pour désigner en pratique une \textbf{densité de nombre surfacique}
\begin{equation}
 n = \frac{N_c}{V_c},
\end{equation}
et non une masse volumique au sens strict. La masse volumique physique en 2D s'obtiendrait par
\begin{equation}
 \rho_m = m n.
\end{equation}

Cette distinction est importante : les quantités reconstruites et les pressions utilisées dans la fermeture s'expriment dans le code à partir de $n$, avec le choix implicite $m=1$ pour les diagnostics de fermeture.

\subsection{Température et énergie thermique}

La quantité reconstruite dans le code est un champ local que l'on peut noter $k_B T_c$, défini à partir de l'énergie cinétique des fluctuations autour de la vitesse moyenne de cellule. Si $d=2$ désigne la dimension spatiale et si $N_c\ge 2$, alors
\begin{equation}
 k_B T_c = \frac{1}{d(N_c-1)} \sum_{i\in c} m\,\abs{\bm{v}_i-\bm{u}_c}^2.
\end{equation}

Dans l'implémentation courante, $m=1$ et la grandeur stockée sous le nom \verb|kBT| ou \verb|T| correspond directement à $k_B T_c$ dans les unités réduites choisies.

\subsection{Pression cinétique}

La pression cinétique locale reconstruite en post-traitement et utilisée dans la fermeture liquide est
\begin{equation}
 P_{\mathrm{kin},c} = n_c\,k_B T_c,
\end{equation}
où $n_c=N_c/V_c$ est la densité de nombre de la cellule $c$.

Cette pression est une pression locale de type gaz parfait dans les unités réduites du modèle. Elle ne doit pas être confondue avec la pression de paroi obtenue par flux d'impulsion aux frontières.

\subsection{Viscosité et nombre de Reynolds}

La viscosité cinématique $\nu$ peut être exprimée sous forme réduite, puis convertie physiquement par
\begin{equation}
 \nu_{\mathrm{phys}} = \nu^*\,\frac{a_0^2}{\Delta t}.
\end{equation}

Le nombre de Reynolds associé à une vitesse caractéristique $U$, une longueur caractéristique $L$ et une viscosité cinématique $\nu$ vaut
\begin{equation}
 \mathrm{Re} = \frac{U L}{\nu}.
\end{equation}

Dans l'état actuel du code, la viscosité n'est pas seulement pensée comme une grandeur théorique MPCD ; elle est aussi \textbf{identifiée numériquement} via des diagnostics de type Poiseuille. Lorsque l'écoulement est piloté par une force volumique horizontale $f_x$, le profil parabolique de référence s'écrit
\begin{equation}
 u_x(y) = \frac{f_x}{2\nu}\,y(H-y),
\end{equation}
où $H=L_y$ est la hauteur du canal. Le simulateur en déduit une estimation de $\nu$ soit par courbure du profil, soit par les contraintes de paroi. Cette stratégie est particulièrement adaptée lorsque la redistribution et la fermeture liquide modifient l'écoulement au-delà du cas MPCD théorique pur.

\section{Thermostat MPCD}

Après la rotation SRD des vitesses relatives, le code applique un thermostat cellulaire optionnel. Son principe est de conserver la vitesse moyenne de cellule tout en renormalisant la norme quadratique des fluctuations pour imposer une température cible $k_B T$.

Pour une cellule $c$, on calcule d'abord l'énergie cinétique relative après rotation,
\begin{equation}
 E_{\mathrm{rel},c}' = \sum_{i\in c} \abs{\bm{w}_i'}^2,
\end{equation}
puis le nombre effectif de degrés de liberté
\begin{equation}
 \mathrm{dof}_c = d(N_c-1).
\end{equation}

Le thermostat introduit un facteur d'échelle
\begin{equation}
 \lambda_c = \sqrt{\frac{\mathrm{dof}_c\,k_B T}{E_{\mathrm{rel},c}'}}
\end{equation}
quand $N_c>1$ et $E_{\mathrm{rel},c}'>0$, puis remplace
\begin{equation}
 \bm{w}_i' \leftarrow \lambda_c\,\bm{w}_i'.
\end{equation}

L'état final après collision et thermostat est donc
\begin{equation}
 \bm{v}_i^{n+1/2} = \bm{u}_c + \lambda_c\,\mathbf{R}(\theta_c)\,\bigl(\bm{v}_i^{\star}-\bm{u}_c\bigr).
\end{equation}

Le thermostat conserve la quantité de mouvement de chaque cellule, mais ne conserve pas son énergie cinétique totale ; il impose une énergie thermique relative cible. Dans le code actuel, il constitue le principal mécanisme explicite de contrôle de la température au niveau MPCD.


\subsection{Correction du bug historique et thermostat vectorisé}

Les premières campagnes Q6/Q9 avaient été réalisées avec une version historique du thermostat. La correction récente a deux conséquences importantes pour l'interprétation des résultats. D'une part, le diagnostic \(k_BT_c\) doit être calculé à partir des fluctuations relatives autour de la vitesse moyenne cellulaire, avec \(d(N_c-1)\) degrés de liberté effectifs lorsque \(N_c>1\). D'autre part, la renormalisation doit être appliquée cellule par cellule en conservant exactement la vitesse moyenne de collision. Une erreur dans cette étape peut modifier la viscosité effective, la pression cinétique reconstruite et la comparaison entre cas classic, Q6 et Q9.

La version corrigée est aussi plus favorable à MATLAB : les sommes par cellule, les énergies relatives et les facteurs \(\lambda_c\) sont calculés par opérations groupées de type \texttt{accumarray}/indexation vectorisée lorsque cela est possible. Cette vectorisation ne change pas la méthode mathématique, mais elle réduit le coût de la boucle de collision et rend les longues validations piston ou Poiseuille plus praticables. Les résultats historiques antérieurs à cette correction doivent donc être conservés comme résultats de développement, non comme valeurs finales de calibration.

\section{Conditions limites}

Le code gère des conditions limites différentes sur chacun des quatre côtés du domaine. Les modes disponibles peuvent être résumés physiquement ainsi.

\subsection{Périodicité}

Quand un bord est périodique, une particule qui sort par ce bord rentre par le bord opposé avec la même vitesse. Si l'axe $x$ est périodique, cela se traduit par
\begin{equation}
 x_i \leftarrow x_i \pm L_x,
\end{equation}
sans modification de $\bm{v}_i$.

\subsection{Réflexion spéculaire}

Sur un bord solide spéculaire, la composante normale de la vitesse change de signe, tandis que la composante tangentielle est inchangée. Par exemple sur la paroi basse $y=0$,
\begin{equation}
 v_{i,y} \leftarrow -v_{i,y},
\qquad
 v_{i,x} \leftarrow v_{i,x}.
\end{equation}

Ce choix réalise une paroi sans pénétration et sans friction tangentielle.

\subsection{Rebond intégral (\emph{bounce-back})}

Dans le mode \emph{bounce-back}, la composante tangentielle est aussi réfléchie relativement à une vitesse de paroi imposée. Si la paroi inférieure se déplace à la vitesse tangentielle $U_{\mathrm{bottom}}$, on écrit
\begin{equation}
 v_{i,x} \leftarrow 2U_{\mathrm{bottom}} - v_{i,x},
\qquad
 v_{i,y} \leftarrow -v_{i,y}.
\end{equation}

Ce mode mime un non-glissement plus fort qu'une réflexion spéculaire pure.

\subsection{Thermalisation aux parois}

Dans le mode \emph{thermalize}, la composante normale sortante est remplacée par une variable aléatoire orientée vers l'intérieur, et la composante tangentielle est tirée autour de la vitesse de paroi. Pour la paroi inférieure, on a schématiquement
\begin{equation}
 v_{i,y} \sim \abs{\mathcal{N}(0,\sigma_w^2)},
\qquad
 v_{i,x} \sim U_{\mathrm{bottom}} + \mathcal{N}(0,\sigma_w^2),
\end{equation}
où $\sigma_w$ est l'écart-type thermique imposé par la paroi.

Ce mécanisme agit comme un thermostat de bord et injecte ou retire de l'énergie selon l'état incident du fluide.

\subsection{Paroi piston}

Le mode \emph{piston} impose sur le bord supérieur une position $y_{\mathrm{p}}(t)$ et une vitesse normale
\begin{equation}
 U_p(t) = \dot y_p(t).
\end{equation}
Lors d'un impact sur le piston,
\begin{equation}
 v_{i,y} \leftarrow 2U_p - v_{i,y}.
\end{equation}

Le piston réduit la \emph{hauteur active} du domaine,
\begin{equation}
 H_{\mathrm{actif}}(t)=y_{\mathrm{p}}(t)-m_p,
\end{equation}
où $m_p$ est une marge géométrique optionnelle (\verb|pistonActiveMargin|) destinée à définir la zone effectivement disponible sous le piston. Le maillage reste fixe : ce n'est pas la taille des cellules qui change, mais le nombre effectif de cellules actives, éventuellement fractionnaires sur la dernière rangée coupée par le piston.

Si la $j$-ième rangée a pour bornes verticales $y_j^-=(j-1)\Delta y$ et $y_j^+=j\Delta y$, la fraction active de cette rangée vaut
\begin{equation}
 \phi_j(t)=\frac{\max\!\bigl(0,\min(H_{\mathrm{actif}}(t),y_j^+)-y_j^-\bigr)}{\Delta y},
 \qquad 0\le \phi_j\le 1.
\end{equation}
Cette fraction est recopiée sur toutes les colonnes de la rangée, ce qui définit une carte de fraction active \(\phi_c\) au niveau cellulaire.

Dans la version piston \emph{incompressible} du code, la cible de redistribution n'est plus adaptée au volume actif courant : la cible d'occupation de référence \(\gamma_0\) est conservée, et seule la géométrie active change. Le piston introduit alors deux quantités distinctes :
\begin{equation}
 \gamma_{\mathrm{target}} = \gamma_0,
 \qquad
 \gamma_{\mathrm{geom}}(t)=\frac{N_p}{\sum_c \phi_c(t)}.
\end{equation}
La première est la cible incompressible effectivement demandée à la redistribution ; la seconde est l'occupation moyenne qui serait géométriquement requise pour loger toutes les particules dans le volume actif sous le piston. Lorsque \(\gamma_{\mathrm{geom}}>\gamma_{\mathrm{target}}\), le piston pousse le système vers un régime de sur-occupation qui ne peut être compensé que par une forte montée de pression, puis par un blocage progressif de la redistribution lorsque la capacité disponible devient insuffisante.

Si l'option correspondante est activée, le code compacte aussi géométriquement le fluide sous le piston par un repositionnement affine vertical. Ce repositionnement n'est pas une modification de densité cible : il constitue uniquement un transport géométrique imposé, auquel la fermeture liquide réagit ensuite par la montée de la pression virielle.

\subsection{Pression de paroi}

L'impulsion échangée avec les parois est accumulée pour construire des diagnostics de pression instantanée de type flux d'impulsion. Si $\Delta P_y^{\mathrm{bot}}$ est la quantité de mouvement verticale échangée avec la paroi inférieure pendant un pas de temps, la pression instantanée associée est estimée par
\begin{equation}
 p_{\mathrm{bot}}^{\mathrm{inst}} = \frac{\Delta P_y^{\mathrm{bot}}}{L_x\,\Delta t}.
\end{equation}

Cette pression de paroi est une grandeur différente de la pression cinétique locale $P_{\mathrm{kin}}$ ou de la pression motrice $P_{\mathrm{drive}}$ utilisée dans la fermeture liquide.


\subsection{Particules virtuelles de paroi : modèle \texttt{wallVP}}

Les validations récentes de Poiseuille et de piston utilisent une condition limite enrichie par particules virtuelles de paroi. Le principe est de compléter les cellules coupées par une paroi solide avec une population virtuelle représentant la partie solide de la cellule. Cette population participe à la collision SRD afin de transmettre une contrainte tangentielle et normale plus réaliste que la seule réflexion géométrique des particules incidentes.

Pour une cellule de collision contenant \(N_f\) particules fluides et \(N_v\) particules virtuelles, on note
\begin{equation}
  \bm P_f = \sum_{i=1}^{N_f}\bm v_i,
  \qquad
  \bm P_v = \sum_{j=1}^{N_v}\bm v_j^{(v)},
  \qquad
  \bm U = \frac{\bm P_f+\bm P_v}{N_f+N_v}.
\end{equation}
La collision virtuelle reconstruit alors l'impulsion post-collision des particules virtuelles sous la forme
\begin{equation}
  \bm P_v' = N_v\bm U + \mathbf R(\theta)\bigl(\bm P_v-N_v\bm U\bigr),
\end{equation}
avec le même angle de rotation SRD que dans la cellule. L'impulsion transmise par la partie virtuelle est
\begin{equation}
  \Delta \bm P_v = \bm P_v' - \bm P_v.
\end{equation}
Ce terme doit être traité comme une impulsion signée. Pour la paroi supérieure, une compression positive correspond à une impulsion verticale positive dans le fluide ; pour la paroi inférieure, elle correspond à une impulsion verticale négative. Le diagnostic récent sépare donc explicitement : impulsion signée, pression positive de compression, contribution impact, contribution \texttt{wallVP}, total impact+VP, et moyennes cumulées sur fenêtre.

La pression mécanique instantanée est normalisée par
\begin{equation}
  p_{\mathrm{wall}}^{\mathrm{inst}} = \frac{\Delta p_n}{\Delta t\,L_x},
\end{equation}
la simulation étant bidimensionnelle avec profondeur unité. Cette grandeur est très bruitée, car elle divise une impulsion de collision par un pas de temps court. Le diagnostic physiquement exploitable est donc la pression moyenne obtenue par différence d'impulsions cumulées sur une fenêtre suffisamment longue,
\begin{equation}
  \overline p_{\mathrm{wall}}[t_a,t_b]
  = \frac{P_{\mathrm{wall}}^{\mathrm{cum}}(t_b)-P_{\mathrm{wall}}^{\mathrm{cum}}(t_a)}{(t_b-t_a)L_x}.
\end{equation}
Le test gaz à l'équilibre entre deux murs montre que cette moyenne cumulée retrouve la pression cinétique bulk, alors que les valeurs instantanées peuvent osciller de plusieurs milliers d'unités réduites autour d'une moyenne de l'ordre de \(2.1\times10^2\).

\section{Méthode de redistribution : cible, étapes et interprétation}

\subsection{Cible d'occupation et de densité}

La redistribution vise à contraindre l'occupation cellulaire vers une cible locale $N_c^{\star}$. Dans le cas le plus simple sans piston, cette cible vaut
\begin{equation}
 N_c^{\star} = \gamma_0,
\end{equation}
où $\gamma_0$ est le nombre moyen visé de particules par cellule pleinement active.

En présence d'une fraction de volume active $\phi_c\in[0,1]$ dans la cellule $c$, la cible locale devient
\begin{equation}
 N_c^{\star} = \gamma_{\mathrm{target}}\,\phi_c.
\end{equation}
Dans la version \emph{incompressible} actuelle du traitement du piston, on impose
\begin{equation}
 \gamma_{\mathrm{target}} = \gamma_0,
\end{equation}
indépendamment de la position du piston. Ainsi, la densité cible de nombre vaut
\begin{equation}
 n_c^{\star} = \frac{\gamma_0\,\phi_c}{V_c}.
\end{equation}

Il est utile d'introduire en parallèle l'occupation géométriquement requise
\begin{equation}
 \gamma_{\mathrm{geom}} = \frac{N_p}{\sum_c \phi_c},
\end{equation}
qui mesure combien de particules par cellule active seraient nécessaires si la totalité des particules restait confinée dans le volume disponible sous le piston. Le rapport
\begin{equation}
 \mathcal{R}_{\mathrm{cap}} = \frac{\gamma_{\mathrm{geom}}}{\gamma_{\mathrm{target}}}
\end{equation}
sert alors de diagnostic de compression. Lorsque \(\mathcal{R}_{\mathrm{cap}}\) dépasse l'unité, la compression géométrique demandée par le piston excède la capacité incompressible nominale du maillage : la redistribution doit forcer le système au voisinage de sa borne de capacité, et la réaction physique attendue est une forte montée de la pression virielle plutôt qu'une adaptation automatique de la densité cible.

Cette distinction entre \(\gamma_{\mathrm{target}}\) et \(\gamma_{\mathrm{geom}}\) est centrale. Dans l'ancienne logique quasi-compressible, la cible elle-même augmentait avec la diminution du volume actif ; dans la version actuelle, on garde au contraire une densité cible constante et on laisse la pression porter la réaction à la compression.

\subsection{Bande admissible}

La redistribution n'impose pas strictement $N_c=N_c^{\star}$ à chaque pas ; elle cherche plutôt à maintenir $N_c$ dans une bande admissible autour de la cible. Une écriture physique simple est
\begin{equation}
 N_c^{\mathrm{low}} \approx N_c^{\star}(1-\eta),
\qquad
 N_c^{\mathrm{high}} \approx N_c^{\star}(1+\eta),
\end{equation}
où $\eta$ est un paramètre de tolérance relié au champ \verb|coef| dans le code. L'implémentation utilise ensuite des arrondis, des planchers et, dans certaines configurations, des seuils locaux modifiés par la topologie active.

\subsection{Étapes de la redistribution}

La redistribution agit comme un opérateur de projection géométrique sur le champ d'occupation. Elle n'intègre pas une équation d'évolution de type Navier--Stokes ; elle déplace directement des particules pour rapprocher le champ discret $N_c$ de la cible locale $N_c^{\star}$, tout en préservant au mieux la structure spatiale du fluide.

Pour fixer les notations, on introduit pour chaque cellule $c$ :
\begin{itemize}
  \item l'occupation actuelle $N_c$ ;
  \item la cible locale $N_c^{\star}$ ;
  \item les seuils admissibles $N_c^{\mathrm{low}}$ et $N_c^{\mathrm{high}}$ ;
  \item le surplus
  \begin{equation}
    \Delta N_c^{+} = \max\bigl(N_c - N_c^{\mathrm{high}}, 0\bigr),
  \end{equation}
  \item le déficit
  \begin{equation}
    \Delta N_c^{-} = \max\bigl(N_c^{\mathrm{low}} - N_c, 0\bigr).
  \end{equation}
\end{itemize}

Une cellule est dite \emph{trop dense} si $\Delta N_c^{+}>0$, et \emph{trop pauvre} si $\Delta N_c^{-}>0$. Le processus de redistribution est alors organisé en plusieurs sous-opérations.

\paragraph{a) Identification topologique des cellules.}
Le code distingue les cellules de \emph{bulk}, entourées de voisins occupés, des cellules \emph{interfaciales} ou \emph{adjacentes à une paroi}. Cette distinction est importante car la politique de redistribution ne doit pas être la même dans un cœur de fluide homogène, près d'une interface libre, ou au voisinage d'un bord mouillant. La topologie locale sert donc à définir l'espace des cibles admissibles et la direction privilégiée des transferts.

Du point de vue algorithmique, une cellule $c$ n'est déclarée \emph{bulk} que si deux conditions sont satisfaites simultanément :
\begin{enumerate}[label=(\roman*)]
  \item la cellule n'est adjacente à aucune paroi physique ;
  \item tous les voisins du stencil de Moore d'ordre~1, noté $\mathcal{N}_1(c)$, sont occupés.
\end{enumerate}
Si l'on note $\chi_c\in\{0,1\}$ l'indicatrice d'occupation de la cellule $c$, la logique peut s'écrire sous la forme
\begin{equation}
  c\in\mathcal{B}
  \quad \Longleftrightarrow \quad
  c\notin\partial\Omega_{\mathrm{phys}}
  \ \text{et}\ 
  \chi_{c'}=1 \ \text{pour tout}\ c'\in\mathcal{N}_1(c),
\end{equation}
où $\mathcal{B}$ désigne l'ensemble des cellules de bulk et $\partial\Omega_{\mathrm{phys}}$ l'ensemble des cellules collées à une frontière non périodique.

Cette définition est volontairement stricte : dès qu'un voisin orthogonal ou diagonal est vide, ou dès qu'une cellule touche une paroi réelle, la cellule est sortie du bulk et traitée comme cellule interfaciale. Cette convention est cohérente avec la logique physique du code : le bulk est la région où la redistribution peut être effectuée comme un transport quasi isotrope, tandis que les cellules non-bulk doivent respecter la topologie de l'interface, le mouillage ou la géométrie active imposée par un piston.

\paragraph{b) Traitement des cellules trop denses.}
Si $N_c > N_c^{\mathrm{high}}$, la cellule doit céder $\Delta N_c^{+}$ particules. Le transfert ne s'effectue pas vers des cellules arbitraires, mais vers un ensemble de cellules réceptrices $\mathcal{A}_c$ satisfaisant deux contraintes :
\begin{equation}
  N_{c'} < N_{c'}^{\mathrm{target}},
  \qquad c'\in\mathcal{A}_c,
\end{equation}
où $N_{c'}^{\mathrm{target}}$ est soit le seuil bas, soit le seuil haut selon l'étape du processus.

Le code procède d'abord par décharge vers les cellules manifestement sous-peuplées, puis, s'il subsiste un excès, vers des cellules simplement capables d'absorber du fluide sans dépasser leur borne haute. Formellement, si $\mathcal{A}_c^{\mathrm{poor}}$ désigne l'ensemble des cellules sous-denses et $\mathcal{A}_c^{\mathrm{relief}}$ les cellules capables d'absorber un surplus, la redistribution réalise schématiquement
\begin{equation}
  c \longrightarrow \mathcal{A}_c^{\mathrm{poor}} \subseteq \mathcal{A}_c^{\mathrm{relief}}.
\end{equation}

Le classement des receveurs n'est pas isotrope pur : il combine la capacité disponible, la distance, et une direction de transfert privilégiée induite par le gradient local d'occupation ou par la normale intérieure à la paroi. Si $\bm{n}_c$ désigne cette direction privilégiée et $\bm{d}_{cc'}$ le vecteur reliant les centres des cellules $c$ et $c'$, un score typique est de la forme
\begin{equation}
  S_{cc'} = w_1\,\mathcal{C}_{cc'} + w_2\,\frac{\bm{d}_{cc'}\cdot \bm{n}_c}{\norm{\bm{d}_{cc'}}} - w_3\,\norm{\bm{d}_{cc'}},
\end{equation}
où $\mathcal{C}_{cc'}$ représente la capacité d'absorption de la cellule réceptrice, et $w_1,w_2,w_3$ des pondérations implicites dans l'algorithme. L'idée physique est simple : on favorise les transferts vers des régions moins denses et dirigées vers l'intérieur du fluide.

La direction $\bm{n}_c$ n'est pas calculée sur le champ brut $N_c$, mais sur un champ d'occupation lissé $\widetilde{N}_c$ obtenu par convolution locale de type
\begin{equation}
  \widetilde{N}_{i,j} = \frac{1}{16}\Bigl(N_{i-1,j-1}+2N_{i-1,j}+N_{i-1,j+1}+2N_{i,j-1}+4N_{i,j}+2N_{i,j+1}+N_{i+1,j-1}+2N_{i+1,j}+N_{i+1,j+1}\Bigr),
\end{equation}
à renormalisation près lorsque le stencil est tronqué par une frontière non périodique. Le gradient utilisé pour orienter la redistribution s'écrit alors
\begin{equation}
  \nabla \widetilde{N}_c \approx \left(\partial_x \widetilde{N}_c,\, \partial_y \widetilde{N}_c\right),
\end{equation}
avec des différences finies centrées à l'intérieur du domaine et décentrées au premier ordre sur les bords non périodiques. Si l'axe considéré est périodique, la différence finie correspondante est évaluée avec raccordement circulaire des indices ; sinon on utilise une dérivée unilatérale. En notation continue, pour une cellule centrée en $(x_i,y_j)$,
\begin{equation}
  \partial_x \widetilde{N}_{i,j} \approx
  \begin{cases}
    \dfrac{\widetilde{N}_{i+1,j}-\widetilde{N}_{i-1,j}}{2\Delta x}, & 2\le i\le N_x-1,\\[1ex]
    \dfrac{\widetilde{N}_{2,j}-\widetilde{N}_{1,j}}{\Delta x}, & i=1 \text{ et bord non périodique},\\[1ex]
    \dfrac{\widetilde{N}_{N_x,j}-\widetilde{N}_{N_x-1,j}}{\Delta x}, & i=N_x \text{ et bord non périodique},
  \end{cases}
\end{equation}
et de même pour $\partial_y \widetilde{N}_{i,j}$. L'orientation de la redistribution dépend donc explicitement du choix des conditions aux limites.

Près d'une paroi physique, cette direction n'est plus laissée au seul gradient d'occupation : elle est remplacée ou combinée avec la normale intérieure à la paroi, ce qui impose que les particules en excès soient réinjectées vers le bulk plutôt que parallèlement au mur.

\paragraph{c) Traitement des cellules trop pauvres.}
Si $0 < N_c < N_c^{\mathrm{low}}$, la cellule doit recevoir des particules. Le code sélectionne alors des cellules donneuses voisines $d$ telles que
\begin{equation}
  N_d > N_d^{\mathrm{low}},
\end{equation}
avec une préférence pour les cellules du bulk, afin de ne pas déstabiliser l'interface. Si $\mathcal{D}_c$ désigne l'ensemble des donneurs admissibles, le nombre prélevé sur la cellule $d\in\mathcal{D}_c$ est borné par
\begin{equation}
  n_{d\to c} \le \min\bigl(N_d - N_d^{\mathrm{low}},\, N_c^{\mathrm{low}}-N_c\bigr).
\end{equation}

On voit ainsi que la redistribution n'a pas pour but d'égaliser instantanément tous les $N_c$ à $N_c^{\star}$, mais de maintenir chaque cellule dans une bande admissible, avec priorité donnée à la stabilité topologique et à la connectivité du fluide.

\paragraph{d) Redistribution à portée finie.}
Le voisinage principal est d'ordre~1, mais le code ouvre aussi une portée~2 si les cibles immédiates sont insuffisantes. Cette extension peut être vue comme une relaxation du principe local, nécessaire lorsque le fluide se compacte fortement ou lorsqu'une paroi impose une géométrie active réduite. L'opérateur reste cependant fortement local : on ne transporte jamais des particules sur de longues distances en un seul pas.

En pratique, l'algorithme utilise d'abord le voisinage de Moore d'ordre~1,
\begin{equation}
  \mathcal{N}_1(c) = \{c' \neq c\ ;\ \max(|\delta i|,|\delta j|)=1\},
\end{equation}
puis, si l'équilibrage échoue, le voisinage d'ordre~2,
\begin{equation}
  \mathcal{N}_2(c) = \{c' \neq c\ ;\ \max(|\delta i|,|\delta j|)\le 2\}.
\end{equation}
Les candidats sont ensuite ordonnés de manière hiérarchique : d'abord par capacité disponible, puis par alignement avec la direction privilégiée, enfin par distance géométrique. Le code exploite pour cela le pas minimal entre cellules en tenant compte de la périodicité éventuelle sur chaque axe, de sorte que la notion de ``plus proche'' dépend bien de la topologie réelle du domaine.

\paragraph{e) Overflow bulk-interface et saturation.}
Lorsque le cœur du fluide est trop dense et que les voisins proches ne suffisent plus, le code autorise un report du surplus du bulk vers les cellules géométriquement interfaciales les plus proches. Cela correspond à une politique d'\emph{overflow} : l'excès de masse numérique est repoussé vers les régions déjà identifiées comme interface ou bord libre. Si même cette étape ne suffit pas, le code signale une \emph{saturation} de redistribution. Physiquement, cela signifie que la contrainte de quasi-incompressibilité devient plus difficile à satisfaire avec les seules relocalisations locales permises par le maillage et la topologie instantanée.

\paragraph{f) Interprétation physique.}
Du point de vue continu, on peut interpréter la redistribution comme une étape de projection approximative vers l'hypersurface discrète
\begin{equation}
  \mathcal{M}_{\rho^{\star}} = \left\{ \{N_c\}_{c=1}^{N_xN_y} \; ; \; N_c \approx N_c^{\star} \text{ pour tout } c \right\},
\end{equation}
qui représente un régime à compressibilité très faible. L'opérateur ne résout pas une équation de pression au sens analytique ; il agit comme une correction de densité imposée sur le champ particulaire lui-même.

\subsection{Correction locale de quantité de mouvement dans la redistribution}

Si $\bm{P}_c^{\mathrm{before}}$ et $\bm{P}_c^{\mathrm{after}}$ sont les quantités de mouvement de la cellule $c$ avant et après les transferts, le code calcule
\begin{equation}
 \delta \bm{P}_c = \bm{P}_c^{\mathrm{before}} - \bm{P}_c^{\mathrm{after}},
\end{equation}
et applique, quand cela est permis, une correction uniforme par particule
\begin{equation}
 \delta \bm{u}_c = \frac{\delta \bm{P}_c}{N_c},
\end{equation}
qui redonne à la cellule une quantité de mouvement proche de sa valeur précédente. Cette correction reste locale et ne constitue pas encore la fermeture liquide elle-même ; elle sert d'abord à éviter que la redistribution ne détruise inutilement le champ moyen.


\subsection{Traitement de surface, mouillage et réorientation interfaciale}

Le code distingue soigneusement deux familles de mécanismes de surface :
\begin{enumerate}[label=\alph*)]
  \item un \textbf{traitement géométrique de surface} intégré à la redistribution, destiné à préserver une topologie cohérente de l'interface et, si nécessaire, une couche mouillante près des parois ;
  \item une \textbf{réorientation cinématique interfaciale} appliquée après la fermeture liquide, destinée à modéliser une cohésion tangentielle ou une tension superficielle effective.
\end{enumerate}

Ces deux mécanismes ne jouent pas le même rôle.

\paragraph{a) Surface et mouillage dans la redistribution.}
Lorsqu'une paroi porte une condition de type mouillage, la redistribution peut imposer localement une occupation minimale dans les cellules voisines de cette paroi. Si l'on note $N_c^{\mathrm{wet}}$ la cible locale de mouillage, la contrainte s'écrit schématiquement
\begin{equation}
  N_c \gtrsim N_c^{\mathrm{wet}}
  \qquad \text{pour } c \in \mathcal{W},
\end{equation}
où $\mathcal{W}$ désigne l'ensemble des cellules de la couche mouillante.

Le but n'est pas d'introduire une force de surface explicite, mais de maintenir une population suffisante à proximité des bords physiques pour éviter que l'interface ne se décroche numériquement de la paroi. Cette logique est particulièrement utile lorsque le fluide est soumis à la gravité, à une compression par piston, ou à des redistributions répétées.

\paragraph{b) Direction privilégiée et normales de surface.}
Dans les cellules interfaciales, la redistribution utilise une normale intérieure approchée $\bm{n}_c$ obtenue à partir du gradient local d'occupation lissée. Si $\widetilde{N}(x,y)$ désigne un champ d'occupation discrètement lissé, la normale intérieure s'écrit, à normalisation près,
\begin{equation}
  \bm{n}_c \propto \nabla \widetilde{N}_c.
\end{equation}
Le signe est choisi de façon à pointer vers le bulk. Cette normale intervient à la fois dans le classement des cellules receveuses lors de la redistribution et, plus tard, dans la réorientation des vitesses interfaciales.

Du point de vue algorithmique, cette normale est obtenue par les mêmes opérateurs de gradient que ceux utilisés pour l'orientation de redistribution. Le champ lissé $\widetilde N$ est d'abord construit, puis les composantes
\begin{equation}
  g_x = \partial_x \widetilde N,
  \qquad
  g_y = \partial_y \widetilde N
\end{equation}
sont évaluées cellule par cellule. La direction intérieure associée à une cellule interfaciale est alors
\begin{equation}
  \widehat{\bm n}_c = \frac{(g_x,g_y)}{\sqrt{g_x^2+g_y^2}+\varepsilon_n},
\end{equation}
où $\varepsilon_n>0$ évite les divisions par zéro lorsque le gradient local est trop faible. Si la cellule est directement au contact d'une paroi physique, cette normale est remplacée par la normale géométrique intérieure au mur ; si plusieurs parois se rencontrent, par exemple dans un coin, les normales associées sont sommées puis renormalisées. Cette règle garantit que la direction privilégiée reste compatible avec la géométrie du domaine même lorsque le gradient d'occupation local devient ambigu.

\paragraph{c) Réorientation interfaciale.}
Après redistribution, correction de champ et kick liquide, le code peut appliquer une réorientation partielle des vitesses dans les cellules identifiées comme interfaciales. Soit $\bm{v}_i$ la vitesse d'une particule d'interface, et $\widehat{\bm{n}}_c$ la normale intérieure unitaire associée à la cellule $c$. La réorientation construit une vitesse cible
\begin{equation}
  \bm{v}_i^{\mathrm{tar}} = \abs{\bm{v}_i}\,\widehat{\bm{n}}_c,
\end{equation}
puis applique une interpolation linéaire
\begin{equation}
  \bm{v}_i^{\mathrm{new}} = (1-\beta_c)\,\bm{v}_i + \beta_c\,\bm{v}_i^{\mathrm{tar}},
\end{equation}
où $\beta_c\in[0,1]$ est un coefficient local de réorientation.

Le code ne réoriente qu'une fraction $\varphi_c\in[0,1]$ des particules de la cellule, ce qui revient à imposer une correction cinématique partielle. Les coefficients $\beta_c$ et $\varphi_c$ peuvent être renforcés en fonction d'un indicateur de courbure ou de désalignement local. On peut formaliser cela par
\begin{equation}
  \beta_c = \beta_0\,\mathcal{F}_\beta(\kappa_c),
  \qquad
  \varphi_c = \varphi_0\,\mathcal{F}_\varphi(\kappa_c),
\end{equation}
où $\kappa_c$ désigne un indicateur local de courbure ou de désalignement, et $\beta_0$, $\varphi_0$ des paramètres de base.

Dans le code, l'indicateur $\kappa_c$ n'est pas une courbure au sens différentiel strict $\nabla\cdot\widehat{\bm n}$, mais un défaut d'alignement des normales locales calculé à partir du voisinage de la cellule interfaciale. L'effet est le suivant : plus les normales voisines varient fortement, plus l'interface est localement coudée ou irrégulière, et plus les paramètres de réorientation sont augmentés. La réorientation agit donc davantage dans les régions de forte courbure apparente ou de surface mal alignée.

\paragraph{d) Signification physique.}
La réorientation n'est pas une tension superficielle au sens strict d'une force de courbure écrite comme $\sigma \kappa \bm{n}$. Elle agit plutôt comme une régularisation cinématique qui tend à réaligner les vitesses interfaciales vers le bulk, donc à limiter l'étalement tangent et la dispersion de l'interface. On peut la voir comme un analogue discret d'une dissipation interfaciale orientée ou d'une cohésion de surface effective.

Le choix de la placer \emph{après} la fermeture liquide est cohérent avec l'architecture actuelle : la réorientation admet localement des cellules de surface plus pauvres, alors que redistribution, correction et kick cherchent au contraire à maintenir une occupation contrôlée. Il est donc préférable de ne pas inclure la réorientation dans l'état de référence utilisé pour la fermeture.

\section{Correction des champs : vitesse, température, énergie}

\subsection{Référence, état redistribué et défaut à corriger}

Soit $\bm{u}_c^{\mathrm{ref}}$ la vitesse moyenne de la cellule $c$ dans l'état de référence pré-redistribution, et $\bm{u}_c^{\mathrm{red}}$ la vitesse moyenne de cette même cellule après redistribution. On définit le défaut de quantité de mouvement
\begin{equation}
 \delta \bm{P}_c^{\mathrm{rep}} = \bm{P}_c^{\mathrm{ref}} - \bm{P}_c^{\mathrm{red}}.
\end{equation}

La correction externe de fermeture liquide introduit alors un champ de réparation uniforme dans la cellule,
\begin{equation}
 \delta \bm{u}_c^{\mathrm{rep}} = \frac{\delta \bm{P}_c^{\mathrm{rep}}}{N_c^{\star}},
\end{equation}
où $N_c^{\star}$ est la cible locale d'occupation.

Le code multiplie ensuite ce champ par un scalaire $\beta_{\mathrm{rep}}\ge 0$ choisi soit manuellement, soit par minimisation quadratique d'un défaut global sur les vitesses de cellule. La vitesse réparée s'écrit donc
\begin{equation}
 \bm{u}_c^{\mathrm{rep}} = \bm{u}_c^{\mathrm{red}} + \beta_{\mathrm{rep}}\,\delta \bm{u}_c^{\mathrm{rep}}.
\end{equation}

Dans le code, le coefficient $\beta_{\mathrm{rep}}$ peut être choisi par minimisation quadratique de l'écart entre champ redistribué et champ de référence. Si l'on note
\begin{equation}
  \bm e_c = \bm u_c^{\mathrm{red}} - \bm u_c^{\mathrm{ref}},
  \qquad
  \bm a_c = \delta \bm u_c^{\mathrm{rep}},
\end{equation}
alors le problème traité est de minimiser
\begin{equation}
  J(\beta) = \sum_c \abs{\bm e_c + \beta \, \bm a_c}^2,
\end{equation}
d'où l'optimum formel
\begin{equation}
  \beta_{\mathrm{rep}}^{\star} = -\frac{\sum_c \bm e_c\cdot \bm a_c}{\sum_c \abs{\bm a_c}^2},
\end{equation}
ensuite tronqué si nécessaire pour respecter les contraintes de l'utilisateur. Ainsi, la correction de champ n'est pas seulement une translation imposée ; elle est calibrée pour rapprocher au mieux le champ redistribué du champ de référence dans une norme quadratique discrète.

\subsection{Température : ce qui est corrigé et ce qui ne l'est pas}

Dans l'état actuel du couplage intégré, la correction externe agit explicitement sur la \textbf{vitesse moyenne} de chaque cellule, pas sur la température. Cette absence d'une correction thermique explicite n'est pas incohérente : une correction uniforme de type
\begin{equation}
 \bm{v}_i \leftarrow \bm{v}_i + \delta \bm{u}_c,
 \qquad i\in c,
\end{equation}
ne modifie pas les fluctuations relatives $\bm{v}_i-\bm{u}_c$ à appartenance cellulaire fixée, donc ne modifie pas directement $k_B T_c$.

En revanche, la redistribution elle-même déplace des particules d'une cellule à l'autre et peut donc modifier la température locale reconstruite. Dans le code actuel, la régularisation thermique est surtout assurée par le thermostat MPCD des collisions SRD et, éventuellement, par les parois thermalisantes.

\subsection{Effet énergétique d'un kick uniforme par cellule}

Considérons une cellule contenant $N_c$ particules et soumise à une correction uniforme $\delta \bm{u}_c$. Si l'on note $\bm{P}_c = \sum_{i\in c} \bm{v}_i$ sa quantité de mouvement avant correction, alors la variation d'énergie cinétique totale de la cellule est
\begin{equation}
 \Delta E_c = \bm{P}_c\cdot \delta \bm{u}_c + \frac{N_c}{2}\abs{\delta \bm{u}_c}^2.
\end{equation}

Cette variation affecte l'énergie macroscopique liée au mouvement moyen de la cellule. L'énergie interne relative,
\begin{equation}
 E_{\mathrm{int},c} = \frac{1}{2}\sum_{i\in c}\abs{\bm{v}_i-\bm{u}_c}^2,
\end{equation}
reste inchangée par un \emph{kick} uniforme à appartenance cellulaire fixée.

Ainsi, dans l'état actuel du code, la correction et le \emph{kick} incompressible modifient la répartition entre énergie moyenne et énergie potentiellement injectée/retraitée par le forçage, mais ne constituent pas un thermostat supplémentaire.

\section{Kick incompressible et fermeture liquide}

\subsection{Pression motrice}

La fermeture liquide actuelle repose sur une pression motrice construite à partir de l'état de référence pré-redistribution,
\begin{equation}
 P_{\mathrm{drive}} = P_{\mathrm{kin}} + P_{\mathrm{vir}}.
\end{equation}

La partie cinétique est
\begin{equation}
 P_{\mathrm{kin},c} = n_c^{\mathrm{ref}}\,k_B T_c^{\mathrm{ref}},
\end{equation}
où $n_c^{\mathrm{ref}}$ est la densité de nombre de la cellule $c$ dans l'état de référence, et $T_c^{\mathrm{ref}}$ la température locale reconstruite dans ce même état.

La partie virielle est définie par
\begin{equation}
 P_{\mathrm{vir},c} = K_{\mathrm{virial}}\,\bigl(n_c^{\mathrm{ref}} - n_c^{\star}\bigr),
\end{equation}
où $K_{\mathrm{virial}}$ est le coefficient de raideur virielle et $n_c^{\star}$ la densité cible issue de la redistribution.

Ainsi,
\begin{equation}
 P_{\mathrm{drive},c} = n_c^{\mathrm{ref}} k_B T_c^{\mathrm{ref}} + K_{\mathrm{virial}}\,\bigl(n_c^{\mathrm{ref}} - n_c^{\star}\bigr).
\end{equation}

\subsection{Interprétation physique}

Le rôle de $P_{\mathrm{vir}}$ n'est pas de remplacer la pression cinétique, mais d'ajouter une raideur effective au fluide lorsque la redistribution maintient $n_c$ près d'une cible donnée. Dans un régime quasi incompressible, il est naturel que la pente hydrostatique moyenne soit portée principalement par $P_{\mathrm{vir}}$ plutôt que par $P_{\mathrm{kin}}$.

Si la redistribution est très serrée, le gradient de densité accessible au fluide est lui-même borné ; dans ce cas, modifier $K_{\mathrm{virial}}$ change surtout l'amplitude de la contre-pression et la charge algorithmique supportée par la redistribution, sans nécessairement changer fortement l'état stationnaire moyen.


\subsection{Compression par piston et module volumique effectif}

Dans le cas du piston incompressible, la densité cible n'est plus adaptée géométriquement : la compression du volume actif doit donc se traduire par une montée de pression. Si l'on note \(V_{\mathrm{actif}}(t)=L_x H_{\mathrm{actif}}(t)\) l'aire active sous le piston et \(P_{\mathrm{tot,actif}}(t)\) la pression totale moyenne du modèle sur ce domaine actif, on définit le module volumique effectif
\begin{equation}
 B_{\mathrm{eff}}(t)
 =
 -\,V_{\mathrm{actif}}(0)\,
 \frac{P_{\mathrm{tot,actif}}(t)-P_{\mathrm{tot,actif}}(0)}
 {V_{\mathrm{actif}}(t)-V_{\mathrm{actif}}(0)}
 \;\approx\;
 \frac{\Delta P_{\mathrm{tot,actif}}(t)}{-\Delta V_{\mathrm{actif}}(t)/V_{\mathrm{actif}}(0)}.
\end{equation}
Sa réciproque
\begin{equation}
 \kappa_{\mathrm{eff}}(t)=\frac{1}{B_{\mathrm{eff}}(t)}
\end{equation}
joue le rôle d'une compressibilité effective.

Le code calcule aussi un module local apparent dérivé de l'équation d'état discrète,
\begin{equation}
 \left(\frac{\mathrm d P_{\mathrm{tot}}}{\mathrm d \rho}\right)_{\mathrm{eff}}
 \approx
 \frac{P_{\mathrm{tot,actif}}(t)-P_{\mathrm{tot,actif}}(0)}
 {\rho_{\mathrm{actif}}(t)-\rho_{\mathrm{actif}}(0)},
\end{equation}
ainsi que la quantité
\begin{equation}
 \rho_{\mathrm{actif}}(t)\,
 \left(\frac{\mathrm d P_{\mathrm{tot}}}{\mathrm d \rho}\right)_{\mathrm{eff}},
\end{equation}
qui doit être comparable à \(B_{\mathrm{eff}}\) si la montée de pression est effectivement portée par la variation de densité moyenne dans le domaine comprimé. Ces diagnostics permettent de quantifier explicitement la faible compressibilité du modèle plutôt que de la déduire uniquement de la forme des profils.

\subsection{Gradient de pression et \emph{kick} de vitesse}

Le code calcule le gradient de $P_{\mathrm{drive}}$ avec un opérateur compatible avec les conditions aux limites de chaque axe. Le \emph{kick} incompressible de fermeture s'écrit ensuite
\begin{equation}
 \delta \bm{u}_c^{\mathrm{EOS}} = -\beta_{\mathrm{EOS}}\,\frac{\Delta t}{n_c^{\star}}\,\nabla P_{\mathrm{drive},c},
\end{equation}
où $\beta_{\mathrm{EOS}}$ est un facteur scalaire de réglage et $n_c^{\star}$ la densité cible locale.

Dans la version actuelle du couplage, les dérivées spatiales sont évaluées sur la grille d'analyse par différences finies. Si $P_{i,j}$ désigne la pression motrice dans la cellule centrée en $(x_i,y_j)$, alors les composantes du gradient sont approchées par
\begin{equation}
  \partial_x P_{i,j} \approx
  \begin{cases}
    \dfrac{P_{i+1,j}-P_{i-1,j}}{2\Delta x}, & 2\le i\le N_x-1,\\[1ex]
    \dfrac{P_{2,j}-P_{1,j}}{\Delta x}, & i=1,\\[1ex]
    \dfrac{P_{N_x,j}-P_{N_x-1,j}}{\Delta x}, & i=N_x,
  \end{cases}
\end{equation}
\begin{equation}
  \partial_y P_{i,j} \approx
  \begin{cases}
    \dfrac{P_{i,j+1}-P_{i,j-1}}{2\Delta y}, & 2\le j\le N_y-1,\\[1ex]
    \dfrac{P_{i,2}-P_{i,1}}{\Delta y}, & j=1,\\[1ex]
    \dfrac{P_{i,N_y}-P_{i,N_y-1}}{\Delta y}, & j=N_y.
  \end{cases}
\end{equation}
Cette écriture explicite la logique du code : dérivées centrées dans le bulk, décentrées au premier ordre au voisinage des bords. Le choix est cohérent avec le fait que l'état de fermeture n'est pas construit comme un champ périodique abstrait, mais comme un champ défini dans la géométrie physique du canal ou du domaine comprimé par piston.

Le signe est celui attendu pour une force de pression. Si la direction verticale $y$ est croissante vers le haut et si la gravité vaut $g<0$, alors un équilibre hydrostatique moyen correspond à
\begin{equation}
 \frac{\mathrm{d}P}{\mathrm{d}y} \approx n^{\star} g
\end{equation}
dans les unités du code, ou plus généralement à
\begin{equation}
 \frac{\mathrm{d}P}{\mathrm{d}y} \approx \rho_m g
\end{equation}
si l'on réintroduit explicitement la masse particulaire.

\subsection{Résidu hydrostatique}

Le post-traitement reconstruit les profils moyens selon $x$ et calcule le résidu hydrostatique
\begin{equation}
 R_h(y) = \frac{\mathrm{d}P_{\mathrm{tot}}}{\mathrm{d}y} - n(y)g
\end{equation}
ou, selon la convention choisie,
\begin{equation}
 R_h^{\star}(y) = \frac{\mathrm{d}P_{\mathrm{tot}}}{\mathrm{d}y} - n^{\star}(y)g.
\end{equation}

Dans une simulation MPCD bruitée, un résidu instantané local non nul n'est pas en soi rédhibitoire. La quantité significative est plutôt le résidu moyen ou lissé, et la comparaison entre la pente moyenne de $P_{\mathrm{tot}}$ et le forçage $n^{\star}g$ dans le \emph{bulk} du domaine.

\section{Bilan de la chaîne redistribution/liquid-closure : état historique}

L'état actuel du code peut être résumé de la manière suivante.

\begin{enumerate}[label=\arabic*)]
  \item La méthode MPCD proprement dite est cohérente : advection, décalage aléatoire de grille, rotation SRD et thermostat cellulaire sont bien en place.
  \item Les conditions limites sont suffisamment riches pour traiter des parois spéculaires, non glissantes, thermalisantes et un piston mobile.
  \item La redistribution agit comme un opérateur de projection vers un régime faiblement compressible contrôlé par une cible d'occupation locale. Dans la version piston actuelle, cette cible peut être maintenue constante afin de rendre la compression essentiellement virielle.
  \item La correction externe actuellement couplée agit d'abord sur le \emph{champ moyen de vitesse}, puis applique un \emph{kick} de pression construit à partir de la pression motrice \(P_{\mathrm{kin}}+P_{\mathrm{vir}}\) évaluée sur la référence.
  \item Il n'existe pas encore, dans le couplage intégré, de correction thermique explicite séparée ; la température reste principalement contrôlée par le thermostat MPCD et par les conditions limites thermalisantes.
  \item Les diagnostics récents distinguent maintenant les états instantanés, les champs moyens temporels, les diagnostics Poiseuille et les diagnostics piston, ce qui permet d'analyser séparément le bulk, la paroi et la réaction de la fermeture liquide.
\end{enumerate}

\section{Résultats historiques de la redistribution et de la fermeture liquide}

\subsection{Cas Poiseuille}

Le cas Poiseuille a servi de banc d'essai principal pour qualifier à la fois le bulk MPCD, les conditions limites et l'effet de la redistribution/fermeture liquide.

\paragraph{Effet de la résolution et des conditions limites.}
Les premiers essais ont montré un glissement important aux parois lorsque la direction transverse était insuffisamment résolue ou lorsque les parois utilisaient une thermalisation. Le raffinement transverse \((N_y=100)\) et l'emploi de conditions limites de type \emph{bounce-back} ont nettement amélioré la situation : le profil longitudinal \(u_x(y)\) devient quasi parfaitement parabolique, avec un coefficient de détermination du fit quadratique proche de l'unité. Cette observation montre que la qualité des parois et la résolution transverse sont déterminantes avant même d'analyser la fermeture incompressible.

\paragraph{Comparaison compressible / incompressible.}
En présence de redistribution et de fermeture liquide, les fluctuations d'occupation sont fortement réduites, typiquement avec \(\mathrm{std}(N)/\gamma \approx 0.13\) et \(\mathrm{outBand}=0\), alors que le cas compressible laisse réapparaître des fluctuations plus larges et des cellules hors bande. Le profil de Poiseuille reste cependant très propre dans les deux cas une fois la résolution et les parois convenablement choisies. La redistribution et la fermeture liquide agissent alors surtout comme un resserrement du système : elles réduisent la compressibilité numérique, diminuent le glissement effectif et stabilisent l'identification de la viscosité.

\paragraph{Convergence à long temps.}
Sur des temps de simulation suffisamment longs, les estimations de viscosité par courbure du profil et par contrainte pariétale deviennent cohérentes. Dans les runs de référence longs, on observe typiquement
\begin{equation}
 \nu_{\mathrm{curve}} \approx \nu_{\tau},
\end{equation}
avec un fit quadratique excellent (\(R^2 \simeq 0.999\)). Cela indique que le bulk et les parois racontent enfin la même physique, et que le régime de Poiseuille est effectivement établi. Il reste un glissement absolu non nul aux parois, qui doit être interprété comme une propriété effective de la condition limite \emph{bounce-back} à résolution finie, plutôt que comme une rupture de la structure parabolique du profil.

\subsection{Cas \texttt{compressibility} : hydrostatique moyenne et fermeture liquide}

Le cas \texttt{compressibility} a été utilisé pour qualifier la capacité de la fermeture liquide à produire une hydrostatique moyenne. Les diagnostics instantanés sont naturellement bruités, comme attendu en MPCD, mais les champs moyens temporels montrent un comportement très clair :
\begin{itemize}
  \item la température moyenne reste presque uniforme ;
  \item la densité moyenne reste proche de la cible, avec une stratification faible mais non nulle ;
  \item la pression totale moyenne
  \begin{equation}
   \overline{P}_{\mathrm{tot}}(y)=\overline{P}_{\mathrm{kin}}(y)+\overline{P}_{\mathrm{vir}}(y)
  \end{equation}
  suit une loi affine raisonnable dans le bulk ;
  \item le résidu hydrostatique moyen
  \begin{equation}
   \overline{R}_h(y)=\frac{\mathrm d \overline{P}_{\mathrm{tot}}}{\mathrm d y}-\overline{n}(y)\,g
  \end{equation}
  reste faible au centre du domaine et ne se dégrade nettement qu'au voisinage des bords.
\end{itemize}

L'interprétation principale est que l'hydrostatique moyenne est portée majoritairement par le terme viriel. Ce résultat est cohérent avec l'objectif poursuivi : pour un liquide quasi incompressible, on n'attend pas que la seule pression cinétique porte la réaction mécanique, mais qu'une pression de fermeture très raide prenne le relais. La comparaison de différentes valeurs de \(K_{\mathrm{virial}}\) montre par ailleurs qu'une redistribution très serrée borne déjà fortement le défaut de densité accessible ; dans ce régime, augmenter \(K_{\mathrm{virial}}\) modifie surtout l'amplitude de la pression virielle et la charge algorithmique de redistribution, sans changer drastiquement l'état moyen.

\subsection{Cas piston : compression incompressible et montée de pression}

Le traitement piston a été réorienté vers une logique véritablement quasi incompressible. Le piston réduit le volume actif en diminuant le nombre effectif de cellules disponibles, mais sans adapter la cible \(\gamma_{\mathrm{target}}\) à cette réduction. Par conséquent :
\begin{itemize}
  \item l'occupation géométriquement requise \(\gamma_{\mathrm{geom}}\) augmente avec la compression ;
  \item la cible incompressible \(\gamma_{\mathrm{target}}=\gamma_0\) reste constante ;
  \item la redistribution ne peut plus ``absorber'' la compression en augmentant artificiellement la densité cible ;
  \item la réaction attendue est donc une forte montée de \(P_{\mathrm{vir}}\), puis de \(P_{\mathrm{tot}}\).
\end{itemize}

Les diagnostics temporels montrent effectivement ce comportement. Pour des compressions géométriques modérées de l'ordre de quelques pourcents,
\begin{equation}
 \frac{\gamma_{\mathrm{geom}}}{\gamma_{\mathrm{target}}} > 1,
\end{equation}
tandis que la pression cinétique active reste presque constante et que la pression virielle augmente fortement, entraînant avec elle la pression totale active. Dans les exemples de référence, une compression géométrique de l'ordre de \(4\%\) conduit à une augmentation très marquée de \(P_{\mathrm{vir}}\), alors que \(P_{\mathrm{kin}}\) varie peu. Le piston est donc bien le cas où la fermeture liquide révèle le plus clairement sa faible compressibilité.

Le calcul du module volumique effectif
\begin{equation}
 B_{\mathrm{eff}} \approx \frac{\Delta P_{\mathrm{tot,actif}}}{-\Delta V_{\mathrm{actif}}/V_{\mathrm{actif}}(0)}
\end{equation}
et de la compressibilité effective \(\kappa_{\mathrm{eff}}=1/B_{\mathrm{eff}}\) fournit enfin un diagnostic direct de cette raideur. Les valeurs obtenues dans les tests de référence sont cohérentes avec la pente \(\mathrm d P_{\mathrm{tot}}/\mathrm d \rho\) reconstruite sur le domaine actif et confirment que le modèle se comporte comme un fluide très peu compressible. À mesure que la compression progresse, on s'attend naturellement à l'apparition d'un régime de saturation de redistribution : ce n'est pas un artefact, mais l'expression numérique du caractère bloquant d'un liquide quasi incompressible confiné dans un volume trop petit.


\section{Redistribution par zones : principe, découpage et double passe}

\subsection{Objectif de la branche par zones}

La redistribution par zones a été introduite comme une étape intermédiaire entre la redistribution de référence, entièrement séquentielle sur le domaine complet, et une future implémentation parallèle plus ambitieuse. L'idée directrice est de conserver l'opérateur physique actuel de redistribution, mais de l'appliquer localement sur des sous-domaines rectangulaires du maillage, de manière à tester la sensibilité de la physique numérique à un traitement ``simili-parallèle'' avant tout portage compilé.

Dans la branche actuelle, l'ordre général du \emph{step} n'est pas modifié : l'advection, les conditions limites, la collision SRD, la capture de l'état de référence, la redistribution, la fermeture liquide et les diagnostics conservent la même séquence que dans la version de référence. Seul l'opérateur de redistribution peut être remplacé par une variante ``par zones'', commandée par les paramètres \verb|useZoneRedistribution|, \verb|zoneTileNx|, \verb|zoneTileNy|, \verb|zoneUseShiftedSecondPass|, \verb|zoneStoreDiagnostics|, \verb|zonePassthroughMode| et \verb|zoneSingleFullDomainMode|.

Cette branche a été utilisée comme outil d'évaluation : il ne s'agit pas encore d'une redistribution parallèle complète, mais d'un cadre permettant de mesurer ce qui est préservé ou perdu lorsque l'on remplace un traitement global par une suite d'opérateurs locaux.

\subsection{Construction du découpage en zones}

Soit un maillage cartésien de taille $N_x \times N_y$. On choisit deux tailles de zones entières
\begin{equation}
N_x^{(z)} = \texttt{zoneTileNx},
\qquad
N_y^{(z)} = \texttt{zoneTileNy}.
\end{equation}
Le domaine est alors découpé en rectangles alignés sur la grille. Pour la passe dite \emph{de base}, les indices de départ sont
\begin{equation}
 i_x^{(0)} \in \{1,\,1+N_x^{(z)},\,1+2N_x^{(z)},\dots\},
 \qquad
 i_y^{(0)} \in \{1,\,1+N_y^{(z)},\,1+2N_y^{(z)},\dots\},
\end{equation}
chaque zone étant ensuite tronquée au bord du domaine si nécessaire. Si l'on note
\begin{equation}
\mathcal Z_{pq}^{\mathrm{base}} = \{(i,j)\, :\, i_{x,p}^{(0)} \le i \le i_{x,p}^{(1)},\; i_{y,q}^{(0)} \le j \le i_{y,q}^{(1)}\},
\end{equation}
alors chaque cellule appartient à une zone de cette famille de base.

Pour un domaine $50\times 50$ et des zones $25\times 25$, la passe de base produit quatre sous-domaines :
\begin{equation}
 x = 1{:}25 \;\text{ou}\; 26{:}50,
 \qquad
 y = 1{:}25 \;\text{ou}\; 26{:}50.
\end{equation}
L'interface principale du pavage est alors située entre les colonnes $25$ et $26$, et entre les lignes $25$ et $26$. Sur un domaine physique $[0,10]\times[0,10]$, avec $50$ cellules dans chaque direction, cela correspond à $x=5$ et $y=5$.

\subsection{Seconde passe décalée}

Lorsque \verb|zoneUseShiftedSecondPass=true|, une seconde famille de zones est construite avec un décalage d'environ une demi-zone dans chaque direction. Le décalage est pris entier :
\begin{equation}
 o_x = \left\lfloor \frac{N_x^{(z)}}{2} \right\rfloor,
 \qquad
 o_y = \left\lfloor \frac{N_y^{(z)}}{2} \right\rfloor.
\end{equation}
Les points de départ de la famille décalée deviennent alors
\begin{equation}
 i_x^{(s)} \in \{1\} \cup \{1+o_x,\,1+o_x+N_x^{(z)},\,1+o_x+2N_x^{(z)},\dots\},
\end{equation}
\begin{equation}
 i_y^{(s)} \in \{1\} \cup \{1+o_y,\,1+o_y+N_y^{(z)},\,1+o_y+2N_y^{(z)},\dots\}.
\end{equation}
La présence systématique du point de départ $1$ garantit une couverture correcte du bord même lorsque le décalage ne divise pas exactement le domaine.

Pour $N_x^{(z)}=N_y^{(z)}=25$ sur un domaine $50\times 50$, on obtient
\begin{equation}
 o_x=o_y=12,
\end{equation}
et les départs de la passe décalée sont
\begin{equation}
 x : 1,\;13,\;38,
 \qquad
 y : 1,\;13,\;38,
\end{equation}
ce qui produit des zones du type
\begin{equation}
1{:}25,
\qquad
13{:}37,
\qquad
38{:}50.
\end{equation}
La deuxième passe ne constitue donc pas un simple balayage dans l'ordre inverse ; elle reconstruit une nouvelle famille géométrique de zones, toujours rectangulaires et alignées sur la grille, mais décalées d'une demi-zone.

\subsection{Transferts strictement intra-zone}

Le point structurel le plus important est que chaque famille de zones est traitée séquentiellement, zone par zone, avec des transferts strictement intra-zone. Si l'on note $\mathcal Z$ la zone courante, les cellules sources denses et les cellules cibles admissibles sont filtrées par le masque de $\mathcal Z$ ; par conséquent, une zone ne peut pas, au cours d'une même passe, décharger directement un surplus dans une cellule située dans une autre zone, ni recevoir directement des particules d'une autre zone.

La seconde passe n'introduit donc pas une compensation interzone globale ; elle fournit seulement une deuxième occasion de redistribuer les particules dans une autre partition du domaine. Une cellule qui se trouvait près d'une frontière de la famille de base peut ainsi se retrouver au milieu d'une zone de la famille décalée, et être corrigée à ce moment-là. Cette correction reste cependant indirecte.

\subsection{Modes de contrôle et validation}

Deux modes auxiliaires ont été utiles pendant la mise au point :
\begin{itemize}
  \item \verb|zonePassthroughMode=true| : la branche zones est traversée, mais délègue entièrement à la redistribution de référence ; elle sert à vérifier que l'orchestrateur zones ne modifie pas la physique lorsqu'il est neutre.
  \item \verb|zoneSingleFullDomainMode=true| : une seule zone recouvre tout le domaine ; le chemin ``zones'' est alors exercé dans une configuration contrôlée, sans introduire de véritable découpage géométrique.
\end{itemize}
Ces deux étapes ont permis de vérifier que la structure \emph{runner -> simulateur -> step -> redistribution par zones} était correctement branchée avant de tester le vrai mode \verb|multi_zone_active|.

\section{Lissage interzone de $P_{\mathrm{drive}}$}

\subsection{Motivation}

Les premiers diagnostics sur le cas \texttt{compressibility} ont montré que la redistribution par zones ne détruisait pas fortement les champs moyens juste après la redistribution, mais qu'un défaut local persistait ensuite après le \emph{post-kick}. L'analyse détaillée du \emph{step} a montré que la fermeture liquide applique, après une correction de quantité de mouvement, un \emph{kick} de vitesse fondé sur le gradient de la pression motrice
\begin{equation}
 P_{\mathrm{drive}} = P_{\mathrm{kin}} + P_{\mathrm{vir}},
\end{equation}
avec
\begin{equation}
 P_{\mathrm{vir}} = K_{\mathrm{virial}}\,(\rho^{\mathrm{ref}} - \rho^{\star}).
\end{equation}
Le \emph{kick} vertical vaut alors
\begin{equation}
 \delta u_y^{\mathrm{EOS}} = -\beta_{\mathrm{EOS}}\,\frac{\Delta t}{\rho^{\star}}\,\partial_y P_{\mathrm{drive}}.
\end{equation}
En présence d'une petite discontinuité de $P_{\mathrm{drive}}$ à une interface de zones, le gradient discret peut amplifier localement cette discontinuité et créer un saut visible sur $u_y$.

\subsection{Principe du lissage}

Pour limiter cette amplification, un lissage local de $P_{\mathrm{drive}}$ a été introduit avant le calcul du gradient. L'idée est de ne lisser que dans une bande centrée sur les interfaces géométriques du zonage. Trois paramètres pilotent cette correction :
\begin{itemize}
  \item \verb|smoothPdriveInterzoneWidthCells| : demi-largeur de la bande de lissage, exprimée en nombre de cellules ;
  \item \verb|smoothPdriveInterzoneBlend| : coefficient de mélange entre le champ brut et le champ lissé ;
  \item \verb|smoothPdriveIncludeShiftedLayouts| : choix d'inclure ou non les interfaces de la famille décalée dans le masque de lissage.
\end{itemize}
Si l'on note $P_{\mathrm{drive}}^{\mathrm{raw}}$ le champ brut et $\widetilde P_{\mathrm{drive}}$ sa version lissée localement, la correction utilisée est du type
\begin{equation}
 P_{\mathrm{drive}}^{\mathrm{corr}}
 =
 (1-\lambda) P_{\mathrm{drive}}^{\mathrm{raw}}
 + \lambda \, \widetilde P_{\mathrm{drive}},
\end{equation}
avec
\begin{equation}
 \lambda = \texttt{smoothPdriveInterzoneBlend}.
\end{equation}
Le bulk éloigné des interfaces est laissé inchangé.

\subsection{Comportement observé}

Sur le cas \texttt{compressibility}, ce lissage améliore effectivement la régularité du saut interzone de $u_y$, mais ne le fait pas disparaître totalement. Les essais montrent un comportement non monotone :
\begin{itemize}
  \item une largeur de lissage de $1$ cellule améliore légèrement le défaut ;
  \item une largeur de $2$ cellules constitue le meilleur compromis testé ;
  \item une largeur de $3$ cellules sur-lisse le champ et dégrade de nouveau les profils moyens.
\end{itemize}
Le réglage le plus satisfaisant parmi ceux testés est donc apparu être, pour le cas hydrostatique, un lissage de largeur $2$ cellules avec un mélange proche de l'unité. Ce point doit être lu comme une correction ciblée des interfaces, et non comme un nouveau modèle physique autonome.

\section{Bilan des essais de la redistribution par zones}

\subsection{Cas \texttt{compressibility} : cas sévère pour le zonage}

Le cas \texttt{compressibility} a été l'outil principal pour révéler les limites de l'approche par zones. Il combine en effet trois ingrédients défavorables :
\begin{itemize}
  \item un gradient de pression essentiellement vertical ;
  \item une pression virielle très raide, donc sensible aux petites erreurs de densité ;
  \item un \emph{kick} vertical directement proportionnel à $\partial_y P_{\mathrm{drive}}$.
\end{itemize}
Dans ce contexte, la moindre discontinuité de redistribution aux interfaces horizontales du pavage se répercute dans le gradient de pression et apparaît comme une cassure de $u_y$ ou comme un pic dans les diagnostics hydrostatiques.

Les tests ont montré que :
\begin{enumerate}[label=\alph*)]
  \item avec de grandes zones ($25\times25$ sur $50\times50$), la seconde passe décalée atténue partiellement les défauts, mais ne supprime pas la cicatrice interzone ;
  \item avec de petites zones ($10\times10$), la situation se dégrade fortement : la combinaison du pavage de base et du pavage décalé introduit alors des interfaces horizontales tous les $5$ points environ, ce qui revient pratiquement à traiter l'hydrostatique par sous-couches horizontales ;
  \item le lissage de $P_{\mathrm{drive}}$ améliore la situation, mais n'efface pas totalement la stratification résiduelle.
\end{enumerate}
Il faut donc considérer \texttt{compressibility} comme un cas très sévère pour l'approche par zones : il est excellent pour révéler les défauts, mais probablement trop exigeant pour juger à lui seul de la pertinence de la méthode dans des cas d'usage plus courants.

\subsection{Cas Poiseuille : robustesse sur un écoulement cisaillé}

Le cas Poiseuille s'est révélé beaucoup plus favorable. Les profils moyens $U_x(y)$ obtenus avec la redistribution par zones restent très proches de la référence, et les quantités intégrées telles que le débit $Q$, la viscosité identifiée par courbure ou la viscosité de paroi restent dans la variabilité naturelle du bruit MPCD. Des tests répétés sur plusieurs graines aléatoires ont montré que l'écart entre référence et redistribution par zones, avec lissage interzone, est du même ordre que la dispersion statistique observée entre deux réalisations indépendantes de la méthode MPCD.

Ce résultat est important : il indique que, pour un écoulement de type canal où le \emph{kick} de pression n'est pas dominé par une hydrostatique très raide, la redistribution par zones ne modifie pas de manière manifeste les observables principales du problème.

\subsection{Cas piston : compression quasi incompressible en volume actif variable}

Le cas piston est lui aussi nettement plus rassurant que \texttt{compressibility}. La cinématique imposée par le piston, la compression géométrique, la densité active requise et les pressions moyennes du modèle restent très proches entre la référence et la version par zones. Les diagnostics de pression active, de module effectif et de redistribution évoluent dans la même enveloppe, ce qui suggère que la compression quasi incompressible elle-même n'est pas disqualifiée par le zonage.

En pratique, le cas piston confirme que le principal verrou de l'approche par zones n'est pas la seule présence d'une pression virielle, mais bien la combinaison particulière ``gradient hydrostatique vertical fort + interfaces horizontales répétées + dérivation de $P_{\mathrm{drive}}$'' qui caractérise le cas \texttt{compressibility}.

\section{Discussion sur l'intérêt de l'approche par zones pour une parallélisation future}

\subsection{Pourquoi l'approche par zones reste intéressante}

L'approche par zones a été conçue pour préparer une future parallélisation de l'opérateur de redistribution, sans toucher à l'advection, aux collisions, aux conditions limites, à la fermeture liquide ou à la réorientation interfaciale. De ce point de vue, elle remplit déjà plusieurs objectifs :
\begin{itemize}
  \item elle isole clairement l'opérateur de redistribution dans une branche contrôlée ;
  \item elle permet de comparer directement une redistribution globale de référence et une redistribution locale ``simili-parallèle'' ;
  \item elle fournit déjà les paramètres géométriques essentiels d'un futur découpage parallèle : taille de sous-domaines, familles de zones, diagnostics de passe.
\end{itemize}
Pour les cas usuels sans surface libre tels que \texttt{poiseuille}, \texttt{piston} et vraisemblablement \texttt{diffusion}, les résultats obtenus jusqu'à présent suggèrent que le coût de développement peut être justifié : les écarts avec la référence semblent du même ordre que la variabilité statistique naturelle de la méthode dans plusieurs observables pertinentes.

\subsection{Limites actuelles}

Trois limites majeures demeurent.

Premièrement, la compensation entre zones reste indirecte : une zone ne peut pas, au cours d'une même passe, être alimentée ou déchargée par une autre zone ; seule la seconde passe décalée offre une nouvelle occasion de corriger les défauts dans un autre pavage.

Deuxièmement, la seconde passe actuelle garde la même orientation géométrique que la première ; elle est simplement décalée d'une demi-zone. Pour des problèmes à fort gradient vertical, cela peut donc préserver, voire multiplier, les interfaces les plus nuisibles. Une perspective naturelle serait d'explorer une seconde passe réellement orientée différemment, par exemple sous forme de bandes verticales plutôt que de blocs rectangulaires de même orientation.

Troisièmement, le noyau multi-zones ne supporte actuellement que les cas sans surface libre et sans mouillage. Les géométries à interface libre (ellipse, barrage, mouillage explicite) exigeraient de réintroduire dans le noyau zones la logique bulk/interface déjà présente dans la redistribution de référence, ce qui constitue un chantier distinct.

\subsection{Jugement provisoire}

Le bilan actuel conduit à la lecture suivante :
\begin{itemize}
  \item l'approche par zones n'est pas encore suffisamment robuste pour servir de référence générale sur des cas hydrostatiques sévères de type \texttt{compressibility} ;
  \item en revanche, elle paraît suffisamment prometteuse sur des cas d'écoulement et de compression plus usuels pour justifier la poursuite du développement ;
  \item la suite la plus raisonnable n'est pas de raffiner indéfiniment le cas \texttt{compressibility}, mais d'évaluer l'intérêt du zonage sur les cas sans surface libre qui correspondent aux usages visés du code.
\end{itemize}


\section{Méthodes de projection incompressible Q6/Q9}

La chaîne la plus récente du dépôt ne repose plus sur la redistribution dure comme opérateur principal d'incompressibilité. Elle introduit une famille de méthodes particules--grille dont l'objectif est de conserver le caractère mésoscopique SRC/MPCD tout en supprimant les modes compressifs les plus nuisibles. Les deux briques essentielles sont :
\begin{itemize}
  \item Q6, qui projette le champ de vitesse cellulaire sur un champ discret quasi solénoïdal ;
  \item Q9, qui complète Q6 par une projection basse fréquence du flux de masse $Nu$, afin de réduire les grandes structures cohérentes de densité.
\end{itemize}

Ces méthodes sont moins intrusives que la redistribution : elles ne déplacent pas directement les particules d'une cellule vers une autre pour imposer une occupation cible. Elles modifient les vitesses particulaires après dépôt sur grille, résolution d'un problème de projection et interpolation de la correction vers les particules. Cette structure est aussi plus favorable à une future parallélisation.

\subsection{Dépôt particules--grille}

À chaque pas, après l'étape SRC/MPCD classique, les particules sont déposées sur une grille cartésienne. Pour une cellule $c$, on définit l'occupation
\begin{equation}
  N_c = \#\{i\in c\},
\end{equation}
la densité de nombre discrète
\begin{equation}
  n_c = \frac{N_c}{V_c},
\end{equation}
et la vitesse moyenne cellulaire
\begin{equation}
  \bm{u}_c = \frac{1}{N_c}\sum_{i\in c}\bm{v}_i,
  \qquad N_c>0.
\end{equation}
Dans les cellules vides, la vitesse est laissée nulle ou reconstruite par les conventions de l'opérateur de projection. Pour les géométries avec obstacles ou régions solides, un masque fluide/solide peut neutraliser les cellules non fluides dans les projections.

La grandeur de densité manipulée dans Q9 est généralement l'occupation $N_c$ plutôt que la masse volumique au sens strict. Lorsque $m=1$ et $V_c$ est constant, utiliser $N_c$, $n_c$ ou $\rho_c$ ne change que des facteurs de normalisation. Dans les expressions ci-dessous, on écrit donc indifféremment le flux de masse discret sous la forme
\begin{equation}
  \bm{M}_c = N_c\bm{u}_c
  \qquad \text{ou} \qquad
  \bm{M}_c = \rho_c\bm{u}_c,
\end{equation}
selon la convention de normalisation retenue.

\subsection{Projection de vitesse Q6 : contrainte $\nabla\cdot u\simeq 0$}

Q6 applique une projection de Helmholtz--Hodge discrète au champ cellulaire $\bm{u}^{\star}$ issu de l'étape SRC/MPCD. On cherche un potentiel scalaire $\phi$ tel que
\begin{equation}
  \bm{u}^{Q6} = \bm{u}^{\star} - \nabla_h \phi
\end{equation}
vérifie
\begin{equation}
  \nabla_h\cdot \bm{u}^{Q6} \simeq 0.
\end{equation}
On obtient donc l'équation de Poisson discrète
\begin{equation}
  \Delta_h \phi = \nabla_h\cdot \bm{u}^{\star}.
\end{equation}

Selon la géométrie, cette équation est résolue par une variante périodique FFT ou par un opérateur périodique en $x$ et borné en $y$. Le cas canal utilise une périodicité longitudinale et une condition compatible avec l'absence de flux normal aux parois. Le cas Taylor--Green utilise une projection entièrement périodique.

La correction de vitesse cellulaire est
\begin{equation}
  \delta\bm{u}_c^{Q6}=\bm{u}_c^{Q6}-\bm{u}_c^{\star},
\end{equation}
puis elle est interpolée vers les particules :
\begin{equation}
  \bm{v}_i \leftarrow \bm{v}_i + \lambda_u\,\delta\bm{u}^{Q6}(\bm{x}_i),
\end{equation}
où $\lambda_u=\verb|projectionStrength|$ est le paramètre de force de projection. Dans les validations de référence, on utilise généralement
\begin{equation}
  \lambda_u=1.
\end{equation}

Q6 doit être interprétée comme une correction cinématique. Elle supprime efficacement la divergence reconstruite de $\bm{u}$, mais elle ne garantit pas à elle seule que les modes basse fréquence de densité ne se développent pas, car la densité particulaire reste fluctuante.


\subsection{Opérateur elliptique \texttt{general\_bc} et fixation du mode constant}

La projection Q6 repose sur la résolution d'une équation de Poisson discrète ou, plus généralement, d'un problème elliptique compatible avec les conditions aux limites. La version récente utilise un opérateur \texttt{general\_bc} destiné à traiter dans un même cadre les directions périodiques, les murs et les géométries piston. Pour une correction de pression scalaire \(\phi\), on écrit schématiquement
\begin{equation}
  \Delta_h \phi = \frac{1}{\Delta t}\nabla_h\cdot\bm u^\star,
  \qquad
  \bm u^{n+1}=\bm u^\star-\Delta t\,\nabla_h\phi.
\end{equation}

Lorsque le problème possède des conditions de Neumann ou des directions périodiques, le mode constant de \(\phi\) n'est pas déterminé. La correction récente impose explicitement un \emph{gauge-fix} du mode constant nul. Cette opération ne change pas le gradient \(\nabla_h\phi\), donc ne change pas la correction de vitesse, mais elle rend le solveur bien posé numériquement et évite les singularités ou mauvais conditionnements artificiels lors des runs piston et canal.

\subsection{Projection du flux de masse Q9 : contrainte sur $\nabla\cdot(Nu)$}

Q9 ajoute à Q6 une correction du flux de masse. La motivation est que la contrainte
\begin{equation}
  \nabla\cdot \bm{u}=0
\end{equation}
reste insuffisante dans une méthode particulaire à occupation fluctuante : le transport de population est contrôlé par
\begin{equation}
  \nabla\cdot (N\bm{u}),
\end{equation}
qui peut rester compressif même si $\nabla\cdot\bm{u}$ est faible.

On définit donc le flux discret
\begin{equation}
  \bm{M}_c^{\star}=N_c\bm{u}_c^{Q6}.
\end{equation}
Q9 cherche une correction de type gradient
\begin{equation}
  \bm{M}_c^{Q9}=\bm{M}_c^{\star}-\nabla_h \psi_c
\end{equation}
telle que la divergence du flux relaxe les modes de densité visés :
\begin{equation}
  \nabla_h\cdot \bm{M}^{Q9}=S_N.
\end{equation}
Le second membre $S_N$ n'est pas une contrainte de densité locale dure. Il est construit à partir du défaut d'occupation filtré basse fréquence. Si
\begin{equation}
  \delta N_c = N_c-\gamma,
\end{equation}
et si $\mathcal{P}_{\mathrm{low}k}$ désigne le filtre spectral retenu, alors
\begin{equation}
  \delta N_c^{\mathrm{low}k}=\mathcal{P}_{\mathrm{low}k}(\delta N)_c.
\end{equation}
La cible de relaxation peut s'écrire, à un facteur de normalisation absorbé dans le paramètre numérique,
\begin{equation}
  S_{N,c} \simeq -\beta_N\,\delta N_c^{\mathrm{low}k},
\end{equation}
où
\begin{equation}
  \beta_N = \verb|massFluxDensityRelaxationBeta|.
\end{equation}
On obtient l'équation de Poisson
\begin{equation}
  \Delta_h\psi = \nabla_h\cdot \bm{M}^{\star}-S_N.
\end{equation}

Une fois le flux corrigé, la correction de vitesse associée est reconstruite par
\begin{equation}
  \delta\bm{u}_c^{Q9}=\frac{\bm{M}_c^{Q9}-\bm{M}_c^{\star}}{\max(N_c,N_{\min})},
\end{equation}
puis appliquée aux particules :
\begin{equation}
  \bm{v}_i\leftarrow \bm{v}_i + \lambda_M\,\delta\bm{u}^{Q9}(\bm{x}_i),
\end{equation}
avec
\begin{equation}
  \lambda_M=\verb|massFluxProjectionStrength|.
\end{equation}

Le choix central de Q9 est de ne corriger que les grandes longueurs d'onde de densité. Le filtre courant est
\begin{equation}
  \verb|massFluxTargetFilter|=\verb|lowpass_fft|,
  \qquad
  \verb|massFluxLowKMaxIndex|=2.
\end{equation}
Les modes locaux, assimilables au bruit cellulaire naturel du MPCD, ne sont donc pas forcés brutalement.

\subsection{Projection finale de nettoyage}

Après interpolation de la correction Q9 vers les particules, la reconstruction cellulaire peut réintroduire une divergence faible mais non nulle, en particulier lorsque l'interpolation et le dépôt ne sont pas des opérateurs adjoints exacts. Un nettoyage optionnel est donc appliqué :
\begin{equation}
  \verb|massFluxFinalVelocityProjectionCleanup|=\verb|true|.
\end{equation}
Il consiste à appliquer une nouvelle projection de vitesse de force réduite
\begin{equation}
  \lambda_{\mathrm{clean}}=\verb|massFluxFinalVelocityProjectionStrength|.
\end{equation}
Le compromis retenu dans les cas validés est généralement
\begin{equation}
  \lambda_{\mathrm{clean}}=0.5.
\end{equation}
Un nettoyage trop fort rend $\nabla\cdot u$ très petit, mais peut réduire l'effet du contrôle low-$k$ ; un nettoyage trop faible préserve mieux la correction de flux de masse, mais laisse une divergence reconstruite plus élevée.

\subsection{Paramètres de référence}

Les paramètres Q9 de référence, avant adaptation à certaines géométries, sont :
\begin{center}
\begin{tabular}{ll}
\toprule
Paramètre & Valeur typique \\
\midrule
\verb|projectionStrength| & $1.0$ \\
\verb|massFluxProjectionMode| & \verb|relax_to_uniform_lowk| \\
\verb|massFluxProjectionStrength| & $1.0$ \\
\verb|massFluxDensityRelaxationBeta| & $2\times10^{-3}$ \\
\verb|massFluxApplyAfterVelocityProjection| & \verb|true| \\
\verb|massFluxTargetFilter| & \verb|lowpass_fft| \\
\verb|massFluxLowKMaxIndex| & $2$ \\
\verb|massFluxFinalVelocityProjectionCleanup| & \verb|true| \\
\verb|massFluxFinalVelocityProjectionStrength| & $0.5$ \\
\bottomrule
\end{tabular}
\end{center}

Pour la marche descendante, le paramètre de relaxation a dû être réduit à
\begin{equation}
  \texttt{massFluxDensityRelaxationBeta}=5\times10^{-4},
\end{equation}
afin d'éviter une sur-correction au voisinage du coin solide. Cette adaptation est importante : Q9 possède un paramètre de raideur effectif, mais sa valeur optimale dépend de la géométrie et du niveau de bruit mécanique.

\section{Diagnostics utilisés pour Q6/Q9}

\subsection{Densité locale et densité basse fréquence}

La densité locale est suivie par l'écart-type de l'occupation cellulaire,
\begin{equation}
  \mathrm{std}(N)=\sqrt{\langle (N_c-\langle N\rangle)^2\rangle_c},
\end{equation}
et par la fraction de cellules hors bande
\begin{equation}
  N_c\notin \gamma(1\pm \varepsilon_N),
\end{equation}
où $\varepsilon_N$ vaut typiquement $0.2$ dans les validations.

Le diagnostic le plus spécifique de Q9 est l'énergie basse fréquence du défaut relatif de densité. En posant
\begin{equation}
  r_c = \frac{N_c}{\gamma}-1,
\end{equation}
et en notant $\widehat r_{k_x,k_y}$ sa transformée de Fourier discrète, on mesure une quantité du type
\begin{equation}
  E_{\mathrm{low}k}=\sum_{0<|k_x|,|k_y|\le K}\abs{\widehat r_{k_x,k_y}}^2,
\end{equation}
avec $K=\verb|lowKMaxIndex|$. La normalisation exacte dépend du script de diagnostic, mais les comparaisons Q6/Q9 utilisent toujours la même convention.

\subsection{Divergence et flux de masse}

Après projection, on reconstruit
\begin{equation}
  D_u = \left\|\nabla_h\cdot\bm{u}\right\|_{\mathrm{RMS}}
\end{equation}
et, pour Q9,
\begin{equation}
  D_M = \left\|\nabla_h\cdot(N\bm{u})-S_N\right\|_{\mathrm{RMS}}.
\end{equation}
On distingue le résidu de projection du flux de masse, calculé immédiatement au niveau grille, et le défaut reconstruit après retour aux particules. Cette distinction est essentielle : le solveur de projection peut être très précis sur grille, tandis que l'interpolation particulaire réintroduit une erreur effective.

\subsection{Diagnostics de structures}

Pour les cas vorticitaires, on utilise la vorticité discrète
\begin{equation}
  \omega = \partial_x u_y - \partial_y u_x
\end{equation}
et l'enstrophie moyenne
\begin{equation}
  \mathcal{E}_{\omega}=\langle \omega^2\rangle.
\end{equation}

Dans le cas Taylor--Green, on projette le champ de vitesse sur le mode analytique initial afin d'obtenir une amplitude de mode, une énergie modale et une cohérence. Dans le cas marche, on suit la longueur de recirculation, l'aire de recirculation, le RMS de vorticité dans la couche de cisaillement et des sondes aval en $\omega$ et $u_y$.

\subsection{Pression effective au piston}

Pour l'évaluation d'une équation d'état effective, trois pressions sont distinguées :
\begin{align}
  P_{\mathrm{kin}} &= \langle n_c k_B T_c\rangle,\\
  P_{\mathrm{wall}} &= \frac{\Delta p_y^{\mathrm{top}}}{\Delta t_{\mathrm{sample}} L_x},\\
  P_{\mathrm{excess}} &= P_{\mathrm{wall}}-P_{\mathrm{kin}}.
\end{align}
La pression mécanique au mur est estimée par l'impulsion cumulée échangée avec le piston pendant la fenêtre de mesure retenue. C'est cette pression qui permet de définir une raideur mécanique effective.

En posant
\begin{equation}
  \chi=\frac{\rho}{\rho_0}=\frac{y_{\mathrm{top},0}}{y_{\mathrm{top}}},
\end{equation}
on écrit localement
\begin{equation}
  P(\chi)=P_0+K_{\mathrm{eff}}\ln \chi,
\end{equation}
ou, pour de faibles compressions,
\begin{equation}
  \Delta P \simeq K_{\mathrm{eff}}(\chi-1).
\end{equation}


\section{Validations Q6/Q9 corrigées}

\subsection{Statut des anciennes validations}

Les résultats Q6/Q9 historiques du rapport précédent sont conservés comme jalons méthodologiques, mais ils ne doivent plus être utilisés seuls pour calibrer la méthode. Trois éléments ont changé : le thermostat cellulaire a été corrigé, les conditions limites de canal ont été renforcées par particules virtuelles \texttt{wallVP}, et plusieurs routines MATLAB ont été réorganisées/vectorisées. Les comparaisons qualitatives restent utiles : Q9 réduit les modes compressifs basse fréquence et ne détruit pas les structures moyennes principales. En revanche, les valeurs de viscosité effective, de raideur piston et de pression de paroi doivent être réévaluées avec les scripts récents.

\subsection{Taylor--Green périodique : préservation des structures lisses}

Le cas Taylor--Green reste le test le plus propre du bulk périodique, car il élimine les complications de paroi. La configuration de référence récente utilise typiquement un domaine \(64\times64\), une occupation \(\gamma=20\), \(\Delta t=10^{-3}\), \(k_BT=10^{-2}\) et une amplitude initiale de l'ordre de \(0.1\). Les runs forcés avec Q6/Q9 complet donnent une viscosité effective bulk de l'ordre de
\begin{equation}
  \nu_{\mathrm{eff,TG}} \simeq 2.0\times10^{-2},
\end{equation}
avec une température stabilisée par le thermostat corrigé et une diminution des composantes basse fréquence du champ de densité.

Le rôle de ce test n'est pas d'identifier une loi d'état, mais de vérifier que les corrections Q6/Q9 ne projettent pas abusivement la dynamique vorticitaire. Le résultat important est que l'amplitude du mode Taylor--Green, l'énergie modale et la cohérence de phase restent exploitables : Q9 agit sur les défauts compressifs cohérents plutôt que comme un filtre destructeur de vorticité.

\subsection{Poiseuille : viscosité effective, murs et particules virtuelles}

Poiseuille reste le test principal de canal. Dans les versions anciennes, les profils pouvaient être bons en apparence mais dépendaient fortement du traitement de paroi et du thermostat. La version récente impose donc une lecture plus stricte : un run Poiseuille valide doit fournir simultanément un profil moyen parabolique, une température contrôlée, une densité basse fréquence réduite par Q9, une estimation de viscosité robuste après exclusion des cellules proches des murs, et une réponse pariétale compatible avec la pression cinétique.

L'introduction des particules virtuelles de paroi modifie sensiblement l'interprétation. Le \emph{bounce-back} géométrique impose une contrainte d'imperméabilité et une forme de non-glissement, mais il ne suffit pas toujours à reconstruire correctement la contrainte collisionnelle au mur. Le modèle \texttt{wallVP} ajoute une population solide effective dans les cellules coupées ; il améliore le couplage fluide--paroi, mais impose de diagnostiquer la pression par impulsions cumulées, et non par échantillons instantanés isolés.

\subsection{Poiseuille multi-résolutions : viscosité brute et correction en taille de grille}

Une difficulté apparue dans les runs Poiseuille récents est la variation apparente de la viscosité identifiée lorsque l'on modifie le nombre de mailles dans la direction transverse. À paramètres MPCD identiques, la viscosité mesurée par fit quadratique du profil peut changer lorsque l'on diminue ou augmente \(N_y\), même lorsque le profil reste visuellement parabolique. Cette observation ne doit pas être interprétée immédiatement comme une modification physique du fluide par Q6/Q9 : une part importante vient de la manière dont le profil discret est ajusté sur une hauteur de canal représentée par un nombre fini de cellules.

Le point de départ est le profil continu de Poiseuille,
\begin{equation}
  u_x(y)=\frac{f_x}{2\nu} y(H-y),
\end{equation}
où \(H\) est la hauteur utile du canal. Si le fit est effectué dans une coordonnée cellulaire normalisée de manière implicite par \(N_y\), la courbure discrète porte un facteur d'échelle en \(H^2\), donc en première approximation en \(N_y^2\) lorsque \(L_y\) et le domaine physique de référence sont conservés. Une viscosité brute extraite directement de la courbure discrète peut donc suivre une loi apparente
\begin{equation}
  \nu_{\mathrm{raw}}(N_y) \propto \frac{1}{N_y^2},
\end{equation}
si l'on ne réintroduit pas correctement l'échelle de hauteur du canal dans l'identification.

La correction utilisée dans les campagnes récentes consiste à rapporter les viscosités à une résolution transverse de référence \(N_{y,\mathrm{ref}}=64\), selon
\begin{equation}
  \nu_{\mathrm{corr}}
  =
  \nu_{\mathrm{raw}}
  \left(\frac{N_y}{N_{y,\mathrm{ref}}}\right)^2.
  \label{eq:nu-corr-ny}
\end{equation}
Cette écriture est volontairement pratique : elle ne prétend pas remplacer une dérivation analytique complète de la viscosité MPCD avec parois, mais elle corrige le biais de lecture introduit par la discrétisation transverse du profil.

Les campagnes comparant plusieurs grilles, notamment des canaux de type \(32\times48\), \(48\times64\) et \(64\times96\), ont montré que les viscosités brutes issues du fit pouvaient varier fortement avec \(N_y\), alors que les viscosités corrigées par \eqref{eq:nu-corr-ny} se regroupaient autour de la viscosité bulk de référence obtenue sur Taylor--Green, typiquement
\begin{equation}
  \nu_{\mathrm{eff,TG}} \simeq 0.02.
\end{equation}
Dans ces runs, les valeurs corrigées se plaçaient typiquement dans une bande de l'ordre de \(0.020\)--\(0.023\) pour Q9 et de \(0.020\)--\(0.021\) pour le cas classic, hors cas manifestement insuffisamment convergé ou mal ajusté. Ce regroupement est un résultat important : la diminution de \(N_y\) ne signifie pas nécessairement que la méthode change de fluide ; elle peut surtout changer l'échelle effective de lecture du profil de Poiseuille.

La conséquence méthodologique est que tout tableau Poiseuille doit désormais distinguer explicitement :
\begin{enumerate}[label=\roman*)]
  \item la viscosité brute de fit, \(\nu_{\mathrm{raw}}\), directement extraite du profil discret ;
  \item la résolution transverse \(N_y\), la hauteur utile \(H\), et le nombre de cellules exclues près des murs ;
  \item la viscosité corrigée \(\nu_{\mathrm{corr}}\), calculée avec \eqref{eq:nu-corr-ny} ou, plus généralement, avec le facteur \((H/H_{\mathrm{ref}})^2\) si la hauteur physique effective varie ;
  \item la comparaison avec la référence bulk périodique Taylor--Green, qui sert de contrôle indépendant des effets de paroi.
\end{enumerate}

Cette correction ne dispense pas de traiter correctement les conditions limites. Les particules virtuelles \texttt{wallVP}, l'exclusion de quelques cellules adjacentes aux murs et la durée du moyennage restent nécessaires pour obtenir un profil exploitable. La correction en \(N_y^2\) répond à un problème différent : elle évite de confondre une erreur d'échelle géométrique dans l'identification de \(\nu\) avec une vraie modification de transport due à Q6/Q9. Les prochains rapports de validation devront donc toujours présenter les deux colonnes \(\nu_{\mathrm{raw}}\) et \(\nu_{\mathrm{corr}}\), et ne comparer à Taylor--Green que la viscosité corrigée.

\subsection{Piston staircase : EOS bulk Q6/Q9 avec fermeture virielle}

Le piston staircase récent constitue le résultat bulk le plus net. Dans le run \texttt{Q9\_FULL} à \(N_p=20480\), grille \(32\times32\), \(\gamma=20\), \(\Delta t=10^{-3}\), \(k_BT=10^{-2}\), compression finale de \(5\%\), la densité physique finale est \(2.1558\times10^4\), la température cellule finale vaut \(0.01\), la densité basse fréquence finale est de l'ordre de \(10^{-7}\), et le limiteur viriel n'est pas déclenché dans le régime doux testé.

Le résultat essentiel est l'EOS bulk. Les fits staircase donnent
\begin{align}
  \frac{\mathrm d P_{\mathrm{kin}}}{\mathrm d\rho} &\simeq 0.010534, & R^2&=1,\\
  \frac{\mathrm d P_{\mathrm{vir}}}{\mathrm d\rho} &\simeq 0.05, & R^2&=1,\\
  \frac{\mathrm d P_{\mathrm{tot}}}{\mathrm d\rho} &\simeq 0.060534, & R^2&=1,\\
  \frac{\mathrm d P_{\mathrm{drive}}}{\mathrm d\rho} &\simeq 0.060534, & R^2&=1.
\end{align}
La valeur finale correspond à
\begin{equation}
  P_{\mathrm{kin}}\simeq 226.95,
  \qquad
  P_{\mathrm{vir}}\simeq 53.89,
  \qquad
  P_{\mathrm{tot}}\simeq 280.84.
\end{equation}
Ce résultat valide la fermeture virielle bulk : le code reproduit exactement la pente imposée par \(K_{\mathrm{virial}}\), et la pression totale suit la somme \(P_{\mathrm{kin}}+P_{\mathrm{vir}}\). Le résidu global de quantité de mouvement du kick viriel reste au niveau du bruit machine et le limiteur de kick n'est pas activé dans ce cas.

\subsection{Validation séparée du diagnostic mécanique \texttt{wallVP}}

Le même run piston a montré que les pressions instantanées de paroi pouvaient être algébriquement cohérentes mais physiquement inutilisables sur des plateaux courts : les résidus de reconstruction \(P_{\mathrm{total}}=P_{\mathrm{impact}}+P_{\mathrm{VP}}\) étaient nuls, mais les moyennes par plateau oscillaient fortement. Cette observation a conduit à isoler un test plus simple : gaz à l'équilibre entre deux murs, sans piston et sans terme viriel.

Dans ce test d'équilibre \(32\times32\), \(N_p=20480\), \(k_BT=0.01\), la pression idéale nominale vaut \(\rho k_BT=204.8\), tandis que la pression cinétique bulk mesurée vaut environ \(215.6\). Sans \texttt{wallVP}, les pressions cumulées haut/bas valent environ \(208.48\) et \(207.53\). Avec \texttt{wallVP}, les contributions impact seules montent à environ \(226.92\) et \(223.67\), les contributions VP signées valent environ \(-12.28\) et \(-9.07\), et les totaux cumulés reviennent à \(214.64\) et \(214.60\). L'écart relatif top/bottom tombe à environ \(2\times10^{-4}\).

Ce test valide la convention de signe, la normalisation \(\Delta p/(\Delta t L_x)\) et l'usage des impulsions cumulées. Il montre aussi que la contribution VP ne doit pas être rendue positive par construction : elle peut corriger à la baisse une contribution impact trop grande. La pression physique de paroi est le total signé impact+VP, moyenné sur une fenêtre longue.

\subsection{Retour au staircase : prudence sur les fenêtres courtes}

Les premiers retours staircase après correction du diagnostic montrent encore des moyennes de fenêtre instables lorsque les plateaux sont courts. Ce résultat n'invalide pas \texttt{wallVP} : le test d'équilibre a précisément montré que les pressions instantanées peuvent fluctuer entre des valeurs de plusieurs milliers d'unités autour d'une moyenne de l'ordre de \(2.1\times10^2\). La règle de validation actuelle est donc la suivante : comparer \(P_{\mathrm{tot,bulk}}\) aux moyennes de pression pariétale obtenues par différence d'impulsions cumulées sur des fenêtres longues, et non aux moyennes d'échantillons instantanés peu nombreux.


\section{Équation d'état bulk et statut mécanique de paroi}

\subsection{EOS bulk actuellement validée}

L'état validé côté piston doit être formulé de manière plus stricte que dans les versions précédentes du rapport. La loi d'état actuellement fiable est une loi bulk mesurée sur les cellules actives, pas encore une loi mécanique pariétale complète. Elle repose sur
\begin{equation}
  P_{\mathrm{vir}} = K_{\mathrm{virial}}(\rho-\rho_{\mathrm{EOSRef}}),
  \qquad
  P_{\mathrm{tot}} = P_{\mathrm{kin}}+P_{\mathrm{vir}},
  \qquad
  P_{\mathrm{drive}}=P_{\mathrm{tot}}.
\end{equation}

Le protocole staircase impose une suite de compressions à nombre de particules constant et hauteur active variable. Pour chaque plateau, on élimine le transitoire de déplacement, puis on moyenne les grandeurs bulk sur la partie finale du plateau. Les fits récents montrent que \(P_{\mathrm{kin}}\), \(P_{\mathrm{vir}}\), \(P_{\mathrm{tot}}\) et \(P_{\mathrm{drive}}\) sont linéaires en densité moyenne avec \(R^2\simeq1\). La pente de \(P_{\mathrm{vir}}\) retrouve directement \(K_{\mathrm{virial}}\), tandis que la pente de \(P_{\mathrm{tot}}\) est la somme de la pente cinétique effective et de la pente virielle imposée.

Cette validation est importante parce qu'elle dissocie deux questions. La première est : le modèle bulk produit-il une réponse pression--densité contrôlée ? La réponse est oui dans le régime doux testé. La seconde est : la pression mécanique mesurée sur le piston restitue-t-elle immédiatement cette même loi ? La réponse est plus prudente : pas sur des plateaux courts, car l'estimateur de paroi est fortement bruité.

\subsection{Kick viriel et conservation de la quantité de mouvement}

Le kick viriel actif est appliqué sous une forme faible et limitée. La version actuelle diagnostique le RMS de correction de vitesse, le rapport à la vitesse thermique, la fraction de cellules limitées, le maximum autorisé et le résidu de quantité de mouvement après correction globale. Dans le run de référence récent, le limiteur n'est pas déclenché, le RMS du kick reste de l'ordre de \(10^{-5}\), et le résidu global de moment reste proche de \(10^{-15}\). Ces valeurs indiquent que le terme viriel modifie la réponse EOS sans introduire de dérive globale détectable de la quantité de mouvement.

\subsection{Statut du diagnostic piston mécanique}

Le diagnostic mécanique piston n'est pas encore considéré comme validé pour les runs staircase courts. Le point acquis est algébrique : les champs de diagnostic reconstruisent correctement
\begin{equation}
  P_{\mathrm{wall,total}} = P_{\mathrm{wall,impact}} + P_{\mathrm{wall,VP}},
\end{equation}
avec un résidu de cohérence nul. Le point encore en cours de validation est statistique et physique : sur une fenêtre de plateau courte, le total instantané peut changer de signe ou prendre des valeurs trop grandes, alors que le bulk suit parfaitement l'EOS. La correction adoptée consiste à stocker les impulsions cumulées et à utiliser des moyennes de fenêtre longues.

La lecture correcte des prochains runs staircase sera donc :
\begin{enumerate}[label=\arabic*)]
  \item vérifier d'abord \(P_{\mathrm{kin}}\), \(P_{\mathrm{vir}}\), \(P_{\mathrm{tot}}\) et \(P_{\mathrm{drive}}\) dans le bulk ;
  \item vérifier que top et bottom donnent des pressions mécaniques cumulées compatibles à compression nulle ;
  \item n'interpréter les plateaux comprimés qu'après des temps de maintien longs ;
  \item comparer séparément la pression mécanique cumulée à \(P_{\mathrm{kin}}\) et à \(P_{\mathrm{tot}}\), afin de déterminer si le terme viriel bulk est effectivement transmis mécaniquement au piston ou seulement utilisé comme pression motrice interne.
\end{enumerate}

\section{Parallélisation et vectorisation MATLAB}

\subsection{Objectif}

La motivation initiale de Q6/Q9 est aussi algorithmique : remplacer une redistribution particulaire dure, difficile à paralléliser proprement, par des opérations plus régulières. Les briques principales de Q6/Q9 sont favorables à une exécution parallèle ou vectorisée : dépôt particules--grille, sommes par cellule, reconstruction de champs, gradients, projections elliptique/FFT, interpolation grille--particules et diagnostics.

\subsection{Ce qui est déjà favorable dans MATLAB}

La version MATLAB récente a commencé à réduire le coût de plusieurs boucles critiques par des opérations groupées. Le thermostat corrigé et les reconstructions cellulaires utilisent des index de cellule et des accumulations vectorisées lorsque c'est possible. Les diagnostics de densité, température, flux de masse, moyenne de profil, pression bulk et impulsions cumulées sont eux aussi structurés comme opérations sur tableaux.

Le gain n'est pas seulement un gain de temps. Une écriture vectorisée rend les diagnostics plus reproductibles : le même champ cellulaire sert à calculer la température, la pression, la divergence, les corrections Q6/Q9 et les métriques de conservation. Cela limite les divergences entre scripts de validation.

\subsection{Limites de MATLAB et direction C++/OpenMP}

MATLAB reste toutefois limitant pour les longs runs piston, les sweeps multi-graines et les géométries plus sévères. Les opérations particulaires avec conditions limites complexes, particules virtuelles, interpolation et stockage d'historiques deviennent coûteuses. La stratégie rationnelle est donc : conserver MATLAB comme banc de validation et de prototypage, puis porter les briques stabilisées vers C++/OpenMP, voire MPI/OpenMP ou GPU.

Les briques prioritaires pour un portage sont :
\begin{enumerate}[label=\arabic*)]
  \item collision SRC/MPCD avec thermostat corrigé ;
  \item conditions limites bounce-back, piston et \texttt{wallVP} ;
  \item projection Q6 avec opérateur \texttt{general\_bc} ;
  \item correction Q9 basse fréquence du flux de masse ;
  \item diagnostics cumulés de moment et de pression de paroi ;
  \item sorties légères pour profils moyens, vorticité, densité basse fréquence et EOS piston.
\end{enumerate}


\section{Synthèse actuelle et conséquences pour la parallélisation}

Le bilan actuel est plus clair que dans la version précédente du rapport. La chaîne historique par redistribution dure a démontré qu'il était possible d'imposer une faible compressibilité effective, mais elle est intrusive et difficile à paralléliser. La chaîne Q6/Q9 corrigée est moins intrusive : elle agit sur les champs projetés, réduit les modes compressifs cohérents et préserve mieux la structure particulaire naturelle du MPCD.

Les points considérés comme validés sont :
\begin{itemize}
  \item la projection Q6 avec opérateur compatible conditions limites et mode constant fixé ;
  \item la correction Q9 basse fréquence du flux de masse ;
  \item le thermostat cellulaire corrigé ;
  \item la stabilité thermique dans les cas de référence ;
  \item la réduction des modes basse fréquence de densité ;
  \item la préservation qualitative et quantitative des structures Taylor--Green ;
  \item la cohérence bulk de l'EOS piston virielle avec \(R^2\simeq1\) ;
  \item la cohérence du diagnostic \texttt{wallVP} en gaz d'équilibre lorsqu'il est évalué par impulsions cumulées.
\end{itemize}

Les points encore ouverts sont :
\begin{itemize}
  \item la statistique des pressions mécaniques de paroi dans les runs staircase comprimés ;
  \item la calibration multi-graines de la viscosité effective en canal avec \texttt{wallVP}, en reportant systématiquement \(\nu_{\mathrm{raw}}\) et \(\nu_{\mathrm{corr}}=\nu_{\mathrm{raw}}(N_y/N_{y,\mathrm{ref}})^2\) ;
  \item la transmission mécanique exacte de la pression virielle bulk vers le piston ;
  \item la robustesse sur écoulements plus complexes, notamment von Kármán et géométries à interfaces libres ;
  \item le coût des longs runs MATLAB, qui justifie un portage parallèle.
\end{itemize}

La prochaine étape scientifique n'est donc pas de modifier à nouveau le bulk Q6/Q9, mais de consolider la statistique de paroi : équilibre gaz long, Poiseuille long avec \texttt{wallVP}, puis staircase piston à plateaux longs. La prochaine étape numérique est de factoriser les briques validées pour préparer un portage C++/OpenMP.




\section{Addendum C++/OpenMP : inlet/outlet, solides immergés et cellules pauvres}

Les développements ultérieurs au socle MATLAB ont été poursuivis dans le dépôt C++/OpenMP \texttt{SRC\_incompressible\_openMP}. Ils ont d'abord concerné les conditions limites inlet/outlet Q6/Q9, puis le traitement des solides immergés fixes. Cette section met à jour le statut méthodologique du rapport : les solides immergés sont désormais séparés du problème plus général des cellules pauvres.

\subsection{Conditions limites inlet/outlet retenues}

La politique inlet/outlet retenue pour les cas ouverts est la suivante. L'inlet est alimenté par un réservoir de type \texttt{hard\_cell\_density}, de manière à imposer une occupation cible. L'outlet peut fonctionner en mode \texttt{balanced\_flux}, \texttt{neumann} ou \texttt{hybrid}. Le mode \texttt{hybrid} est devenu le choix nominal pour les cas d'injection plus difficiles : il part d'un flux local de type Neumann, ajoute éventuellement un faible mélange vers un profil équilibré, et applique une correction de balance globale faible au niveau de l'outlet.

Un point important des validations inlet/outlet est l'abandon des exclusions volumétriques près des frontières ouvertes comme stratégie nominale. Les réglages validés sont
\[
  \texttt{q9OpenBoundaryExclusionCells}=0,\qquad
  \texttt{virialOpenBoundaryExclusionCells}=0.
\]
Les exclusions non nulles créaient une interface active/inactive se comportant comme une paroi numérique près de l'outlet. La rampe d'inlet doit en outre être appliquée de façon cohérente à l'injection particulaire, au flux Q6 et au flux massique Q9.

Les apertures segmentées inlet/outlet sont maintenant interprétées comme des géométries physiques, par exemple une fente ou une buse, et non comme un correctif numérique destiné à stabiliser un canal canonique. Les diagnostics Q9 spatialisés sont écrits sous forme de sidecars binaires \texttt{q9\_diagnostics\_step\_*.q9bin}, lisibles par les outils MATLAB de visualisation.

\subsection{Formulation cellule/face pour les solides immergés}

Les essais précédents avaient montré que désactiver Q9 dans un halo autour d'un obstacle produisait une couche morte : les cellules adjacentes au solide devenaient incapables de produire une correction de pression ou de flux massique, précisément là où les accumulations de densité apparaissent. La formulation retenue est donc une formulation cellule/face.

On distingue :
\begin{itemize}
  \item les cellules solides, inactives pour Q6, Q9 et le viriel ;
  \item les cellules fluides, actives même lorsqu'elles sont adjacentes au solide ;
  \item les cellules coupées, actives si leur fraction fluide dépasse le seuil fixé ;
  \item les faces fluide--fluide, ouvertes au flux correctif ;
  \item les faces fluide--solide ou coupées fermées, sur lesquelles le flux normal correctif est imposé nul.
\end{itemize}

Le réglage nominal est donc
\[
  \texttt{projectionImmersedSolidMaskEnable}=\texttt{true},\quad
  \texttt{projectionImmersedSolidCloseCutFaces}=\texttt{true},\quad
  \texttt{q9ImmersedSolidHaloCells}=0.
\]
Le paramètre de halo Q9 reste éventuellement disponible comme mode conservateur ou de débogage, mais il ne doit plus être considéré comme une stratégie physique.

\subsection{Fermeture géométrique des flux Q6/Q9 aux faces solides}

Les premiers tests avec cercle immergé ont révélé que l'opérateur elliptique fermait bien les faces solides, mais que le diagnostic de fuite Q9 pouvait rester non nul lorsque la face fermée était portée par une cellule propriétaire solide. Une fermeture géométrique explicite a donc été ajoutée après reconstruction des flux :
\[
  F_{\mathrm{proj}}^{Q9}(\text{face solide})=0,\qquad
  F_{\mathrm{corr}}^{Q9}(\text{face solide})=-F_{\mathrm{base}}(\text{face solide}).
\]
La même convention a ensuite été appliquée au diagnostic Q6. Les grandeurs de fuite effectives
\[
  \texttt{q6ImmersedSolidLeakProjectedFluxRms},\qquad
  \texttt{q9ImmersedSolidLeakMassFluxRms}
\]
désignent désormais les flux après fermeture géométrique. Les diagnostics avant fermeture restent disponibles pour vérifier l'amplitude de la correction imposée.

Cette correction a été validée sur cercle périodique, cylindre périodique et cylindre en canal ouvert. Les résultats caractéristiques sont :
\[
  \texttt{q6ImmersedSolidLeakProjectedFluxRms}=0,\qquad
  \texttt{q9ImmersedSolidLeakMassFluxRms}=0,\qquad
  \texttt{q9ImmersedHaloExcludedCells}=0.
\]
Les cellules adjacentes et coupées restent actives, et les sidecars Q9 confirment l'absence de couronne inactive autour du solide.

\subsection{Viriel près des solides}

Le kick viriel a été modifié pour ne pas pousser artificiellement dans le solide. Dans les cellules fluides adjacentes au solide, le viriel reste actif, mais la composante normale du kick orientée vers une face solide fermée est neutralisée. Les composantes tangentielles sont conservées. Les diagnostics ajoutés comptent notamment les cellules adjacentes actives, le nombre de composantes normales supprimées et les normes du kick supprimé.

La validation périodique avec cylindre montre que le viriel reste actif au voisinage du solide sans créer de fuite de flux normal. La formulation actuelle est donc :
\begin{itemize}
  \item viriel actif en cellule fluide adjacente ;
  \item composante normale entrante dans le solide supprimée ;
  \item pas de halo viriel mort autour du solide ;
  \item cellules coupées traitées par fraction fluide.
\end{itemize}

\subsection{Validation visuelle des solides immergés}

La validation s'est faite par paliers afin de limiter les coûts :
\begin{enumerate}
  \item smoke-tests courts avec cercle ou cylindre périodique ;
  \item visualisation du démarrage sur quelques centaines de pas ;
  \item run périodique plus long avec forçage modéré ;
  \item premier cas inlet/outlet avec cylindre fixe.
\end{enumerate}

Les visualisations périodiques montrent un sillage cohérent, un déficit de vitesse en aval du cylindre et l'absence de couche inactive autour de l'obstacle. En inlet/outlet à vitesse modérée, les diagnostics restent propres : fuite Q6/Q9 solide nulle, halo Q9 nul, cellules adjacentes actives, limiteur Q9 peu ou pas sollicité. À ce stade, le traitement des solides immergés fixes peut donc être considéré comme validé pour les besoins du développement suivant.

\section{Limite identifiée : cellules pauvres et défauts de population}

Les essais plus sévères ont mis en évidence une difficulté indépendante de la présence d'un solide. Dans des cas ouverts ou avec poche initiale de cellules vides, certaines cellules pauvres peuvent se comporter comme des défauts persistants de population. L'effet est particulièrement visible lorsque l'on introduit une poche vide \(6\times6\) dans un canal ouvert sans solide.

Le témoin SRC classique convecte et diffuse progressivement la poche. En revanche, les variantes Q9/viriel enrichies par une garde locale de population ont tendance à faire l'inverse : la poche devient un défaut eulérien actif, dévie le champ de vitesse et se comporte comme un obstacle mou ou un puits fixé à la grille.

\subsection{Essais de garde locale de population}

Une première idée a consisté à écrire
\[
  P_{\mathrm{vir}} = P_{\mathrm{bulk}}(\rho) + P_{\mathrm{guard}}\left(\frac{N}{N_\star}\right),
  \qquad
  N_\star=\gamma_{\mathrm{ref}}\,\theta_f,
\]
où \(\theta_f\) est la fraction fluide de la cellule. Le terme \(P_{\mathrm{guard}}\) devait être nul dans une bande \([r_{\mathrm{low}},r_{\mathrm{high}}]\), positif dans les cellules surpeuplées et négatif dans les cellules sous-peuplées. L'objectif était de créer une pression locale dense--vers--pauvre sans redistribution dure de particules.

Cette approche a été testée sous plusieurs formes :
\begin{itemize}
  \item mode \texttt{pressure\_gradient}, injectant \(P_{\mathrm{guard}}\) dans le gradient cell-centered existant ;
  \item mode \texttt{face\_donor}, appliquant le kick depuis la cellule donneuse vers la cellule receveuse face par face ;
  \item mode \texttt{face\_donor\_surplus}, n'autorisant que les cellules réellement surpeuplées à donner ;
  \item mode \texttt{face\_donor\_capacity}, autorisant les cellules normales à donner jusqu'à un plancher de population.
\end{itemize}

Ces essais ont clarifié le mécanisme mais n'ont pas fourni de solution robuste. Le mode \texttt{face\_donor} simple peut transformer une cellule vide en puits aspirant ses voisines. Le mode \texttt{face\_donor\_surplus} évite cette aspiration mais laisse la poche vide agir comme un obstacle mou. Le mode \texttt{face\_donor\_capacity}, intermédiaire, n'a pas supprimé ce comportement et a même accentué l'effet dans certains cas.

\subsection{Interprétation}

Le résultat principal est le suivant : une cellule pauvre ne doit pas devenir une source de pression attractive. Une accumulation peut légitimement produire une pression répulsive, mais un vide de particules ne doit pas imposer un champ de vitesse eulérien fort. Dans la méthode particulaire, un trou local doit plutôt être convecté, diffusé et progressivement réparé par le transport SRC, l'inlet ou les voisins ; le transformer en objet actif le rend plus persistant.

La difficulté vient du couplage entre trois faits :
\begin{enumerate}
  \item Q9 corrige un flux de masse mais ne crée pas de porteurs dans une cellule vide ;
  \item les sécurités low-mass protègent les cellules peu peuplées, mais peuvent les rendre passives ;
  \item une garde attractive locale peut fixer un défaut de population sur la grille.
\end{enumerate}

La conclusion provisoire est donc de ne plus poursuivre les gardes locales attractives dense--vers--pauvre sous cette forme. Le problème des cellules pauvres doit être repris dans une branche dédiée, séparée de la branche des solides immergés, en partant du constat que le SRC classique transporte correctement les poches pauvres alors que certaines corrections Q9/viriel les eulérianisent.

\subsection{Pistes pour la branche suivante}

Les pistes les plus raisonnables pour la branche suivante sont :
\begin{itemize}
  \item rendre les cellules low-mass transparentes pour Q9, afin qu'elles ne deviennent pas des murs numériques ;
  \item éviter tout terme de pression négative attaché aux cellules pauvres ;
  \item tester une pression de garde positive seulement, agissant uniquement sur les accumulations ;
  \item séparer correction des accumulations et transport des trous ;
  \item examiner si la logique \texttt{ramp\_floor} de Q9 doit être remplacée par un mode \texttt{transparent\_floor} ;
  \item conserver la chaîne classic comme témoin systématique pour toute poche de cellules pauvres.
\end{itemize}

Cette difficulté ne remet pas en cause la validation des solides immergés. Elle identifie un problème plus général de robustesse locale de population dans les cas ouverts ou fortement hétérogènes.


\section{Branche resampling pondéré : réponse aux cellules pauvres et domaine fluide évolutif}

Les derniers développements MATLAB ont repris le problème des cellules pauvres sous un angle différent. Les essais Q9/viriel/guard avaient pour point commun d'essayer de réparer un défaut de population par une correction du champ de vitesse, du flux de masse ou de la pression. Or les tests avec poches vides ont montré que cette approche devient instable conceptuellement lorsque la cellule contient trop peu de particules : le champ hydrodynamique local n'est plus fiable, mais la correction eulérienne continue à l'interpréter comme un degré de liberté robuste. La poche pauvre peut alors devenir un objet actif de grille, au lieu d'être simplement transportée et diffusée par la dynamique particulaire.

La nouvelle branche, créée à partir de la base MATLAB Q6/Q9 intégrée, introduit une formulation \textbf{SRC/MPCD pondérée avec resampling local conservatif}. L'objectif n'est plus de forcer une cellule pauvre par une pression attractive, mais de maintenir un support particulaire local suffisant pour que les champs $\rho$, $u$ et $\rho u$ restent représentatifs. Cette section décrit le principe, les briques algorithmiques ajoutées et les premiers résultats obtenus.

\subsection{Difficulté fondamentale des cellules pauvres}

La projection Q6 porte sur la contrainte cinématique
\begin{equation}
  \nabla\cdot u \simeq 0,
\end{equation}
et la correction Q9 sur une contrainte de flux massique
\begin{equation}
  \nabla\cdot(\rho u) \simeq 0.
\end{equation}
Ces contraintes ont un sens hydrodynamique fort lorsque les champs sont reconstruits à partir d'une population statistiquement suffisante. En revanche, dans une cellule vide ou presque vide, la vitesse moyenne
\begin{equation}
  u_c = \frac{\sum_{i\in c} m_i v_i}{\sum_{i\in c} m_i}
\end{equation}
devient mal définie ou dominée par très peu de particules. La difficulté n'est donc pas seulement numérique : l'opérateur agit alors sur une variable hydrodynamique qui n'est plus correctement représentée.

Les essais de garde locale ont clarifié la contradiction. Une correction assez faible laisse le trou persister. Une correction assez forte pour remplir rapidement le trou agit comme une force interne artificielle. Les variantes \texttt{pressure\_gradient}, \texttt{face\_donor}, \texttt{face\_donor\_surplus} et \texttt{face\_donor\_capacity} ont donc produit soit des oscillations, soit une aspiration des voisins, soit un obstacle mou. Le principe retenu est dès lors le suivant :
\begin{quote}
Une cellule pauvre ne doit pas être corrigée par une attraction eulérienne. Elle doit être traitée comme une insuffisance locale du support particulaire.
\end{quote}

\subsection{Passage à des particules pondérées}

La branche resampling introduit un état particulaire de la forme
\begin{equation}
  \{x_i, v_i, m_i\}_{i=1}^{N_{\mathrm{cap}}},
\end{equation}
où $m_i$ est désormais une masse particulaire variable et $N_{\mathrm{cap}}$ la capacité mémoire préallouée. Le cas $m_i=m_0$ pour toutes les particules doit reproduire le comportement SRC/MPCD classique pondéré de façon neutre.

Les dépôts particules--grille sont remplacés par des dépôts massiques :
\begin{equation}
  M_c = \sum_{i\in c} m_i,
  \qquad
  P_c = \sum_{i\in c} m_i v_i,
  \qquad
  U_c = \frac{P_c}{M_c}.
\end{equation}
La collision SRC pondérée se fait autour du centre de masse massique $U_c$ :
\begin{equation}
  v_i' = U_c + R(\theta_c)(v_i-U_c),
\end{equation}
ce qui conserve exactement le moment pondéré de la cellule tant que la rotation est appliquée aux particules fluides de cette cellule.

La projection Q6 reste une projection de vitesse. Elle est appliquée au champ $U_c$ reconstruit par dépôt pondéré, puis la correction $\delta U$ est interpolée aux particules fluides. Une correction globale exacte de moment pondéré est appliquée lorsque nécessaire, afin d'éviter une dérive du moment total.

\subsection{Pool préalloué et rôles de particules}

Pour rester compatible avec un futur portage OpenMP, l'ajout de particules n'est pas traité par une croissance dynamique non bornée des tableaux. Le code utilise un pool préalloué :
\begin{equation}
  N_{\mathrm{cap}} = \alpha_{\mathrm{cap}}\,\gamma N_xN_y,
\end{equation}
avec typiquement $\alpha_{\mathrm{cap}}=2$. Les particules possèdent un rôle logique :
\begin{itemize}
  \item rôle 0 : particule inactive, disponible dans le pool libre ;
  \item rôle 1 : particule fluide réelle, participant au dépôt, à la collision, au thermostat, à Q6 et au remap ;
  \item rôle 2 : particule latente, stockée dans le domaine mais ne représentant pas encore du fluide.
\end{itemize}
La fonction de masque actif ne sélectionne que les particules de rôle fluide. Les particules inactives ou latentes sont exclues des dépôts, de la collision, du thermostat, de la projection et des diagnostics hydrodynamiques.

Cette distinction est importante pour les futurs solides mobiles, fronts libres ou domaines partiellement remplis : une particule peut exister dans les tableaux sans être un porteur de masse fluide.

\subsection{Extraction, insertion et remap local conservatif}

La population cible d'une cellule mouillée est définie par trois seuils :
\begin{equation}
  N_{\min},\qquad N_{\mathrm{target}}=\gamma,
  \qquad N_{\max}.
\end{equation}
Si $N_c<N_{\min}$, des particules sont insérées depuis le pool libre jusqu'à $N_{\mathrm{target}}$. Si $N_c>N_{\max}$, des particules sont extraites et replacées dans le pool libre. Le choix nominal des particules extraites est \texttt{closest\_to\_cell\_mean}, c'est-à-dire que l'on retire d'abord les particules dont la vitesse est la plus proche de la vitesse moyenne locale, afin de minimiser la perturbation du moment avant remap.

Le remap local impose ensuite, dans chaque cellule mouillée non vide,
\begin{equation}
  \sum_{i\in c} m_i^{\mathrm{new}} = M_c^{\mathrm{target}},
  \qquad
  \sum_{i\in c} m_i^{\mathrm{new}} v_i^{\mathrm{new}} = M_c^{\mathrm{target}} U_c^{\mathrm{target}}.
\end{equation}
Dans le mode simple \texttt{scale\_preserve\_velocity}, les masses d'une cellule sont d'abord multipliées par un facteur commun. Ce mode impose exactement $M_c^{\mathrm{target}}$ et conserve la vitesse moyenne pondérée lorsque la cible de moment est la vitesse avant remap. Un solveur \texttt{min\_change} existe également, mais le mode de mise à l'échelle s'est montré plus robuste pour les premiers tests.

L'ordre nominal après SRC/Q6 est le suivant :
\begin{enumerate}[label=\arabic*)]
  \item reconstruire le champ cible $U_c$ avant édition de population ;
  \item extraire les particules des cellules surpeuplées ;
  \item insérer les particules dans les cellules sous-peuplées ;
  \item remapper les masses pour imposer masse et moment cibles ;
  \item appliquer éventuellement un thermostat pondéré post-remap.
\end{enumerate}
Ainsi, l'extraction et l'insertion modifient le support particulaire, tandis que le remap restaure les grandeurs hydrodynamiques de cellule.

\subsection{Masque mouillé/non mouillé et domaine fluide évolutif}

La branche ajoute un masque de cellule
\begin{equation}
  \chi_c \in \{0,1\},
\end{equation}
où $\chi_c=1$ désigne une cellule mouillée, donc active pour le remap de masse, et $\chi_c=0$ une cellule non mouillée. La masse cible devient alors conceptuellement
\begin{equation}
  M_c^{\mathrm{target}} = \chi_c\,\rho_0 V_c.
\end{equation}
Dans la version actuelle, les cellules non mouillées ne sont pas forcées vers une masse cible et ne déclenchent ni insertion, ni extraction, ni remap de masse fluide. Le masque peut être fixe ou mis à jour par hystérésis de masse, ce qui prépare le traitement de domaines fluides évolutifs, d'interfaces libres ou de solides mobiles.

La couche latente complète ce formalisme. Un domaine peut être initialisé avec des particules latentes de masse et vitesse nulles ; une cellule source peut ensuite activer des particules fluides de masse non nulle. Le remap n'est appliqué qu'aux cellules mouillées, de sorte qu'un domaine vide n'est plus rempli artificiellement au premier pas.

Cette structure sera également utile pour les solides immergés : les cellules solides ou non mouillées auront $M_c^{\mathrm{target}}=0$, tandis que les cellules fluides ou partiellement fluides auront une cible ajustée à leur fraction fluide. Les particules virtuelles de paroi restent, elles, des contributions collisionnelles temporaires et non des particules fluides transportées.

\subsection{Sécurité de masse et renormalisation conservatrice}

Les premiers tests hétérogènes ont montré que le remap masse/moment peut satisfaire les contraintes algébriques en créant des particules très lourdes ou très légères lorsque l'état initial est extrêmement déséquilibré. Une sécurité conditionnelle a donc été ajoutée : si les masses candidates sortent d'une bande admissible, par exemple
\begin{equation}
  m_i \notin [0.25m_0,4m_0],
\end{equation}
la cellule bascule vers un mode \texttt{uniform\_mass\_velocity\_shift}. Les masses sont remises à
\begin{equation}
  m_i = \frac{M_c^{\mathrm{target}}}{N_c},
\end{equation}
puis les vitesses sont décalées uniformément pour restaurer le moment cible.

Une renormalisation plus douce a ensuite été conçue pour limiter la dérive long terme de la distribution des masses. Pour chaque cellule, on mémorise
\begin{equation}
  M_c,
  \qquad
  U_c,
  \qquad
  E_{\mathrm{th},c}=\frac12\sum_{i\in c}m_i\abs{v_i-U_c}^2.
\end{equation}
Les masses sont relaxées vers $M_c/N_c$, puis renormalisées pour conserver $M_c$. Les vitesses sont ensuite décalées pour restaurer $U_c$, et les fluctuations sont rescalées pour restaurer $E_{\mathrm{th},c}$. La correction préserve donc localement
\begin{equation}
  M_c,
  \qquad
  U_c,
  \qquad
  E_{\mathrm{th},c},
\end{equation}
tout en réduisant l'hétérogénéité des masses.

Deux validations récentes ont précisé le statut de cette correction. D'abord, un smoke Poiseuille court de 500 pas avec une renormalisation unique au pas 300 a montré une chute nette de la dispersion globale des masses, sans rupture des autres contraintes. À pas communs, la dispersion relative finale passe typiquement de
\begin{equation}
  \frac{\sigma_m}{\langle m\rangle}\simeq 0.079
  \quad\text{sans renormalisation}
  \qquad \text{à} \qquad
  \frac{\sigma_m}{\langle m\rangle}\simeq 0.059
  \quad\text{avec renormalisation},
\end{equation}
à 500 pas. Au pas de renormalisation, les résidus de moment et d'énergie thermique locale restent au niveau du roundoff, ce qui confirme que la correction ne se comporte pas comme une force volumique cachée.

Ensuite, le protocole long de référence a consisté à appliquer une seule renormalisation au pas 3000 d'un Poiseuille Q6+resampling wallVP de 5000 pas, avec tous les autres paramètres inchangés. Le résultat final est
\begin{equation}
  \frac{\sigma_m}{\langle m\rangle}=0.200,
  \qquad
  M_{\mathrm{relRMS}}\simeq 10^{-16},
  \qquad
  k_BT\simeq 0.00987,
\end{equation}
contre environ $0.234$ de dispersion des masses dans le run sans renormalisation. Le profil Poiseuille reste excellent, avec $R^2\simeq0.9978$ et $\nu_{\mathrm{eff}}\simeq18.0$, à comparer à $R^2\simeq0.9979$ et $\nu_{\mathrm{eff}}\simeq17.65$ sans renormalisation. La correction est donc validée comme sécurité numérique ponctuelle ou conditionnelle : elle réduit l'hétérogénéité des masses sans dégrader visiblement la masse cellulaire, la température ou le profil hydrodynamique.

\subsection{Tests fonctionnels de population}

Le premier test conçu pour faire tourner explicitement la boucle complète est le cas \texttt{paired\_pockets}. Il crée simultanément une poche vide et une poche surpeuplée en déplaçant les particules de l'une vers l'autre. Un résultat typique est :
\begin{itemize}
  \item au pas 0 : $N_{\min}=0$, $N_{\max}=40$, $M_{\mathrm{relRMS}}\simeq 0.221$ ;
  \item au pas 1 : 500 particules extraites, 500 insérées, $N\in[19,21]$, $N_{\mathrm{active}}$ constant ;
  \item après 300 pas : $N\in[17,23]$, $M_{\mathrm{relRMS}}\simeq 10^{-16}$ et les compteurs d'extraction/insertion restent fixés à 500.
\end{itemize}
Ce test démontre que le pool fonctionne comme un tampon de recyclage et non comme une source non bornée de particules.

Un second protocole impose une hétérogénéité distribuée de population. Les cas avec écart-type initial modéré sont ramenés vers la bande $N\in[14,26]$ avec $M_{\mathrm{relRMS}}$ au roundoff. Pour des hétérogénéités extrêmes, la sécurité de masse empêche les masses aberrantes observées avant son introduction. Dans un cas sévère représentatif, on observe après 300 pas
\begin{equation}
  N\in[14,26],
  \qquad
  M_{\mathrm{relRMS}}\simeq 10^{-16},
  \qquad
  \frac{\sigma_m}{\langle m\rangle}\simeq 0.13,
\end{equation}
avec des masses restant dans une plage contrôlée.

Le test d'injection/remplissage part d'un domaine entièrement latent, puis active progressivement des particules dans une cellule source. Dans un run accéléré, le nombre de cellules mouillées croît de 0 à 2048, la fraction mouillée atteint 1 vers le pas 105, et la masse des cellules mouillées reste contrôlée avec $M_{\mathrm{relRMS}}\simeq 10^{-14}$. Ce test valide la séparation entre particules latentes et particules fluides. Il montre aussi qu'un remplissage brutal est exigeant pour la distribution des masses : dans le run analysé, $\sigma_m/\langle m\rangle$ atteint environ 0.40 avec des masses proches de la borne de sécurité supérieure. Ce cas doit donc être lu comme un stress-test du mécanisme, non comme un réglage physique nominal.

\subsection{Validation hydrodynamique préliminaire}

Une fois la boucle de population fonctionnelle, la priorité est de vérifier que la méthode produit encore un comportement fluide crédible. Deux cas ont été retenus : Taylor--Green forcé en domaine périodique et Poiseuille forcé avec parois et particules virtuelles de paroi.

Sur Taylor--Green, le cas Q6 resamplé donne typiquement
\begin{equation}
  \mathrm{rmsDiv}\simeq 4\times 10^{-17},
  \qquad
  M_{\mathrm{relRMS}}\simeq 10^{-16},
  \qquad
  k_BT\simeq 0.0113,
\end{equation}
avec $N\in[14,26]$. La cohérence du mode Taylor--Green est meilleure que dans le témoin classic, avec une valeur représentative d'environ 0.893 contre 0.747, et l'enstrophie bruitée est réduite. Ce résultat suggère que Q6+resampling ne détruit pas les structures cohérentes de bulk.

Sur Poiseuille, les premiers essais sans particules virtuelles de paroi n'étaient pas exploitables : le profil était trop influencé par le glissement pariétal. Les particules virtuelles \texttt{wallVP} ont donc été réintroduites en version pondérée. Elles sont ajoutées seulement comme contributions agrégées à la collision près des murs :
\begin{equation}
  U_c^{\mathrm{coll}}=
  \frac{P_{\mathrm{real}}+P_{\mathrm{virt}}}{M_{\mathrm{real}}+M_{\mathrm{virt}}},
\end{equation}
et seules les particules fluides réelles sont mises à jour. Les particules virtuelles ne participent ni au transport, ni au pool, ni au remap.

Le run Poiseuille wallVP à 5000 pas fournit une comparaison directe :
\begin{center}
\begin{tabular}{lccc}
\toprule
Cas & $R^2$ parabolique & $\nu_{\mathrm{eff}}$ & $M_{\mathrm{relRMS}}$ \\
\midrule
SRC classic wallVP pur & $0.9959$ & $16.99$ & $0.15$ \\
Q6+resampling wallVP & $0.9979$ & $17.65$ & $\simeq10^{-16}$ \\
Q6+resampling wallVP + renorm. ponctuelle & $0.9978$ & $18.01$ & $\simeq10^{-16}$ \\
\bottomrule
\end{tabular}
\end{center}
Le résultat important n'est pas la valeur absolue de la viscosité, mais la proximité entre le témoin SRC wallVP pur et le cas Q6+resampling. La méthode resamplée conserve donc la physique de transport du cas wallVP de référence, avec un écart d'environ 4 \%, tout en supprimant les fluctuations de masse cellulaire. Le coût du run sans renormalisation est une dispersion des masses particulaires d'environ
\begin{equation}
  \frac{\sigma_m}{\langle m\rangle}\simeq 0.23.
\end{equation}
La renormalisation ponctuelle au pas 3000 abaisse cette valeur à environ $0.20$ à 5000 pas, tout en conservant un fit parabolique de qualité comparable et une viscosité effective très proche. Cette comparaison montre que la dispersion des masses peut être contrôlée sans disqualifier la dynamique Poiseuille.

\subsection{Interprétation et statut de la branche resampling}

La nouvelle méthode ne doit pas être interprétée comme une simple variante de Q9. Elle remplace le traitement des trous de population par une régularisation du support particulaire. Le principe devient :
\begin{equation}
\begin{aligned}
  &\text{support particulaire contrôlé}
  + \text{masse cellulaire imposée}
  + \text{projection Q6} \\
  &\hspace{3em}\Longrightarrow
  \text{fluide quasi incompressible pondéré}.
\end{aligned}
\end{equation}
Q9 peut rester utile comme diagnostic ou correction basse fréquence, mais il n'est plus nécessairement le mécanisme principal de contrôle de la densité lorsque le remap impose $M_c$ cellule par cellule.

Le statut actuel est donc le suivant :
\begin{itemize}
  \item la boucle extraction--insertion--remap est fonctionnelle ;
  \item les cellules non mouillées et les particules latentes permettent un domaine fluide évolutif ;
  \item les tests de population démontrent la réparation rapide des défauts de support sans force attractive ;
  \item Taylor--Green valide qualitativement le comportement bulk ;
  \item Poiseuille wallVP valide un premier comportement canal/paroi avec viscosité effective proche du SRC de référence ;
  \item la distribution des masses reste le principal point de vigilance, mais la renormalisation locale préservant $M,u,E$ est maintenant validée comme sécurité ponctuelle ou conditionnelle.
\end{itemize}
La branche resampling devient donc la voie la plus prometteuse pour traiter les cellules pauvres, les fronts de remplissage, les solides immergés et, à terme, les solides mobiles ou interfaces libres. Elle doit encore être validée statistiquement sur des runs plus longs et plusieurs graines, puis portée vers C++/OpenMP une fois la formulation stabilisée.




\section{Portage C++/OpenMP du resampling pondéré}

Les développements précédents ont stabilisé la formulation MATLAB du resampling pondéré. La dernière phase du projet a consisté à en traduire le cœur dans le dépôt \texttt{SRC\_incompressible\_openMP}, sans transformer le code C++ en générateur de cas physiques. La règle de workflow retenue est volontairement stricte : les états initiaux \texttt{.smpcd} sont préparés par des scripts MATLAB \texttt{prepare\_\dots.m} dans \texttt{init/}, les scripts \texttt{.sh} écrivent les fichiers \texttt{params*.kv} et lancent les exécutables OpenMP dans \texttt{runs/}, puis les scripts MATLAB \texttt{analyze\_\dots.m} post-traitent les sorties. Le noyau C++ ne fabrique donc pas les cas physiques : il lit un état initial, applique la dynamique SRC/Q6/resampling et écrit des diagnostics et dumps reproductibles.

Cette séparation est importante. Elle évite de multiplier les chaînes de génération, conserve MATLAB comme référence pour les scénarios de validation et permet de comparer directement les mêmes configurations entre prototype et code parallèle. Le C++ reste focalisé sur les opérations à fort coût : streaming, conditions limites, collision, projection elliptique, resampling, dépôt pondéré, diagnostics de masse et écriture d'états.

\subsection{Infrastructure OpenMP ajoutée}

Le portage a été introduit par jalons courts et testables. L'état particulaire OpenMP a été étendu avec :
\begin{itemize}
  \item une masse individuelle $m_i$ ;
  \item un rôle logique \texttt{Fluid}, \texttt{Inactive} ou \texttt{Latent} ;
  \item une distinction explicite entre \texttt{role} et \texttt{type}, afin de conserver la possibilité future d'espèces fluides multiples ;
  \item un format \texttt{.smpcd} V2 persistant les masses et les rôles ;
  \item un dépôt pondéré ne comptant que les particules \texttt{Fluid} réelles ;
  \item des diagnostics de masse, population, pool, remap, sécurité de masse et garde de population.
\end{itemize}

La séparation rôle/type est une différence structurante par rapport à certains prototypes MATLAB. Le rôle indique si une particule participe à la dynamique fluide courante. Le type est réservé à la nature physique future de la particule : espèce, matériau, traceur ou famille de paroi. Ainsi, une particule peut être présente dans les tableaux mais inactive, latente ou fluide. Les particules \texttt{Inactive} et \texttt{Latent} sont exclues du dépôt, de la collision, de Q6, du thermostat, du remap, du \texttt{mass guard} et des diagnostics hydrodynamiques.

Le dépôt pondéré utilisé par la collision et les diagnostics OpenMP est
\begin{equation}
  N_c = \sum_{i\in c,\;r_i=\mathrm{Fluid}} 1,\qquad
  M_c = \sum_{i\in c,\;r_i=\mathrm{Fluid}} m_i,
\end{equation}
\begin{equation}
  P_c = \sum_{i\in c,\;r_i=\mathrm{Fluid}} m_i v_i,
  \qquad
  U_c = \frac{P_c}{M_c}.
\end{equation}
Les particules virtuelles de paroi \texttt{wallVP} sont conservées dans leur rôle collisionnel : elles modifient le centre de collision près des murs, mais ne sont pas versées dans le pool, ne sont pas remappées et ne sont pas interprétées comme des particules fluides transportées.

\subsection{Correction d'une première divergence avec MATLAB}

Une étape importante du portage a été la mise en évidence d'une divergence conceptuelle. La première version OpenMP du resampling était surtout \emph{mass-driven} : les cellules pauvres ou riches étaient détectées à partir de seuils sur $M_c$. Cette version pouvait donc imposer
\begin{equation}
  M_c \simeq M_{\mathrm{target}}
\end{equation}
tout en laissant une cellule mouillée avec un nombre très faible de particules actives. Le champ de masse devenait alors parfait, mais le support particulaire $N_c$ pouvait rester insuffisant. Ce défaut est apparu dans le cas d'obstacle en canal ouvert : la masse était ramenée au roundoff, mais le domaine aval se dépeuplait localement et les masses individuelles compensaient trop fortement le manque de particules.

Ce comportement n'était pas conforme au prototype MATLAB. Dans la méthode de référence, le resampling est d'abord un contrôle de population :
\begin{equation}
  N_c < N_{\min} \quad\Longrightarrow\quad \text{repeupler la cellule},
\end{equation}
\begin{equation}
  N_c > N_{\max} \quad\Longrightarrow\quad \text{dépeupler la cellule},
\end{equation}
puis seulement ensuite une correction de masse/moment :
\begin{equation}
  M_c \rightarrow M_{\mathrm{target}},\qquad
  P_c \rightarrow M_{\mathrm{target}}U_c.
\end{equation}

Le patch OpenMP 0140 a donc réintroduit le garde de population $N_{\min}/N_{\mathrm{target}}/N_{\max}$ comme mécanisme principal. Pour $\gamma=20$, les scripts de validation utilisent typiquement
\begin{equation}
  N_{\min}=14,\qquad N_{\mathrm{target}}=20,\qquad N_{\max}=26.
\end{equation}
La logique devient conforme au développement MATLAB : on maintient d'abord un support particulaire suffisant, puis on utilise les masses pondérées pour fermer la contrainte de masse.

\subsection{Variante hybride efficace pour OpenMP}

Le mécanisme retenu est une variante hybride pensée pour éviter un appariement global donneur--receveur coûteux. Trois familles d'opérations sont distinguées.

Premièrement, si une cellule mouillée est sous-peuplée mais non vide, le code effectue un \emph{split local} de particules fluides lourdes. Une particule de masse $m$ et vitesse $v$ est remplacée par deux particules fluides de masse $m/2$ et vitesse $v$ :
\begin{equation}
  (m,v)\quad\longrightarrow\quad
  \left(\frac{m}{2},v\right)+
  \left(\frac{m}{2},v\right).
\end{equation}
Cette opération augmente $N_c$ tout en conservant localement masse, moment et énergie cinétique de translation. Elle utilise un slot \texttt{Inactive} du pool, mais ne crée pas de masse nouvelle.

Deuxièmement, si une cellule est surpeuplée, le code applique une extraction/merge locale qui renvoie des slots vers le pool inactif et évite que $N_c$ ne dérive trop au-dessus de $N_{\max}$. Cette opération reste locale à la cellule ; elle évite les tris globaux et les recherches donneur--receveur à coût quadratique.

Troisièmement, les mécanismes d'activation latente, d'extraction et d'insertion préexistants restent disponibles pour les cas de front mouillé, d'inlet, de cellules vides ou de pool déjà préparé. La logique générale est donc :
\begin{enumerate}[label=\arabic*)]
  \item dépôt pondéré des particules \texttt{Fluid} ;
  \item classification wet/dry et classification $N_c$ ;
  \item garde de population $N_{\min}/N_{\mathrm{target}}/N_{\max}$ ;
  \item activation latente, extraction et insertion si activées ;
  \item remap masse/moment tous les $K$ pas ;
  \item \texttt{mass guard} avec la même cadence que le remap de masse ;
  \item renormalisation thermique à chaque pas lorsque le resampling est actif.
\end{enumerate}

Ce choix est plus favorable à OpenMP que le transfert global donneur--receveur. Les décisions sont essentiellement locales à la cellule ; l'accès au pool peut être organisé par listes compactes et préallocations ; le coût principal reste linéaire en nombre de particules et de cellules. Les optimisations déjà intégrées évitent aussi les reconstructions coûteuses de plans de mutation à chaque dépôt et suppriment un diagnostic thermique quadratique qui provoquait un ralentissement massif sur Poiseuille.

\subsection{Fréquence des opérations dans la version OpenMP}

Le portage OpenMP distingue explicitement les fréquences des opérations. Le pas SRC classique et Q6, lorsqu'il est activé, restent exécutés à chaque pas. Le resampling ajoute les opérations suivantes :
\begin{center}
\small
\begin{tabular}{p{0.33\linewidth}p{0.36\linewidth}p{0.22\linewidth}}
\toprule
Opération & Rôle & Fréquence / paramètre \\
\midrule
\texttt{resamplingEnable} & verrou global du chemin mutateur resampling & lu à chaque pas \\
Dépôt pondéré & calcule $N_c,M_c,P_c,U_c$ sur \texttt{Fluid} & chaque pas \\
Classification wet/dry & détermine les cellules mouillées & chaque pas, \texttt{resamplingWetMaskMode} \\
Garde de population & impose la bande $N_{\min}/N_{\mathrm{target}}/N_{\max}$ & chaque pas, \texttt{resamplingPopulationGuardEnable} \\
Activation latente & \texttt{Latent}$\rightarrow$\texttt{Fluid} & chaque pas si activée \\
Extraction/insertion & basculement \texttt{Fluid}$\leftrightarrow$\texttt{Inactive} & chaque pas si activé \\
Remap de masse & impose $M_c\rightarrow M_{\mathrm{target}}$ & tous les $K$ pas, \texttt{resamplingMassRenormalizationPeriod} \\
\texttt{mass guard} & borne $m_i$ dans $[m_{\min},m_{\max}]$ & même cadence que le remap \\
Renormalisation thermique & restaure/cohère $U_c$ et $E_{\mathrm{th},c}$ & chaque pas si activée \\
\bottomrule
\end{tabular}
\end{center}
La convention pour $K$ est :
\begin{equation}
  K=1 : \text{remap à chaque pas},\qquad
  K>1 : \text{remap tous les }K\text{ pas},\qquad
  K=0 : \text{remap désactivé}.
\end{equation}
Pour les validateurs post-0140, le réglage nominal est
\begin{equation}
  K=10,\qquad m_{\min}=0.5,\qquad m_{\max}=2.0.
\end{equation}

\subsection{Différences pratiques entre MATLAB et OpenMP}

La version OpenMP ne doit pas être vue comme une simple copie ligne à ligne du prototype MATLAB. Elle en conserve la logique physique, mais avec des choix adaptés au coût parallèle :
\begin{itemize}
  \item le C++ lit des états \texttt{.smpcd} générés par MATLAB ; il ne génère pas les cas physiques ;
  \item les rôles sont persistés dans le fichier d'état, ce qui rend possible un redémarrage exact avec pool, particules latentes et inactives ;
  \item les opérations locales de garde de population privilégient split/merge cellulaire plutôt qu'un appariement global donneur--receveur ;
  \item la renormalisation de masse est cadencée par un paramètre, alors que la garde de population reste à chaque pas ;
  \item les diagnostics runtime sont plus nombreux que dans MATLAB : masses, population wet, pool, guard, solid leak, divergence Q6, opérations d'extraction/insertion et coût mural ;
  \item les restrictions de sécurité du noyau C++ sont plus fortes : par exemple, le cas \texttt{Q6 + inlet/outlet + solide immergé} n'autorise encore que certains solides rectangulaires masqués, afin d'éviter un désaccord entre le masque solide de projection et la dynamique particulaire.
\end{itemize}

Cette dernière restriction explique pourquoi le vrai cylindre en canal ouvert n'a pas encore été activé avec Q6/resampling. Les cylindres périodiques et les obstacles rectangulaires ouverts sont validables avec le noyau actuel ; le cylindre inlet/outlet demandera un futur chantier C++ spécifique sur les cellules coupées et les faces solide--fluide dans l'opérateur elliptique.

\subsection{Validations OpenMP actuellement obtenues}

Les validations post-portage montrent que les briques essentielles fonctionnent dans plusieurs régimes. Le tableau suivant résume le statut, en mettant l'accent sur le gain par rapport au témoin \texttt{classic}.

\begin{center}
\small
\begin{tabular}{p{0.23\linewidth}p{0.32\linewidth}p{0.34\linewidth}}
\toprule
Cas & Observation principale & Gain par rapport à \texttt{classic} \\
\midrule
Taylor--Green périodique et forcé & Q6 et Q6+resampling conservent le motif cohérent ; le cas forcé maintient une amplitude stationnaire lorsque le forçage est calibré. & Q6 réduit la divergence de grille ; Q6+resampling impose la masse cellulaire au roundoff et maintient une population bornée. \\
Population aléatoire distribuée & Le cas déclenche un désordre initial de population comparable au prototype MATLAB. & Le resampling ramène le support vers la bande $[14,26]$ et supprime l'erreur de masse, alors que \texttt{classic}/Q6 gardent des fluctuations massiques. \\
Poiseuille \texttt{wallVP} & Les trois cas donnent un profil parabolique cohérent ; Q6+resampling ne détruit pas le transport. & La masse cellule est ramenée au roundoff et la population reste bornée. Après optimisation 0132--0133, le surcoût court observé était d'environ $6.4\,\mathrm{s}$ contre $4.4\,\mathrm{s}$ pour Q6 seul sur 1000 pas, soit environ $1.45\times$ au-dessus de Q6. \\
Cylindre périodique immergé & Le solide circulaire périodique reste propre et le sillage reste stable. & Le \texttt{solidLeakMass} représentatif passe d'environ 76 en \texttt{classic} à 22 en Q6 et 18 en Q6+resampling, avec masse cellulaire contrôlée. \\
Backward step & Le cas inlet/outlet avec solide rectangulaire reste crédible ; la recirculation n'est pas encore très marquée, mais l'écoulement ne présente pas de pathologie. & Q6+resampling stabilise la masse et le support sans casser les conditions ouvertes ni le masque solide. \\
Injection/fill depuis domaine inactif & Le domaine part vide/inactif puis se remplit progressivement par l'inlet localisé. & \texttt{classic} et Q6 seuls accumulent localement de très fortes masses ; Q6+resampling propage le front mouillé et maintient la masse des cellules mouillées au niveau cible. \\
Obstacle rectangulaire en canal ouvert & Le cas a révélé le défaut du premier portage : masse correcte mais support particulaire appauvri. & Après 0140, ce cas devient le test décisif du garde $N_{\min}/N_{\mathrm{target}}/N_{\max}$ et doit remplacer les conclusions de la variante mass-driven. \\
\bottomrule
\end{tabular}
\end{center}

Les gains ne sont donc pas seulement quantitatifs sur un indicateur. Ils sont surtout structurels : le cas resamplé maintient simultanément un champ de masse quasi incompressible, un support particulaire suffisant et la possibilité de domaines mouillés évolutifs. C'est ce point qui rend la méthode SRC/MPCD pondérée beaucoup plus puissante que les corrections eulériennes de cellules pauvres testées auparavant.

\subsection{Points de vigilance restants}

Le portage OpenMP est désormais plus conforme à la méthode MATLAB, mais plusieurs points restent ouverts :
\begin{itemize}
  \item la campagne post-0140 doit confirmer que le garde de population supprime le dépeuplement observé dans les canaux ouverts ;
  \item la fréquence optimale du remap de masse $K$ doit être balayée selon les cas ;
  \item les bornes $m_{\min},m_{\max}$ doivent être interprétées comme des sécurités numériques, non comme des paramètres physiques ;
  \item le vrai cylindre inlet/outlet avec Q6 exigera un masque circulaire compatible projection, faces coupées et diagnostics de fuite de flux projeté ;
  \item le coût de Q6 elliptique dans les géométries ouvertes reste un chantier d'optimisation séparé ; l'essai de warm-start a été abandonné pour ne pas polluer les validations physiques.
\end{itemize}

Le prochain objectif scientifique est donc de terminer la revalidation post-0140 sur les cas Taylor--Green, Poiseuille, remplissage, obstacle ouvert et marche, avant de chercher une allée tourbillonnaire plus nette. Le prochain objectif algorithmique, séparé, sera l'optimisation du solveur Q6 et la généralisation propre du masque solide courbe en inlet/outlet.




\section{Développements OpenMP post-0143 : segments ouverts, configuration explicite et capacité fermée}

Les développements postérieurs au portage OpenMP du resampling pondéré ont poursuivi deux objectifs complémentaires. Le premier est logiciel : réduire les états de configuration ambigus et permettre des conditions limites ouvertes plus générales, notamment plusieurs segments \texttt{inlet}/\texttt{outlet} sur une même face. Le second est physique : traiter le cas limite d'un domaine fermé déjà plein soumis à un flux entrant net, situation incompatible avec une fermeture incompressible stricte mais naturellement interprétable comme une compression faible, puis forte, d'un liquide dans une enceinte rigide.

Cette section documente les patchs 0143 à 0154. Elle ne remplace pas les sections précédentes : les cas Taylor--Green, Poiseuille, solides immergés et remplissage restent les validations de référence du couple Q6+resampling. Les nouveaux développements ajoutent un chemin d'étude pour les domaines fermés surpressurisés et pour la future interaction avec une paroi déformable.

\subsection{Conditions limites ouvertes segmentées en coordonnées relatives}

La première simplification concerne les conditions limites ouvertes. L'ancien formalisme associait un mode à une face entière, par exemple \texttt{bcLeft=inlet} ou \texttt{bcRight=outlet}, avec une option d'ouverture partielle limitée à un unique intervalle. Ce formalisme ne permettait pas d'écrire proprement un cas du type : entrée sur la partie haute de la face gauche, sortie sur la partie basse de la même face, et paroi solide ailleurs.

Le nouveau formalisme introduit une liste de segments ouverts :
\begin{verbatim}
openBoundarySegmentsEnable = true
openBoundarySegmentCount = N
openBoundarySegmentK = face mode sMin sMax ux uy type mass
\end{verbatim}
où \texttt{face}\(\in\{\texttt{left},\texttt{right},\texttt{bottom},\texttt{top}\}\), \texttt{mode}\(\in\{\texttt{inlet},\texttt{outlet}\}\), et \(s\in[0,1]\) est la coordonnée tangentielle relative à la face. Sur les faces gauche/droite, \(s\) désigne la coordonnée verticale relative ; sur les faces bas/haut, il désigne la coordonnée horizontale relative. Les champs \texttt{ux}, \texttt{uy}, \texttt{type} et \texttt{mass} servent à l'injection ; ils sont présents dans la syntaxe des \texttt{outlet} pour garder un parseur à longueur fixe.

La convention retenue est volontairement stricte : les segments ouverts sont posés sur une face solide par défaut, et les portions non couvertes restent solides. Par exemple,
\begin{verbatim}
bcLeft   = solid
bcRight  = solid
bcBottom = solid
bcTop    = solid

openBoundarySegmentsEnable = true
openBoundarySegmentCount = 2
openBoundarySegment0 = left inlet  0.65 0.90  0.10 0.00  0 1.0
openBoundarySegment1 = left outlet 0.10 0.35  0.00 0.00  0 1.0
\end{verbatim}
représente une face gauche solide percée localement d'un inlet haut et d'un outlet bas. Les anciens paramètres d'ouverture partielle de type \texttt{leftOpenYMin}/\texttt{leftOpenYMax} ont été supprimés afin d'éviter que de vieux scripts produisent des comportements implicites non maîtrisés.

Du point de vue algorithmique, cette évolution a trois conséquences. Premièrement, le traitement lagrangien des frontières est effectué par face et par segment : une particule sortant par une portion solide est réfléchie, tandis qu'une particule sortant par un segment \texttt{outlet} est convertie en particule inactive. Deuxièmement, le réservoir \texttt{inlet} parcourt les segments d'injection et non plus une face unique. Troisièmement, Q6 reconstruit des profils de flux ouverts par face à partir des segments ; le solveur elliptique n'a pas besoin d'être modifié, car les profils discrets de flux étaient déjà stockés par rangée ou colonne de face.

Cette abstraction prépare aussi l'injection multi-type. Le champ \texttt{type} peut d'abord être interprété comme un traceur passif de provenance, par exemple pour distinguer un jet central et une gaine coaxiale. Une vraie physique multi-espèces exigerait toutefois des collisions, thermostats, dépôts massiques et remaps dépendants du type ; cette extension reste séparée.

\subsection{Modernisation de la configuration : suppression de \texttt{method}}

Le second nettoyage supprime l'ancien aiguillage \texttt{method=classic/q6}. Dans les versions précédentes, \texttt{method=q6} activait implicitement la projection, tandis que \texttt{projectionEnable=true} l'activait aussi. Les deux options étaient donc redondantes, et des configurations comme
\begin{verbatim}
method = q6
projectionEnable = false
\end{verbatim}
étaient sémantiquement ambiguës. Or Q6 n'est pas un modèle séparé de la projection : Q6 est précisément l'opérateur de projection de vitesse.

La configuration moderne repose donc sur deux interrupteurs principaux :
\begin{equation}
  \texttt{projectionEnable}\in\{\texttt{false},\texttt{true}\},
  \qquad
  \texttt{resamplingEnable}\in\{\texttt{false},\texttt{true}\}.
\end{equation}
Ils définissent quatre cas d'ablation :
\begin{center}
\begin{tabular}{lll}
\toprule
\texttt{projectionEnable} & \texttt{resamplingEnable} & Interprétation \\
\midrule
\texttt{false} & \texttt{false} & SRC/MPCD classique \\
\texttt{false} & \texttt{true}  & SRC classique avec garde numérique de support \\
\texttt{true}  & \texttt{false} & SRC avec projection Q6, sans resampling \\
\texttt{true}  & \texttt{true}  & Q6+resampling pondéré \\
\bottomrule
\end{tabular}
\end{center}

Le paramètre \texttt{resamplingPopulationGuardEnable} a également été supprimé. Lorsque \texttt{resamplingEnable=true}, le garde de population est considéré comme le coeur du resampling, et son comportement est piloté par les bornes
\begin{equation}
  N_{
m min} < N_{
m target} < N_{
m max}.
\end{equation}
Cette simplification est importante pour l'interprétation physique : le resampling peut être utilisé soit avec Q6, comme élément de la chaîne quasi-incompressible, soit sans Q6, comme outil purement numérique de maintien du support particulaire SRC classique.

\subsection{Masque mouillé pour domaines pleins et correction du dépeuplement pariétal}

Les essais en cuve pleine avec inlet/outlet localisés ont mis en évidence un défaut du mode \texttt{resamplingWetMaskMode=occupied}. Dans ce mode, une cellule est considérée mouillée uniquement si elle contient déjà de la masse. Près d'une paroi solide, une fluctuation peut vider temporairement une cellule ; celle-ci devient alors \emph{dry}, sort du champ du garde de population, et peut rester durablement vide. Le problème ne vient pas d'une fuite solide directe, mais d'un effet de cliquet wet/dry.

Pour une cuve initialement pleine, le bon choix est \texttt{resamplingWetMaskMode=active\_domain}. La cellule reste mouillée tant qu'elle appartient au domaine fluide actif, même si son occupation instantanée tombe à zéro. Le garde de population et l'insertion locale peuvent alors la réparer. Le passage à \texttt{active\_domain} a supprimé les cellules vides persistantes près des parois dans les runs de cuve pleine.

La règle pratique devient donc :
\begin{itemize}
  \item pour un remplissage depuis domaine vide ou latent, \texttt{occupied} reste adapté, car les cellules non atteintes par le front doivent rester sèches ;
  \item pour un domaine plein ou une cuve déjà mouillée, \texttt{active\_domain} est nécessaire pour empêcher le dépeuplement pariétal irréversible.
\end{itemize}

\section{Réponse de capacité fermée et viriel faiblement compressible}

\subsection{Motivation physique}

Un domaine fermé, plein, à parois infiniment rigides et soumis à un inlet sans outlet n'est pas compatible avec une incompressibilité stricte. Pour un fluide incompressible dans un volume fixe,
\begin{equation}
  \nabla\cdot u = 0,
  \qquad
  \int_{\partial\Omega} u\cdot n\,\mathrm{d}S = 0.
\end{equation}
Un flux entrant net viole cette compatibilité. Un liquide réel répondrait par sa compressibilité finie : la masse volumique augmenterait, la pression monterait, et la température pourrait monter si la compression est suffisamment adiabatique. La limite incompressible stricte correspondrait à une pression infinie pour imposer une quantité finie de liquide supplémentaire dans un volume rigide.

L'objectif n'est donc pas de stabiliser artificiellement ce cas, mais de permettre une transition continue : tant que le domaine reste sous sa capacité nominale, le comportement quasi-incompressible est conservé ; si un flux entrant maintenu produit une surmasse, l'intensité effective de Q6 diminue, la fermeture virielle se raidit, et le remap de masse cesse d'effacer la compression globale.

\subsection{Mesure de surcapacité}

On définit une masse de référence fermée
\begin{equation}
  M_{\rm ref} = N_{\rm ref}\,M_{\rm cell,ref},
\end{equation}
où \(N_{\rm ref}\) est le nombre de cellules du domaine fluide de référence et \(M_{\rm cell,ref}\) la masse cellulaire nominale, typiquement \(\gamma m_0\). La surcapacité relative est
\begin{equation}
  \eta_M = \max\left(0,\frac{M_{\rm tot}-M_{\rm ref}}{M_{\rm ref}}\right).
\end{equation}
Toutes les réponses de capacité sont continues en \(\eta_M\). Si \(\eta_M=0\), les cas validés précédemment restent strictement inchangés.

\subsection{Modulation de Q6 par \texttt{q6ProjectionStrength}}

La projection Q6 n'est pas remplacée par un nouveau solveur. Le coefficient d'application de la projection est simplement modulé :
\begin{equation}
  \alpha_{Q6}^{\rm eff}
  = \alpha_{Q6}^{0}\,
  \exp\left[-\left(\frac{\eta_M}{\eta_{Q6}}\right)^{p_{Q6}}\right],
\end{equation}
où \(\alpha_{Q6}^{0}\) est \texttt{q6ProjectionStrength}. Lorsque la cuve n'est pas sur-remplie, \(\alpha_{Q6}^{\rm eff}=\alpha_{Q6}^{0}\). Lorsque la surmasse augmente, Q6 s'affaiblit progressivement. Cette écriture évite un basculement brutal entre un solveur incompressible et un solveur compressible.

Dans les diagnostics OpenMP, on suit séparément la force nominale de Q6, le facteur de capacité, et la force effective réellement appliquée. Les premiers tests montrent que, dans un cas inlet seul maintenu, \(\alpha_{Q6}^{\rm eff}\) chute rapidement vers zéro lorsque la surcapacité atteint quelques pourcents : le système quitte alors volontairement le régime quasi-incompressible.

\subsection{Renforcement viriel continu}

La fermeture virielle est renforcée selon une loi de capacité, par exemple
\begin{equation}
  K_{\rm vir}^{\rm eff}
  = K_{\rm vir}^{0}\left[1+G_{\rm vir}
  \left(\frac{\eta_M}{\eta_{\rm vir}}\right)^{p_{\rm vir}}\right].
\end{equation}
La pression virielle cellulaire est reconstruite relativement à la masse cellulaire nominale :
\begin{equation}
  P_{{\rm vir},c}
  = K_{\rm vir}^{\rm eff}\left(\frac{M_c}{M_{\rm cell,ref}}-1\right).
\end{equation}
La pression totale utilisée pour les diagnostics bulk et pariétaux est
\begin{equation}
  P_{{\rm tot},c}=P_{{\rm kin},c}+P_{{\rm vir},c}.
\end{equation}

Un kick viriel optionnel applique une correction de vitesse proportionnelle à \(-\nabla P_{\rm vir}/\rho\), avec correction globale de moment si elle est activée. Dans un état uniformément pressurisé, ce gradient est faible ; la pression peut donc monter sans générer de kick volumique important. En revanche, un inlet localisé crée des gradients viriels et peut produire des accélérations violentes, voire une rupture numérique, si le flux est maintenu.

\subsection{Cible effective du remap de masse}

Le premier prototype de réponse fermée montrait un verrou : le remap de masse ramenait chaque cellule vers \(M_{\rm cell,ref}\), donc il effaçait presque toute surmasse avant que le viriel ne puisse vraiment monter. Le correctif consiste à conserver la surmasse globale dans la cible de remap :
\begin{equation}
  M_{\rm cell,target}^{\rm eff}
  = M_{\rm cell,ref} + \frac{\max(0,M_{\rm tot}-M_{\rm ref})}{N_{\rm ref}}.
\end{equation}
Ainsi, le resampling conserve son rôle numérique : homogénéiser la masse, le moment et le support particulaire. Mais il ne détruit plus la compression globale imposée par un flux entrant net. La masse totale peut alors augmenter, \(\eta_M\) croît, Q6 s'affaiblit et \(K_{\rm vir}^{\rm eff}\) augmente.

Cette modification préserve les validations antérieures car elle n'est active que lorsque la réponse de capacité est explicitement activée et que \(M_{\rm tot}>M_{\rm ref}\). Sinon, le remap conserve sa cible fixe historique.

\subsection{Injection additive et diagnostics de moment}

Pour qu'une cuve pleine puisse réellement se pressuriser, l'inlet ne doit pas seulement remplacer localement les particules du réservoir. Une option d'injection additive associe au flux normal imposé un budget de masse positif. Si le domaine est fermé, cette masse reste dans le domaine et augmente \(M_{\rm tot}\). Les diagnostics suivent alors : masse totale, surcapacité, force Q6 effective, module viriel effectif, pression virielle moyenne, kick viriel et activité du resampling.

Les runs de cuve fermée avec inlet seul montrent une montée de masse jusqu'à plusieurs pourcents de surcapacité, une extinction rapide de Q6, une croissance forte de \(K_{\rm vir}^{\rm eff}\) et une montée de l'énergie. Ils montrent aussi un budget de quantité de mouvement encore incomplet : les contributions Q6 et resampling ne suffisent pas à expliquer les sauts tardifs du moment global. Les contributions de paroi, d'inlet et de viriel doivent donc être instrumentées plus finement avant d'interpréter complètement la dynamique. Cette limite ne remet pas en cause la validation statique de la charge pariétale, mais elle encadre l'usage dynamique du modèle.

\section{Pression pariétale virielle et objectif parois déformables}

\subsection{Pourquoi la pression bulk ne suffit pas}

Dans un continuum, une pression uniforme exerce une traction normale sur une paroi, même si \(\nabla p=0\) dans le bulk. Dans le code particulaire, une fermeture virielle utilisée seulement comme kick volumique \(-\nabla P_{\rm vir}\) ne transmet pas automatiquement une pression uniforme à une paroi solide. Pour préparer un couplage avec une paroi déformable ou un modèle de rupture, il faut donc construire explicitement une charge pariétale.

Le diagnostic ajouté convertit la pression cellulaire reconstruite,
\begin{equation}
  P_{{\rm tot},c}=P_{{\rm kin},c}+P_{{\rm vir},c},
\end{equation}
en charge normale sur les faces solides. Pour une face solide adjacente à une cellule \(c\), la contribution élémentaire est
\begin{equation}
  \Delta \bm{F}_{\rm wall} = P_{{\rm tot},c}\,A_f\,\bm{n}_{\rm wall},
\end{equation}
où \(A_f\) vaut \(\Delta y\) pour les faces verticales et \(\Delta x\) pour les faces horizontales en 2D de profondeur unité. Les segments \texttt{inlet}/\texttt{outlet} sont exclus de cette intégration : seule la portion solide de la paroi porte la charge.

Les diagnostics OpenMP séparent la pression cinétique, la pression virielle, la pression totale par face, les longueurs solides par face, et les forces nettes \(F_x,F_y\). Cette séparation est essentielle pour distinguer une pression hydrostatique uniforme d'un chargement fortement asymétrique induit par un inlet localisé.

\subsection{Validation statique : boîte uniformément surcompressée}

Le test de validation 0152 initialise une boîte fermée, sans inlet, avec une masse uniforme
\begin{equation}
  M = \lambda M_{\rm ref},
  \qquad \lambda\in\{1.00,1.02,1.05,1.10\}.
\end{equation}
La dynamique est presque statique et virielle pure : pas de thermostat de paroi, pas de \texttt{wallVP} thermique et pas d'inlet. Le résultat attendu est
\begin{equation}
  p_{\rm wall,left}\simeq p_{\rm wall,right}
  \simeq p_{\rm wall,bottom}\simeq p_{\rm wall,top},
  \qquad
  \bm{F}_{\rm wall}^{\rm net}\simeq 0.
\end{equation}

Les résultats obtenus valident ce chaînage. La pression moyenne pariétale suit la loi attendue \(K_{\rm vir}^{\rm eff}\eta_M\), la pression par face est confondue pour les quatre parois, et la force nette relative est au niveau du bruit numérique. Ce test démontre que la pression bulk virielle est bien convertible en charge normale pariétale dans le cas de référence uniformément surcompressé.

\subsection{Cas dynamique inlet seul et limite actuelle}

Dans un cas inlet seul, la pression pariétale augmente fortement et devient non uniforme : la face proche de l'injection et la face supérieure peuvent recevoir une charge plus forte que les autres. Cette non-uniformité n'est pas forcément un défaut ; elle est attendue dans un régime transitoire de cuve fermée alimentée localement. En revanche, le budget de quantité de mouvement doit encore être complété. Pour un futur couplage solide, il faudra exposer dans les fichiers \texttt{summary\_runtime.csv} les impulsions associées aux parois, à l'inlet, au kick viriel brut, au kick viriel corrigé, au thermostat et aux corrections de projection.

Le statut actuel est donc le suivant : la conversion \(P_{\rm kin}+P_{\rm vir}\) vers une charge normale pariétale est validée en statique uniforme ; l'usage dynamique dans un domaine fermé alimenté par un inlet est prometteur, mais demande encore un budget d'impulsion complet.

\section{Positionnement bibliographique : MPCD quasi-incompressible, projection et resampling}

\subsection{Contexte MPCD/SRD}

La méthode MPCD/SRD a été introduite comme modèle mésoscopique de solvant par Malevanets et Kapral \cite{MalevanetsKapral1999}. La revue de Gompper, Ihle, Kroll et Winkler \cite{GompperIhleKrollWinkler2009} synthétise les variantes principales, l'interprétation hydrodynamique, le rôle du streaming/collision, les coefficients de transport et les extensions vers fluides complexes. Dans ce cadre, le fluide standard est un solvant fluctuant, localement conservatif en masse et moment, mais il n'est pas strictement incompressible. En pratique, on se place souvent à bas Mach ou à densité particulaire suffisante pour réduire les effets compressibles.

Les travaux sur MPCD non idéal traitent explicitement cette limite. Ihle, Tüzel et Kroll ont proposé un algorithme particulaire avec équation d'état non idéale fondé sur des collisions multiparticulaires biaisées dépendant de la densité locale \cite{IhleTuzelKroll2006}. Plus récemment, Zantop et Stark ont introduit une règle de collision MPCD modifiée produisant une équation d'état non idéale et une compressibilité fortement diminuée \cite{ZantopStark2021a,ZantopStark2021b}. Ces travaux sont proches dans l'objectif général -- réduire la compressibilité MPCD -- mais ils agissent principalement sur la collision et l'équation d'état.

\subsection{Projection et méthodes particulaires incompressibles}

Le choix Q6 s'inscrit dans la tradition des méthodes de projection de pression pour les équations incompressibles, initiée par Chorin \cite{Chorin1968}. Dans les méthodes particulaires, la même tension entre incompressibilité et compressibilité artificielle se retrouve en SPH. PCISPH impose l'incompressibilité par une procédure prédictive-corrective de pression \cite{SolenthalerPajarola2009}, tandis que les travaux récents sur l'ACSPH et le \(\delta\)-SPH mettent en évidence les liens entre correction pression/vitesse et compressibilité artificielle \cite{DeCourcy2024}. Des solveurs Navier--Stokes faiblement compressibles à projection de pression évolutive explorent également l'idée de ne pas imposer strictement \(\nabla\cdot u=0\), mais d'utiliser une projection pour amortir les ondes acoustiques et contrôler la pression \cite{Yang2021}.

La contribution Q6+resampling est différente. La projection Q6 traite la partie cinématique : elle retire la composante compressive du champ de vitesse lorsque le bilan global de flux est compatible avec l'incompressibilité. Le resampling pondéré traite une difficulté distincte : la qualité du support particulaire. Une cellule pauvre ou vide près d'une paroi ne doit pas être interprétée comme une dépression physique ; elle doit d'abord être réparée comme défaut de représentation. Le garde \(N_{
m min}/N_{
m target}/N_{
m max}\), les rôles \texttt{Fluid}/\texttt{Latent}/\texttt{Inactive} et le remap conservatif masse--moment--énergie séparent donc clairement
\begin{equation}
  \text{incompressibilité cinématique}
  \quad \text{et} \quad
  \text{support particulaire local}.
\end{equation}

Cet avantage est particulièrement net par rapport à deux stratégies concurrentes : augmenter fortement \(\gamma\), ce qui réduit le bruit mais augmente le coût, ou modifier la collision pour obtenir une équation d'état non idéale. L'approche développée ici conserve le coeur SRC/MPCD classique, permet des ablations propres \texttt{classic}/\texttt{q6}/\texttt{classic\_resampling}/\texttt{q6\_resampling}, et laisse la fermeture de pression comme un module séparé. Cette modularité est importante pour la parallélisation OpenMP et, à terme, GPU.

\section{Positionnement bibliographique : viriel faiblement compressible et pression pariétale}

\subsection{Contraintes, tenseurs de stress et forces pariétales en MPCD}

Les contraintes et tenseurs de stress MPCD ont été analysés par Winkler et Huang \cite{WinklerHuang2009}. Ils distinguent notamment des formulations internes et externes du tenseur de contrainte, et considèrent des fluides confinés. Imperio, Padding et Briels ont proposé une méthode de calcul des forces sur parois et particules incluses en MPCD avec conditions de non-glissement, en évitant de reconstruire explicitement le profil de vitesse près de la surface \cite{ImperioPaddingBriels2011}. Ces travaux constituent les références les plus proches pour la question des forces mécaniques exercées sur un solide dans MPCD.

La nouveauté de la fermeture actuelle ne réside pas dans l'existence d'une force de paroi MPCD, mais dans le fait que la pression bulk reconstruite,
\begin{equation}
  P_{\rm tot}=P_{\rm kin}+P_{\rm vir},
\end{equation}
est explicitement projetée en charge normale sur les portions solides des parois dans un domaine fermé surcompressé. Le test statique 0152 montre que cette charge donne une pression identique sur les quatre faces et une force nette nulle lorsque la surpression est uniforme. Cette étape est indispensable pour envisager un couplage avec un solide déformable : une pression uniforme de bulk doit charger une paroi même lorsque \(\nabla P=0\) dans le domaine.

\subsection{Parenté avec SPH faible compressibilité et FSI}

Les méthodes SPH faiblement compressibles utilisent depuis longtemps une équation d'état raide pour convertir une variation de densité en pression. Les modèles de paroi WCSPH et la qualité des pressions pariétales ont été étudiés par Valizadeh et Monaghan \cite{ValizadehMonaghan2015}. Côté fluide--structure et rupture, des cadres SPH récents couplent WCSPH à des modèles de solides déformables ou fracturables, par exemple le cadre d'Islam pour les grandes déformations et la fracture \cite{Islam2024}, ou des modèles 3D combinant WCSPH et pseudo-ressorts pour la rupture fragile sous chargement hydrodynamique \cite{Singh2025}. Ces références montrent que la chaîne \emph{pression fluide -- charge pariétale -- déformation/rupture} est une voie active dans les méthodes particulaires.

La différence est que la méthode présente n'est pas une formulation SPH de pression. Elle conserve le coeur SRC/MPCD, ajoute une projection de vitesse modulable, maintient le support particulaire par resampling pondéré, et introduit une pression virielle qui devient dominante seulement lorsque la capacité fermée est dépassée. À notre connaissance, cette combinaison -- projection MPCD affaiblie par surcapacité, remap conservant la surmasse globale, renforcement viriel continu et projection pariétale de \(P_{\rm kin}+P_{\rm vir}\) -- n'a pas été rapportée sous cette forme dans la littérature MPCD.

La revendication doit rester prudente : les ingrédients individuels existent dans des familles voisines, notamment MPCD non idéal, projection incompressible, artificial compressibility, stress MPCD et SPH-FSI. L'originalité se situe dans leur assemblage modulaire pour passer continûment d'un SRC/MPCD quasi-incompressible validé à un régime faiblement compressible capable de charger des parois fermées.



\section{Optimisation OpenMP, profilage contrôlé et branche de production allégée}

Les développements précédents ont établi la formulation physique et numérique de la chaîne SRC/Q6/resampling/viriel. Une étape séparée a ensuite été menée sur le dépôt C++/OpenMP afin de réduire le temps d'exécution sans modifier les fonctionnalités ni la trajectoire numérique. Cette étape a suivi une règle stricte : toute optimisation modifiant un chemin algorithmique devait être validée par une comparaison discriminante entre la version de référence et la version optimisée. Les cas utilisés pour cette comparaison étaient les quatre cas mono-configuration déjà employés dans le dépôt : Taylor--Green périodique, Poiseuille avec parois, obstacle rectangulaire ouvert, et piston avec fermeture virielle. Le critère d'acceptation était l'égalité, à tolérance fixée, des métriques physiques, numériques et de resampling sur l'ensemble des cas.

\subsection{Méthode de travail et séparation des branches}

La phase d'optimisation a conduit à distinguer trois usages du code OpenMP. La branche \texttt{SRC\_openMP\_optimized} est conservée comme branche de développement et de préparation GPU : elle garde les diagnostics internes fins, les scripts de profilage et les outils de comparaison. La branche \texttt{clean/openmp-light} dérive de cette version optimisée mais désactive les diagnostics internes par défaut afin de fournir une version OpenMP portable, plus proche d'un code de production. Enfin, une future branche \texttt{clean/openmp-production-stripped} pourra retirer complètement le code de profilage interne lorsque la version allégée aura été suffisamment stabilisée.

Cette séparation évite de confondre deux objectifs. D'une part, la branche instrumentée reste nécessaire pour identifier les noyaux dominants, préparer un futur backend GPU et comprendre les régressions. D'autre part, la branche light vise un exécutable qui ne crée pas de fichiers de profilage internes en usage normal et qui laisse les analyses de cas au post-traitement externe, conformément à l'architecture initiale du code.

\subsection{Instrumentation initiale et identification des points chauds}

Une première série de patches a ajouté des profils internes contrôlés : temps par phase du pas SRC/MPCD, profil Q6/CG, profil des dépôts particules--cellules, profil des gardes de population et de masse, et diagnostics post-guard. Ces profils ont été utiles pour localiser précisément les verrous. Ils ont montré que le coût n'était pas concentré dans un seul module, mais dans plusieurs chemins distincts : projection elliptique Q6, dépôts pondérés répétés, garde de population, garde de masse, et rafraîchissement thermique. Ils ont aussi permis d'écarter certaines optimisations intuitives mais non pertinentes, par exemple la seule réutilisation des buffers temporaires du \texttt{mass\_guard}, dont le gain s'est révélé marginal.

Le profilage a également révélé plusieurs coûts anormaux. Le plus net concernait la réponse de capacité fermée : lorsque \texttt{closedCapacityResponseEnable=false}, le code effectuait encore un dépôt de masse cellulaire avant de constater que la fermeture était désactivée. Un retour anticipé au début du chemin viriel a supprimé ce coût parasite. Un second coût anormal provenait du changement de rôle dans le \texttt{population\_guard} : les fonctions de mutation appelaient indirectement des vérifications globales de rôles à chaque activation ou inactivation locale. Le profil profond des mutations a montré que ces appels dominaient presque totalement le coût du \texttt{population\_guard}.

\subsection{Optimisations OpenMP validées}

Les optimisations finalement conservées sont celles qui ont passé les validations discriminantes. Elles sont résumées ci-dessous.

\begin{center}
\small
\begin{tabular}{p{0.24\linewidth}p{0.44\linewidth}p{0.24\linewidth}}
\toprule
Zone du code & Modification retenue & Effet principal \\
\midrule
Réponse virielle fermée & Retour immédiat si la fermeture de capacité est désactivée, avant tout dépôt masse/cellule. & Suppression d'un coût inutile dans les cas sans viriel. \\
Projection Q6/CG & Réorganisation du noyau elliptique, pré-calculs et fusion de certaines opérations vectorielles. & Réduction du coût du solveur elliptique tout en conservant les résidus Q6. \\
Resampling sans édition de rôles & Suppression des reconstructions de pool après remap ou renormalisation thermique lorsque seules les masses ou vitesses changent. & Moins de scans particulaires inutiles. \\
\texttt{population\_guard} & Listes compactes de cellules candidates overfull/underfull, conservant l'ordre croissant des cellules. & Évite de parcourir toutes les cellules lorsque seules quelques cellules sont candidates. \\
Mutations de rôles & Écriture locale préconditionnée des rôles \texttt{Fluid}/\texttt{Inactive} dans les mutations, sans revalidation globale répétée. & Suppression du verrou principal du \texttt{population\_guard}. \\
Dépôt post-thermal & Rafraîchissement spécialisé des moments et vitesses après renormalisation thermique, sans reconstruire les listes candidates ni les plans. & Forte baisse du coût \texttt{post\_thermal\_deposit}. \\
Dépôt post-guard & Réutilisation stricte des \texttt{cellId} après \texttt{population\_guard}, avec mise à jour locale des particules activées/inactivées, mais dépôt complet conservé. & Réduction du coût de \texttt{post\_guard\_deposit} sans changer les listes/plans. \\
Diagnostics internes & Mode light par défaut, profils réactivables par \texttt{MPCD\_INTERNAL\_PROFILES=1}, puis suppression du coût runtime des profils lorsque désactivés. & Version portable sans fichiers de profilage internes en usage normal. \\
\bottomrule
\end{tabular}
\end{center}

Deux optimisations ont été explicitement rejetées malgré des gains apparents. Le fast path du \texttt{mass\_guard}, qui évitait la projection de masse lorsque la cellule paraissait déjà admissible, a modifié la trajectoire discrète du resampling dans le test discriminant. De même, une première tentative de refresh local complet du dépôt post-guard a accéléré le code mais a changé certaines décisions de resampling sur Taylor--Green, obstacle ouvert et piston viriel. Ces rejets sont importants : ils montrent que le critère de validation n'était pas seulement temporel, mais bien fonctionnel.

\subsection{Résultats de validation et gains mesurés}

La validation discriminante finale du jalon OpenMP optimisé a comparé la version de référence et la version optimisée sur les quatre cas standard, avec \(64\times64\) cellules, \(\gamma=20\), \(1000\) pas et \(8\) threads. Les quatre cas ont passé la comparaison avec zéro métrique échouée sur 76 métriques par cas, soit 304 métriques comparées sans divergence. Les speedups mesurés sur cette campagne étaient de l'ordre de \(1.45\times\) pour l'obstacle ouvert, \(1.52\times\) pour le piston viriel, \(1.96\times\) pour Poiseuille et \(1.47\times\) pour Taylor--Green. Ces valeurs dépendent naturellement de la machine, du bruit d'ordonnancement OpenMP et des scripts utilisés, mais elles donnent l'ordre de grandeur du gain obtenu sans changer la logique physique.

Sur le profil \texttt{q6\_resampling} court, le verrou \texttt{population\_guard} a été presque éliminé comme poste dominant. Après l'écriture locale des rôles, son coût est devenu secondaire devant Q6, \texttt{mass\_guard}, collision et dépôts. Le dépôt post-thermal a également cessé d'être un verrou significatif après spécialisation. Le dépôt post-guard a été réduit sans recourir au refresh local non équivalent : la réutilisation stricte des \texttt{cellId} conserve le dépôt complet et donc les classifications/listes/plans, mais évite une partie substantielle du coût de calcul cellulaire.

\subsection{Réorganisation des diagnostics}

Une seconde étape a consisté à séparer les diagnostics de développement des diagnostics fluides de production. Les fichiers de profilage internes ont été désactivés par défaut dans la branche \texttt{clean/openmp-light}. Sont notamment concernés :
\begin{itemize}[nosep]
  \item \texttt{phase\_profile} et \texttt{q6\_cg\_profile} ;
  \item \texttt{deposit\_profile} et \texttt{resampling\_guard\_profile} ;
  \item \texttt{post\_guard\_profile} et les diagnostics d'équivalence post-guard.
\end{itemize}
 Les sorties physiques et numériques générales restent conservées : masse, nombres de particules fluides/inactives, moment global, température effective, statistiques de population, résidus Q6 utiles, compteurs globaux de resampling et diagnostics viriels lorsque la fermeture est active.

La variable d'environnement \texttt{MPCD\_INTERNAL\_PROFILES=1} permet encore de réactiver les profils internes dans la branche light. Cela préserve une possibilité de diagnostic ponctuel sans encombrer les runs de production. Un patch complémentaire a ensuite évité le coût runtime des timers et accumulateurs internes lorsque cette variable n'est pas active. Des scripts de smoke test vérifient les deux comportements : aucun fichier de profilage interne en mode light, et production effective de profils lorsque \texttt{MPCD\_INTERNAL\_PROFILES=1} est explicitement défini.

Enfin, les scripts de profilage historiques, les rapports intermédiaires et les CSV de comparaison ont été déplacés vers \texttt{dev\_history/}. La branche light garde donc les scripts de build, de validation générique et de smoke test, mais n'expose plus les nombreux scripts de développement liés aux jalons d'optimisation. Une validation finale \texttt{clean/openmp-light} contre \texttt{SRC\_openMP\_optimized} a confirmé l'absence de modification fonctionnelle : les quatre cas standard ont passé la comparaison avec zéro métrique échouée.

\subsection{Conséquences pour la suite GPU}

Cette phase d'optimisation OpenMP fournit aussi une base utile pour le port GPU. Les profils conservés dans la branche instrumentée ont identifié les noyaux qui deviendront prioritaires sur de grandes grilles : projection Q6, dépôt particules--cellules, collision par cellule et éventuellement construction d'index cellule--particules. À l'échelle \(1000\times1000\) avec \(\gamma\simeq20\), le nombre de particules atteint l'ordre de \(2\times10^7\), et les gains purement OpenMP auront probablement des rendements décroissants. La branche optimisée/instrumentée doit donc rester le support du développement GPU, tandis que la branche \texttt{clean/openmp-light} fournit une référence CPU portable et validée.



\section{Portage CUDA du SRC classic, architecture résidente et thermostat boundary-aware}

\subsection{Définition du périmètre : SRC classic et fermeture liquide}

Le chantier GPU a été mené sur la branche \texttt{SRC\_GPU}, dérivée de la branche OpenMP allégée. Il est essentiel de fixer la terminologie, car elle conditionne l'architecture du code. Dans cette branche, on appelle \textbf{SRC classic} le pas physique complet
\begin{equation}
  \text{advection/streaming}
  + \text{décalage aléatoire de grille}
  + \text{rotation/collision SRC}
  + \text{thermostat}.
\end{equation}
La fermeture liquide incompressible ou faiblement compressible est un module additionnel :
\begin{equation}
  \text{SRC classic}
  + \text{Q6}
  + \text{resampling pondéré}
  + \text{fermeture virielle/capacité}.
\end{equation}
Le résultat CUDA obtenu à ce stade concerne donc le SRC classic. Il ne doit pas être confondu avec une migration CUDA complète de Q6, du resampling ou du viriel. Cette séparation a été conservée volontairement afin de ne pas verrouiller les chantiers futurs, en particulier la migration Q6 CUDA.

Le choix de CUDA, plutôt qu'un backend abstrait OpenCL/CUDA, a été retenu pour ce cycle de développement afin de réduire le coût d'intégration et de concentrer l'effort sur la résidence mémoire, les transferts et les chemins physiques. Le backend OpenMP reste disponible explicitement, ce qui permet de comparer les trajectoires et de conserver une référence CPU robuste.

\subsection{Principe d'architecture : état persistant et chemins spécialisés}

La difficulté principale du portage n'est pas seulement d'écrire des noyaux CUDA isolés. Un port naïf qui transfère toutes les particules vers le GPU à chaque sous-étape, puis les redescend immédiatement vers l'hôte, reproduit le calcul mais perd l'essentiel du gain. La stratégie retenue est donc une architecture \emph{résidente} : les particules et les espaces de travail cellule restent autant que possible côté GPU pendant la séquence SRC classic.

On distingue plusieurs couches.
\begin{itemize}[nosep]
  \item L'état particulaire persistant contient positions, vitesses, masses, rôles et métadonnées nécessaires aux noyaux de streaming, frontière, collision et thermostat.
  \item L'espace de travail cellule persistant contient les dépôts, moments, vitesses moyennes, énergies thermiques et compteurs de population nécessaires à la collision SRC et au thermostat.
  \item Les chemins de streaming/frontière sont spécialisés par famille géométrique : périodique, parois simples, obstacle rectangle, piston/mobile wall, inlet/outlet full-face et inlet/outlet segmenté.
  \item Le noyau de collision SRC applique la rotation par cellule en utilisant les dépôts cellule, les signes aléatoires et les informations de rôle/matériau.
  \item Le thermostat CUDA reconstruit l'énergie thermique relative et applique le rescale cellulaire lorsque le chemin physique le permet.
\end{itemize}

Cette organisation explique la différence entre deux types de chemins. Pour un cas SRC classic-only, le chemin cible est entièrement fusionné côté GPU :
\begin{equation}
  \text{frontière/streaming CUDA}
  \longrightarrow
  \text{collision SRC CUDA}
  \longrightarrow
  \text{thermostat CUDA}.
\end{equation}
Pour un cas où Q6, resampling, viriel ou réponse de capacité modifient les vitesses ou les masses après la collision SRC, le thermostat fusionné n'est pas correct. L'ordre conservé est alors
\begin{equation}
  \text{collision SRC CUDA}
  \longrightarrow
  \text{étapes CPU/OpenMP de fermeture liquide}
  \longrightarrow
  \text{thermostat CUDA post-CPU}.
\end{equation}
Ce second chemin est notamment celui validé pour le piston/mobile wall. Il préserve la cohérence physique et la possibilité de réactiver Q6 CPU/OpenMP, resampling CPU/CUDA et viriel.

\subsection{Étapes du portage CUDA}

Le développement CUDA a progressé de manière incrémentale. Les premiers jalons ont validé les briques élémentaires : détection de toolchain CUDA, backend de projection expérimental, dépôts cellule, thermostat isolé, collision SRC isolée, puis état particulaire persistant. Les jalons suivants ont introduit les espaces cellule persistants, les chemins streaming/frontières, les obstacles, l'inlet/outlet et les variantes résidentes.

La séquence validée peut être résumée ainsi.
\begin{center}
\begin{tabular}{p{0.18\linewidth}p{0.72\linewidth}}
\toprule
Jalons & Rôle \\
\midrule
0186--0208 & Mise en place du backend CUDA, tests Taylor--Green, dépôts cellule, thermostat isolé et premiers chemins combinés. \\
0209--0225 & Collision SRC CUDA, état particulaire persistant, ponts host/device, workspace cellule persistant et premières validations end-to-end. \\
0227--0244 & Briques CUDA du resampling : garde pauvre/riche, plan de transfert, opérations extraction/insertion et état persistant. Ces briques préparent le futur, mais ne constituent pas encore la fermeture liquide full CUDA. \\
0245--0249b & Streaming/frontières CUDA : périodique, wall-simple, rectangle immergé, piston/mobile wall, inlet/outlet full-face et inlet/outlet segmenté. \\
0251--0259 & État cellule partagé, collision SRC persistante, premières optimisations de transfert et thermostat CUDA validé initialement en périodique. \\
0260--0275 & Passage au SRC classic résident optimisé : périodique, wall-simple, solide rectangle, piston, inlet/outlet full-face, inlet/outlet segmenté, validations consolidées, instrumentation performance et backend CUDA principal. \\
0276--0281 & Extension du thermostat CUDA aux frontières : wall-simple, solide, piston/mobile wall, inlet/outlet full-face et segmenté, puis validation consolidée. \\
0284--0286 & Ajout du solide circulaire CUDA, d'abord en périodique puis avec inlet/outlet full-face pour le cas Von Karman, et consolidation du jalon SRC classic full CUDA. \\
0288--0293 & Stabilisation des ouvertures CUDA : clipping géométrique des segments, modes explicites de sortie \texttt{neumann}, \texttt{equilibrium\_flux} et \texttt{forced\_flux}, commutateur physique unique \texttt{thermostatEnable}, puis pré-extraction outlet avant refill inlet. \\
\bottomrule
\end{tabular}
\end{center}

Le point 0276--0281 est décisif : avant cette étape, le code disposait d'un chemin SRC CUDA résident robuste, mais le thermostat CUDA n'était réellement validé qu'en périodique. Après cette étape, le thermostat appartient bien au chemin SRC classic CUDA pour les familles de conditions limites validées.

\subsection{Thermostat CUDA : séparation collision/thermostat près des frontières}

Le correctif physique principal concerne les particules virtuelles de paroi ou de solide. Pour la collision SRC, les moyennes cellule peuvent inclure des particules virtuelles, car elles servent à imposer le couplage wall/no-slip/thermal attendu. En revanche, le thermostat de fluide ne doit pas utiliser ces particules virtuelles pour reconstruire l'énergie thermique réelle des particules transportées. La règle validée est donc
\begin{align}
  \text{moyenne de collision} & = \text{particules réelles} + \text{particules virtuelles éventuelles}, \\
  \text{moyenne de thermostat} & = \text{particules réelles post-collision uniquement}.
\end{align}

Cette distinction explique l'échec initial du thermostat CUDA hors périodique et le correctif 0276. Dans le chemin fusionné, il ne fallait pas réutiliser aveuglément les moments de collision enrichis par les particules virtuelles ; il fallait reconstruire, après rotation SRC, un dépôt thermostat fondé sur les seules particules réelles. Une fois cette séparation imposée, le même principe s'est étendu aux parois simples et aux solides fixes.

Pour le piston/mobile wall, une autre contrainte apparaît. Le cas validé peut activer Q6, resampling, viriel ou réponse de capacité entre collision et thermostat. Il serait alors incorrect d'appliquer le thermostat immédiatement après la collision CUDA. La validation piston utilise donc le thermostat CUDA en mode post-CPU : les étapes CPU/OpenMP intermédiaires modifient l'état, puis le thermostat CUDA agit sur l'état mis à jour.


\subsection{Conditions limites couvertes}

Les familles de conditions limites validées couvrent maintenant les cas principaux du dépôt. Le tableau suivant utilise des noms de cas abrégés afin de mettre l'accent sur la famille physique plutôt que sur le nom exact des scripts.
\begin{center}
\begin{tabular}{p{0.24\linewidth}p{0.27\linewidth}p{0.39\linewidth}}
\toprule
Famille & Cas de validation & Chemin CUDA validé \\
\midrule
Périodique & Taylor--Green / SRC périodique & Streaming périodique, collision SRC, thermostat CUDA. \\
Wall-simple & Poiseuille wall & Paroi simple, collision SRC, reconstruction thermostat réelle, thermostat CUDA. \\
Solide fixe & Rectangle périodique & Gestion solide/rectangle, collision SRC, thermostat CUDA. \\
Solide circulaire & Cercle périodique & Réflexion cercle CUDA, fraction solide circulaire, collision SRC et thermostat CUDA. \\
Piston/mobile wall & Piston legacy & Mobile wall et SRC CUDA ; thermostat CUDA post-CPU pour préserver Q6/resampling/viriel. \\
Inlet/outlet full-face & Obstacle rectangle ouvert & Inlet/outlet résident full-face, collision SRC, thermostat CUDA fusionné. \\
Inlet/outlet segmenté & U-turn segmenté & Inlet/outlet segmenté résident, état partagé, collision SRC, thermostat CUDA fusionné. \\
Cercle + inlet/outlet & Cylindre/Von Karman & Full-face inlet/outlet, réservoir dur, cercle solide CUDA, collision SRC et thermostat CUDA fusionné. \\
Régimes outlet explicites & 0291--0293 & Sortie passive \texttt{neumann}, équilibrage local \texttt{equilibrium\_flux} et extraction imposée \texttt{forced\_flux}, applicables aux chemins full-face et segmentés résidents. \\
\bottomrule
\end{tabular}
\end{center}

Le correctif 0280c est particulièrement important pour les ouvertures. Les chemins résidents inlet/outlet refusaient initialement \texttt{thermostatEnable=true}. Le garde a été relâché uniquement dans le cas sûr : thermostat CUDA persistant demandé, Q6/resampling/viriel/capacité désactivés, et chemin classic-only. Ainsi, l'activation du thermostat ne désactive plus silencieusement le chemin résident, mais les modules de fermeture liquide restent protégés.

\subsection{Ouvertures CUDA complètes : full-face, segments et régimes de sortie}

Les conditions limites ouvertes CUDA ont été complétées au-delà du simple couple \texttt{inlet}/\texttt{outlet} géométrique. Le code distingue maintenant deux aspects : la \emph{géométrie} de l'ouverture et le \emph{régime physique} appliqué à la sortie. La géométrie peut être une face complète, par exemple \texttt{bcLeft=inlet} et \texttt{bcRight=outlet}, ou bien une liste de segments portés par une face solide. Dans les deux cas, le chemin cible reste résident CUDA lorsque le cas est SRC classic-only : les particules sont advectées, traitées par les frontières ouvertes, éventuellement insérées ou extraites par le réservoir, puis transmises directement à la collision SRC et au thermostat sans retour hôte intermédiaire.

Le traitement segmenté a nécessité un correctif géométrique spécifique. Les premières versions sélectionnaient les cellules de réservoir à partir de leur centre, puis inséraient les particules dans toute la cellule. Lorsqu'une borne de segment coupait une cellule, des particules pouvaient donc être créées légèrement au-delà de la portion réellement ouverte, ce qui produisait une singularité de densité localisée à la lèvre du segment. Le correctif 0288 remplace cette logique par un clipping d'intervalle : une cellule de bord est pondérée par la fraction réellement ouverte, et les positions injectées sont tirées uniquement dans la sous-cellule incluse dans le segment. La suppression de cette singularité confirme que l'inlet segmenté doit être interprété comme une géométrie continue discrétisée, non comme une sélection de centres de cellules.

Les régimes de sortie sont maintenant explicites via
\begin{equation}
  \texttt{openBoundaryOutletMode}
  \in
  \{\texttt{neumann},\texttt{equilibrium\_flux},\texttt{forced\_flux}\}.
\end{equation}
Leur sens est le suivant.
\begin{center}
\begin{tabular}{p{0.22\linewidth}p{0.68\linewidth}}
\toprule
Mode & Interprétation \\
\midrule
\texttt{neumann} & Sortie passive. Le bord laisse partir les particules qui atteignent naturellement l'outlet ; aucune extraction artificielle n'est imposée. Ce mode convient aux canaux ouverts classiques, aux sillages et aux sorties suffisamment aval. \\
\texttt{equilibrium\_flux} & Équilibrage local. Le code tente d'extraire à l'outlet une quantité comparable au gain dû au refill inlet du pas courant. L'extraction reste plafonnée par la masse ou le nombre de particules effectivement disponibles dans la couche outlet ; elle ne garantit donc pas une masse globale constante si la sortie est localement pauvre. \\
\texttt{forced\_flux} & Extraction imposée par l'utilisateur, décorrélée du flux inlet. La quantité extraite est spécifiée par un flux ou par un budget par pas, et le code prélève localement dans la couche outlet tant que des particules y sont disponibles. Ce mode permet de modéliser aspiration, drainage ou pompage local, en laissant à l'utilisateur le choix d'équilibrer ou non les flux entrant et sortant. \\
\bottomrule
\end{tabular}
\end{center}

Les paramètres associés restent volontairement peu nombreux. Un budget d'extraction peut être donné en masse ou en nombre de particules :
\begin{itemize}[nosep]
  \item \texttt{openBoundaryOutletForcedMassFlux} : flux massique demandé ; la masse extraite par pas est proportionnelle à $\Delta t$.
  \item \texttt{openBoundaryOutletForcedMassPerStep} : masse directement demandée par pas, utile pour les essais contrôlés.
  \item \texttt{openBoundaryOutletForcedParticleFlux} : flux en nombre de particules, converti en budget par pas.
  \item \texttt{openBoundaryOutletForcedParticlesPerStep} : nombre de particules à extraire par pas ; prioritaire pour les tests à masse particulaire uniforme.
  \item \texttt{openBoundaryOutletForcedLayerCells} : épaisseur, en cellules, de la couche outlet dans laquelle l'extraction locale est autorisée.
\end{itemize}

Le mode \texttt{equilibrium\_flux} a également conduit à une correction d'ordre d'application. Dans un réservoir hard inlet, le refill consomme des slots \texttt{Inactive}. Si l'extraction outlet équilibrante est effectuée après le refill, le réservoir peut s'épuiser avant que l'outlet n'ait libéré les slots nécessaires. Le correctif 0293 applique donc l'extraction outlet d'équilibre ou forcée avant le refill hard inlet, puis reconstruit le pool disponible. L'ordre pertinent devient
\begin{equation}
  \text{streaming/frontières}
  \longrightarrow
  \text{extraction outlet}
  \longrightarrow
  \text{refill hard inlet}
  \longrightarrow
  \text{collision SRC}
  \longrightarrow
  \text{thermostat}.
\end{equation}
Cette modification ne transforme pas \texttt{equilibrium\_flux} en condition de masse globale parfaite : si la couche outlet ne contient pas assez de particules, l'extraction reste partielle et la masse totale peut augmenter. Elle rend en revanche le recyclage du pool inactif cohérent lorsque l'outlet dispose localement de particules à extraire.

La gestion du thermostat a été clarifiée en parallèle. Le paramètre \texttt{thermostatEnable} est le commutateur physique unique. Lorsqu'il vaut \texttt{false}, le thermostat est désactivé en CPU comme en CUDA, même si des variables d'environnement CUDA sélectionnent des chemins rapides. Lorsqu'il vaut \texttt{true}, les variables d'environnement ne choisissent que le backend et le mode de résidence. Cette règle est importante pour les simulations d'aspiration : \texttt{forced\_flux} avec thermostat actif représente une extraction isotherme, tandis que \texttt{forced\_flux} sans thermostat laisse évoluer l'énergie du gaz restant.

Enfin, les messages d'épuisement du réservoir ont été rendus explicites. Lorsque le pool \texttt{Inactive} ne suffit plus, le code signale maintenant l'étape concernée, le type de réservoir, le fait que l'append GPU dynamique est désactivé, puis les compteurs utiles. Ce choix garde le chemin résident déterministe : les longs runs doivent soit préallouer un pool suffisant, soit régler les flux et l'épaisseur outlet de manière à recycler assez de particules.

\subsection{Validation consolidée 0281}

Le validateur 0281 regroupe les validations thermostat CUDA boundary-aware. Il exécute cinq familles, chacune sur deux tailles de grille, soit dix lignes de synthèse. Toutes les lignes ont passé la comparaison avec zéro métrique échouée sur 76 métriques comparées. La température effective après thermostat est ramenée à la cible \(10^{-3}\) dans tous les cas.


\begin{center}
\small
\begin{tabular}{p{0.27\linewidth}p{0.28\linewidth}p{0.14\linewidth}p{0.10\linewidth}p{0.11\linewidth}}
\toprule
Validateur & Cas & Grille & Statut & Speedup \\
\midrule
wall 0276 & Poiseuille wall & 64x64 s300 & PASS & 0.29 \\
wall 0276 & Poiseuille wall & 128x128 s300 & PASS & 0.42 \\
solid 0277 & Rectangle périodique & 64x64 s300 & PASS & 0.63 \\
solid 0277 & Rectangle périodique & 128x128 s300 & PASS & 0.89 \\
piston 0278 & Piston legacy & 64x64 s300 & PASS & 0.55 \\
piston 0278 & Piston legacy & 128x128 s300 & PASS & 0.72 \\
full-face IO & Obstacle ouvert & 64x64 s300 & PASS & 1.72 \\
full-face IO & Obstacle ouvert & 128x128 s300 & PASS & 3.14 \\
segmenté IO & U-turn segmenté & 64x64 s300 & PASS & 2.55 \\
segmenté IO & U-turn segmenté & 128x128 s300 & PASS & 5.35 \\
\bottomrule
\end{tabular}
\end{center}

Ces speedups doivent être interprétés avec prudence. Les runs wall/solid/piston sont courts, et le coût fixe de lancement, synchronisation, génération d'état et comparaison peut dominer. L'information essentielle est donc d'abord la non-régression physique : mêmes métriques que le chemin CPU de référence et thermostat à la cible. Les gains deviennent nets pour les inlet/outlet résidents, où la résidence GPU et l'état partagé évitent plus de transferts.

\subsection{Politique de garde et conséquences pour les futurs développements}

L'état atteint peut être formulé ainsi : la branche \texttt{SRC\_GPU} fournit maintenant un code \textbf{SRC/MPCD classic full CUDA résident} pour les cas classic-only validés, incluant advection/streaming, décalage de grille, collision/rotation SRC, conditions limites avancées et thermostat. Le chemin OpenMP reste disponible explicitement et sert de référence de validation.

La fermeture liquide reste volontairement séparée. Lorsque \texttt{projectionEnable}, \texttt{resamplingEnable}, la réponse de capacité ou le viriel sont actifs, le code ne doit pas forcer un raccourci CUDA classic-only qui contournerait ces modules. Le bon invariant est :
\begin{itemize}[nosep]
  \item classic-only : chemin résident CUDA fusionné ;
  \item Q6/resampling/viriel : étapes CPU/OpenMP ou hybrides conservées, puis thermostat CUDA post-CPU si activé ;
  \item futur Q6 CUDA : chantier séparé portant sur le solveur elliptique, les opérateurs divergence/gradient, les masques de frontière et la synchronisation avec les autres modules.
\end{itemize}

Ce découplage est une propriété d'architecture, pas une limitation accidentelle. Il permet d'utiliser immédiatement un SRC classic CUDA physique et rapide pour les cas validés, tout en gardant une trajectoire claire vers une fermeture liquide entièrement accélérée.



\section{Resampling CUDA post-SRC : implémentation, paramètres et validation}

\subsection{Principe physique et ordre d'exécution}

Le chantier CUDA du resampling a conduit à préciser un point important d'interprétation. Le resampling n'est pas un processus physique supplémentaire placé avant la collision ; il ne doit pas être lu comme une création ou destruction de particules matérielles. La dynamique physique validée du SRC/MPCD reste celle de l'advection, des conditions limites, du décalage aléatoire de grille et de la collision/rotation SRC. Le resampling CUDA est donc placé \emph{après} le SRC classic, comme dans le développement MATLAB et OpenMP de la méthode pondérée.

L'ordre nominal retenu est :
\begin{equation}
\begin{aligned}
&\text{streaming CUDA} \\
&\quad \to \text{conditions limites CUDA : murs, solides, inlet/outlet, piston} \\
&\quad \to \text{collision SRC CUDA avec shift collisionnel interne} \\
&\quad \to \text{thermostat CUDA si } \texttt{thermostatEnable=true} \\
&\quad \to \text{Q6 CPU/OpenMP éventuel si } \texttt{projectionEnable=true} \\
&\quad \to \text{resampling CUDA post-SRC : survey, reconditionnement, population guard}.
\end{aligned}
\end{equation}

Ce choix garantit que le resampling maintient la validité statistique de l'état fluide produit par SRC sans redéfinir les collisions physiques. Le shift reste attaché à l'opérateur de collision ; le resampling travaille sur la grille physique non shiftée. Pour les domaines ouverts, toute variation physique nette de masse ou d'impulsion doit provenir des conditions inlet/outlet, non du resampling. Dans un domaine fermé, le resampling doit conserver la masse totale, l'impulsion et l'énergie relative locale à la précision numérique.

\subsection{Jalons d'implémentation 0294--0299}

Le développement a été volontairement progressif afin de ne pas casser le chemin SRC classic full CUDA résident validé.

\begin{itemize}[nosep]
  \item \textbf{0294 -- audit.} L'audit a distingué les briques historiques de resampling CUDA déjà présentes, notamment classification poor/rich, plans de transfert, extraction/insertion persistante, rôles \texttt{Fluid}/\texttt{Inactive}/\texttt{Latent}, workspace cellule et état particulaire CUDA partagé. La conclusion d'architecture est que les anciennes briques étaient utiles mais ne formaient pas encore un module post-SRC non destructif.
  \item \textbf{0295 -- survey post-SRC.} Un diagnostic CUDA passif calcule, après SRC, $N_c$, $M_c$, $P_{x,c}$, $P_{y,c}$, $U_{x,c}$, $U_{y,c}$, l'énergie relative et les classes \texttt{empty}/\texttt{poor}/\texttt{target}/\texttt{rich}. Aucun rôle, masse ou vitesse n'est modifié. Les validations strictes TG, Poiseuille, backward step et U-box segmentée sont passées. Von Karman est conservé comme stress test car le thermostat CUDA boundary-aware n'est pas bit-reproductible OFF/OFF dans ce cas, alors que le même cas sans thermostat est reproductible.
  \item \textbf{0296 -- reconditionnement conservatif des masses.} Le module réduit une dispersion excessive des masses particulaires sans changer le support. Dans chaque cellule humide non vide, les masses sont rapprochées de $M_c/N_c$ avec une force \texttt{strength}, puis l'impulsion cellulaire est restaurée. Sur les états à masses uniformes, le module est neutre et produit uniquement son CSV de diagnostic.
  \item \textbf{0297 -- garde local de population.} Le premier module mutateur agit sur les cellules $0<N_c<N_{\min}$ par split de particules représentatives et sur les cellules $N_c>N_{\max}$ par merge local. Le split ne crée pas de fluide : une particule numérique pondérée est remplacée par deux porteurs statistiques dont la masse et l'impulsion totales égalent celles du donneur.
  \item \textbf{0298 -- restauration de l'énergie relative.} Après split/merge, la masse et l'impulsion locales sont conservées, puis les vitesses relatives à $U_c$ peuvent être rescalées pour restaurer $K_{\mathrm{rel},c}=\frac12\sum_i m_i\lVert v_i-U_c\rVert^2$.
  \item \textbf{0299 -- filtrage boundary-aware.} Les cellules candidates sont filtrées près des ouvertures et, sur demande, près des parois ou solides. Le halo ouvert vaut 1 cellule par défaut, afin de ne pas mélanger le contrôle de support interne avec les réservoirs hard inlet/outlet.
\end{itemize}

\subsection{Formulation locale du split, du merge et de la restauration}

Pour une cellule pauvre, le split minimal sélectionne un donneur fluide local de masse $m_d$, position $x_d$ et vitesse $v_d$, puis active un slot \texttt{Inactive} :
\begin{equation}
  m_d' = m_d - \delta m, \qquad
  m_{\mathrm{new}} = \delta m, \qquad
  x_{\mathrm{new}}=x_d, \qquad
  v_{\mathrm{new}}=v_d.
\end{equation}
La masse, l'impulsion et l'énergie portée par la particule donneuse sont alors conservées exactement avant la restauration globale de cellule. Pour une cellule riche, deux particules locales sont fusionnées par moyenne pondérée :
\begin{equation}
  m' = m_1+m_2, \qquad
  v' = \frac{m_1 v_1 + m_2 v_2}{m_1+m_2}.
\end{equation}
Cette opération conserve masse et impulsion mais peut réduire l'énergie relative. C'est la raison du module 0298 : après mutation, le code compare $K_{\mathrm{rel},c}^{\mathrm{before}}$ et $K_{\mathrm{rel},c}^{\mathrm{after}}$, puis applique si nécessaire
\begin{equation}
  v_i \leftarrow U_c + s_c\,(v_i-U_c), \qquad
  s_c = \sqrt{\frac{K_{\mathrm{rel},c}^{\mathrm{target}}}{K_{\mathrm{rel},c}^{\mathrm{current}}}},
\end{equation}
avec un bornage de $s_c$ par \texttt{MPCD\_CUDA\_RESAMPLING\_MOMENT\_RESTORE\_0298\_MAX\_SCALE}. L'opération est donc un remaillage statistique conservatif, non une correction cachée du champ de vitesse.

\subsection{Options utilisateur et variables d'environnement spécifiques 0295--0314}

Les listings de paramètres 0292 fournis avec ce rapport décrivent le socle CUDA 0286/0292 : backend, collisions résidentes, thermostat, inlet/outlet, solides, Q6 préservé et paramètres \texttt{params.kv}. Les options 0295--0303 ci-dessous sont des compléments spécifiques au resampling CUDA post-SRC et aux scripts de démonstration.


\begin{itemize}[nosep]
  \item \path{MPCD_CUDA_RESAMPLING_SUPPORT_SURVEY_0295} : active le survey post-SRC non-mutant ; cadence via \path{MPCD_CUDA_RESAMPLING_SUPPORT_SURVEY_0295_EVERY}.
  \item \path{MPCD_CUDA_RESAMPLING_MASS_RECONDITION_0296} : active le reconditionnement conservatif des masses ; force via \path{MPCD_CUDA_RESAMPLING_MASS_RECONDITION_0296_STRENGTH}.
  \item \path{MPCD_CUDA_RESAMPLING_POPULATION_GUARD_0297} : active le split/merge local post-SRC ; cadence via \path{MPCD_CUDA_RESAMPLING_POPULATION_GUARD_0297_EVERY}.
  \item \path{MPCD_CUDA_RESAMPLING_POPULATION_GUARD_0297_NMIN}, \path{..._NTARGET}, \path{..._NMAX} : triplet de population. Le réglage nominal provisoire est \texttt{12:20:32}.
  \item \path{MPCD_CUDA_RESAMPLING_MOMENT_RESTORE_0298} : restaure l'énergie relative cellulaire ; bornage via \path{MPCD_CUDA_RESAMPLING_MOMENT_RESTORE_0298_MAX_SCALE}.
  \item \path{MPCD_CUDA_RESAMPLING_POPULATION_GUARD_0299_BOUNDARY_AWARE} : active le filtrage boundary-aware ; l'exclusion des ouvertures est contrôlée par \path{MPCD_CUDA_RESAMPLING_POPULATION_GUARD_0299_OPEN_BOUNDARY_HALO_CELLS}.
\end{itemize}


Les scripts de démonstration ajoutés ou mis à jour en 0303 utilisent des alias plus courts afin de faciliter les comparaisons visuelles : \texttt{RESAMPLING\_ENABLE}, \texttt{RESAMPLING\_SURVEY\_ENABLE}, \texttt{GUARD\_NMIN}, \texttt{GUARD\_NTARGET}, \texttt{GUARD\_NMAX}, \texttt{GUARD\_EVERY}, \texttt{RESTORE\_ENABLE}, \texttt{BOUNDARY\_AWARE} et \texttt{OPEN\_BOUNDARY\_HALO\_CELLS}. Ces variables ne sont pas des clés \texttt{params.kv} ; elles configurent les scripts et exportent les variables \texttt{MPCD\_CUDA\_...} appropriées.

\subsection{Sorties et diagnostics produits}

Le module resampling CUDA ajoute plusieurs sorties légères, placées dans le \texttt{outputDir} du run ou dans \texttt{dev\_history/artifacts/} pour les campagnes de validation.

\begin{table}[htbp]
\centering
\footnotesize
\begin{tabular}{p{0.31\linewidth}p{0.59\linewidth}}
\toprule
Fichier & Contenu principal \\
\midrule
\path{cuda_resampling_support_survey_0295.csv} & Population et moments cellulaires post-SRC : \texttt{emptyCells}, \texttt{poorCells}, \texttt{richCells}, \texttt{targetBandCells}, \texttt{meanNActive}, \texttt{stdNActive}, \texttt{minNWet}, \texttt{maxNWet}, \texttt{massRelRmsWet}, $K_{\mathrm{rel}}$ et $k_BT$ pondéré. \\
\path{cuda_resampling_mass_recondition_0296.csv} & Nombre de cellules et particules traitées, cellules sautées ou bornées, erreur de masse et d'impulsion après reconditionnement. \\
\path{cuda_resampling_population_guard_0297.csv} & Compteurs \texttt{splitApplied}, \texttt{mergeApplied}, erreurs maximales de masse, impulsion et $K_{\mathrm{rel}}$, budgets avant/après, exclusions boundary-aware 0299. \\
\path{cuda_resampling_active_sweep_0300_*.csv} & Résumé par run, comparaison au classic et manifest de la campagne active TG/Poiseuille/step/U-box. \\
\path{cuda_resampling_backward_step_long_0301_*.csv} & Séries temporelles et comparaisons longues sur marche arrière pour plusieurs vitesses d'entrée et triplets $N_{\min}/N_{\mathrm{target}}/N_{\max}$. \\
\bottomrule
\end{tabular}
\caption{Diagnostics ajoutés par la validation resampling CUDA 0295--0301.}
\end{table}

\subsection{Scripts de démonstration et post-traitement}

Les quatre scripts de démonstration CUDA historiques ont été rendus compatibles resampling en 0303 : Taylor--Green forcé, Poiseuille périodique forcé, backward step inlet/outlet et U-box segmentée. Ils peuvent être lancés avec \texttt{RESAMPLING\_ENABLE=0} pour obtenir une référence SRC classic avec survey passif, ou avec \texttt{RESAMPLING\_ENABLE=1} pour activer le triplet nominal \texttt{12:20:32}. Les répertoires sont séparés en \texttt{classic\_survey/} et \texttt{resampling\_guard\_nmin12\_nt20\_nmax32/} afin de permettre une comparaison directe par les outils MATLAB/post-traitement existants.

Un script Von Karman séparé, \texttt{run\_demo\_src\_resampling\_cuda\_von\_karman\_cylinder\_0303.sh}, a été ajouté pour ne pas interférer avec le script local \texttt{0285} utilisé pour les essais de développement. Le cas Von Karman avec thermostat est conservé comme stress test physique ; pour les validations bit-à-bit strictes, le même cas est lancé avec thermostat désactivé car le thermostat CUDA boundary-aware présente une non-reproductibilité OFF/OFF mesurable sur ce cas complexe.

\subsection{Validation courte 0300 : quatre cas reproductibles}

La campagne 0300 a comparé \texttt{classic}, \texttt{mass0296} et trois gardes actifs sur quatre cas courts : Taylor--Green périodique, Poiseuille wall, backward step et U-box segmentée. Tous les runs se sont terminés avec \texttt{exitCode=0}. TG et Poiseuille restent essentiellement passifs, ce qui vérifie l'absence d'effet parasite sur des supports déjà bien conditionnés. Le backward step déclenche le guard et fournit le premier cas actif court ; la U-box courte reste peu sollicitante mais confirme la compatibilité avec les ouvertures segmentées.

\begin{table}[htbp]
\centering
\footnotesize
\begin{tabular}{p{0.23\linewidth}rrrrr}
\toprule
Cas 0300 & split & merge & poor & rich & erreur max. $K_{\mathrm{rel}}$ \\
\midrule
TG, \texttt{12:20:32} & 0 & 0 & 1 & 0 & $1.4\times10^{-17}$ \\
Poiseuille, \texttt{12:20:32} & 0 & 0 & 0 & 0 & $1.4\times10^{-17}$ \\
Backward step, \texttt{12:20:32} & 17 & 0 & 17 & 0 & $2.1\times10^{-17}$ \\
U-box segmentée, \texttt{12:20:32} & 0 & 0 & 0 & 0 & $1.4\times10^{-17}$ \\
\bottomrule
\end{tabular}
\caption{Synthèse du sweep actif court 0300 pour le réglage nominal provisoire \texttt{12:20:32}.}
\end{table}

\subsection{Validation longue 0301--0302 : marche arrière}

Le backward step a été choisi comme banc d'essai principal parce que les simulations classic perdent rapidement des particules derrière la marche lorsque la vitesse d'entrée augmente. Sur une grille modérée $96\times48$, avec $3000$ pas et $U_{\mathrm{in}}=0.30$, $0.45$ puis $0.60$, la référence classic développe respectivement $21$, $98$ et $191$ cellules humides vides en fin de run. Le réglage nominal \texttt{12:20:32} supprime ces cellules vides finales dans les trois cas.

\begin{table}[htbp]
\centering
\footnotesize
\begin{tabular}{c r r r r r r}
\toprule
$U_{\mathrm{in}}$ & empty classic & empty guard & poor classic & poor guard & stdN classic & stdN guard \\
\midrule
0.30 & 21 & 0 & 543 & 370 & 8.05 & 5.18 \\
0.45 & 98 & 0 & 845 & 632 & 10.56 & 6.37 \\
0.60 & 191 & 0 & 1074 & 775 & 12.81 & 7.33 \\
\bottomrule
\end{tabular}
\caption{Effet du resampling CUDA nominal \texttt{12:20:32} sur le support particulaire du backward step long 0301.}
\end{table}

Pour $U_{\mathrm{in}}=0.60$, le réglage nominal effectue $10832$ splits et $7759$ merges cumulés. Les erreurs maximales observées restent au bruit numérique : masse $3.6\times10^{-15}$, impulsion $1.8\times10^{-14}$ et énergie relative restaurée $2.8\times10^{-13}$. Le moment longitudinal final diffère du classic d'environ $142$ sur $P_x\simeq1.57\times10^4$, soit moins de un pour cent. Cette différence est acceptable pour un run actif où le support particulaire est volontairement reconditionné et où le classic souffre précisément de cellules vides persistantes.

La conclusion expérimentale de 0301--0302 est donc que \texttt{12:20:32} constitue le meilleur réglage nominal provisoire. Le triplet \texttt{10:20:34} est plus prudent mais moins correcteur ; \texttt{14:20:30} réduit davantage les cellules pauvres mais devient plus intrusif en nombre de mutations et en décalage global du moment. Le réglage nominal retenu est suffisamment actif pour supprimer les vides derrière la marche, tout en gardant les budgets $M/P/K_{\mathrm{rel}}$ propres.


\subsection{Diagnostics 0304--0306 : flags de support, classification géométrique et outliers de vitesse}

Après la validation 0301--0302, les visualisations longues ont montré que le seul comptage des cellules vides ne suffisait pas à qualifier la qualité du support. Certaines cellules, notamment près des parois ou des solides, pouvaient rester peuplées tout en présentant des vitesses cellulaires aberrantes. Cette observation a conduit à distinguer trois défauts différents : la cellule vide ($N_c=0$), la cellule pauvre mais non vide, et la cellule bien peuplée en nombre mais mal représentée du point de vue des masses ou des vitesses.

Le jalon 0304 ajoute un flag passif post-SRC/post-thermostat sur la grille physique non shiftée. Ce flag ne déclenche encore aucune correction ; il compte seulement les cellules humides vides, les cellules sous un seuil $N_c\leq N_{\mathrm{trigger}}$, les extrema de population et les coûts du diagnostic. L'intérêt de cette étape est double. D'une part, elle montre que le déclenchement adaptatif futur peut être fondé sur des informations déjà disponibles après le dépôt cellulaire, sans nouveau survey lourd. D'autre part, elle prépare un mécanisme commun qui pourra plus tard déclencher soit un guard anticipé, soit un refill de cellule vide.

Le jalon 0305 ajoute une classification géométrique des cellules problématiques : bulk, wall-adjacent, solid-adjacent, open-boundary-adjacent et corner-adjacent. Cette classification a permis de clarifier l'interprétation des cellules vides. Dans le cas Taylor--Green périodique avec trou initial, les cellules vides sont bien classées comme bulk et le resampling les résorbe par advection et correction de support. En revanche, les essais de sillage et de marche ont montré que les zones proches paroi ou proches solide exigent un traitement plus prudent que le bulk.

Le jalon 0306 ajoute enfin un diagnostic d'outliers de vitesse par classe de population et par classe géométrique. Les cellules sont réparties en bins $N=1$, $N=2$, $N=3$, $4\leq N\leq 6$, $7\leq N<N_{\min}$ et $N\geq N_{\min}$. Pour chaque bin, le code suit les cellules dont $\lVert U_c\rVert$ dépasse un seuil utilisateur. Un fichier de pires cellules garde, pour chaque pas diagnostiqué, les coordonnées, $N_c$, $M_c$, $U_x$, $U_y$, $\lVert U_c\rVert$, $K_{\mathrm{rel}}$, $k_BT$ et les indicateurs géométriques. Ce diagnostic a montré que les vitesses aberrantes observées avant 0307 n'étaient pas d'abord des excès thermiques : dans les pires cas, $K_{\mathrm{rel}}=0$ et $k_BT=0$, mais $M_c$ est de l'ordre de $10^{-9}$ et $\lVert U_c\rVert$ peut atteindre plusieurs dizaines. Un rethermostatage relatif n'aurait donc pas corrigé cette pathologie, car il conserve la vitesse moyenne cellulaire.

\subsection{Split-safety 0307--0308 : prévention des cascades de masse faible}

Le diagnostic 0306 a conduit à identifier un mécanisme fondamental : le split local conservatif peut être réappliqué plusieurs fois à des particules déjà allégées. Dans une cellule structurellement pauvre, par exemple solid-adjacent ou proche d'une arête de marche, la règle initiale ``si $0<N_c<N_{\min}$, splitter localement'' peut donc produire une cascade géométrique
\begin{equation}
  m=1 \to \frac12 \to \frac14 \to \frac18 \to \cdots,
\end{equation}
qui finit par créer des particules quasi sans masse. Ces particules ne pèsent presque rien dans les bilans globaux, mais elles peuvent porter une vitesse cellulaire énorme par la formule $U_c=P_c/M_c$ lorsque $M_c$ devient dégénéré.

Le jalon 0307 introduit une sécurité de split activable : choix préférentiel du donneur le plus massif, masse minimale du donneur, masse minimale de la nouvelle particule, masse minimale restante après split, et modes de prudence pour les cellules solid-adjacent. Le mode nominal provisoire conserve le split près des solides mais interdit les fragments sous-massifs :
\begin{equation}
  m_{\mathrm{donor}} \geq 0.5, \qquad
  m_{\mathrm{new}} \geq 0.25, \qquad
  m_{\mathrm{donor}}' \geq 0.25.
\end{equation}
L'objectif n'est pas de corriger a posteriori les particules faibles masses, mais d'empêcher leur création.

Les résultats 0307 sont très discriminants. Sur Von Karman, le mode diagnostic sans prévention observe $88863$ splits, dont $87666$ proviennent de donneurs de masse inférieure à $0.5$ et $86758$ de donneurs de masse inférieure à $0.1$ ; la masse minimale de donneur atteint $2.6\times10^{-13}$. Avec le mode \texttt{safe\_floor}, le nombre de splits chute à $1126$ et aucun split ne provient d'un donneur sous $0.5$. Sur backward step, la masse minimale de nouvelle particule passe d'environ $3.3\times10^{-3}$ à $0.2501$. Ces résultats montrent que la très forte activité du guard avant 0307 était en partie une activité pathologique de resplit de fragments, non un vrai contrôle utile du support.

\begin{table}[htbp]
\centering
\footnotesize
\begin{tabular}{p{0.22\linewidth}p{0.20\linewidth}rrrr}
\toprule
Cas & Mode & splits & donneurs $<0.5$ & donneurs $<0.1$ & $m_{\min}^{\mathrm{new}}$ \\
\midrule
Backward step & diagnostic & 3320 & 1605 & 433 & $3.3\times10^{-3}$ \\
Backward step & safe floor & 2008 & 0 & 0 & $2.50\times10^{-1}$ \\
Von Karman & diagnostic & 88863 & 87666 & 86758 & $2.3\times10^{-14}$ \\
Von Karman & safe floor & 1126 & 0 & 0 & $2.50\times10^{-1}$ \\
\bottomrule
\end{tabular}
\caption{Effet du split-safety 0307 sur les cascades de masse faible. Les chiffres sont ceux des runs de validation 0307 ; ils montrent que le mode \texttt{safe\_floor} coupe la cascade de splits sans désactiver globalement le resampling.}
\end{table}

Le jalon 0308 consolide ce comportement comme mode nominal des scripts de resampling CUDA. La validation stricte 0308, avec complétion contrôlée jusqu'à $6000$ pas, donne quatre runs PASS : backward step classic, backward step resampling, Von Karman classic et Von Karman resampling. Les outliers de vitesse disparaissent : \texttt{sumHighUAllBins=0} dans tous les cas, et les pires cellules ont désormais des masses normales. Par exemple, pour backward step resampling, la pire cellule a $N=20$, $M\simeq19.99$ et $\lVert U\rVert\simeq0.67$ ; pour Von Karman resampling, elle a $N=14$, $M\simeq13.25$ et $\lVert U\rVert\simeq0.49$. La cause principale des vitesses folles était donc la dégénérescence de masse par split répété, non le thermostat.

\subsection{Scripts de visualisation 0309 et diagnostic particulaire MATLAB}

Le jalon 0309 ajoute des scripts de démonstration visuelle indépendants pour Poiseuille, backward step et Von Karman. Chaque cas produit deux répertoires distincts : une référence SRC classic et un run \texttt{resampling\_split\_safe}. Les scripts activent par défaut le mode 0307 consolidé et conservent le survey de support afin de permettre une comparaison quantitative et visuelle.

Les outils MATLAB de visualisation ont également été corrigés pour éviter une confusion importante. Dans les états \texttt{.smpcd}, les slots \texttt{Inactive} gardent leurs positions ; il est donc possible qu'une zone ait $N=0$ en particules fluides tout en affichant des points si le visualiseur trace tous les slots. Les fonctions de lecture/affichage distinguent désormais les rôles \texttt{Fluid}, \texttt{Inactive} et \texttt{Latent}, et peuvent colorer les particules par rôle, masse ou vitesse. Des filtres de seuils sur la masse et sur la vitesse permettent de suivre les particules suspectes sans surcharger la figure. Ces outils sont devenus essentiels pour vérifier que le mode split-safe supprime les particules quasi sans masse dans les sillages longs.

\subsection{Coût des slots inactifs, fast path et dumps compacts 0310--0314}

Les runs longs avec hard inlet et resampling ont montré qu'un grand réservoir \texttt{Inactive} est nécessaire pour éviter l'épuisement des slots lorsque l'inlet injecte un flux net. En revanche, il ne faut pas que les kernels physiques ou les dumps paient systématiquement le coût de la capacité totale. L'audit 0310 a d'abord révélé un défaut de script : certains scripts de démonstration forçaient accidentellement un réservoir de cinq millions de slots inactifs. Cette valeur a été supprimée en 0311.

Le jalon 0312 introduit un audit indépendant des scripts de démonstration. Il génère directement l'état initial et le fichier \texttt{params.kv}, puis lance le binaire sans dépendre des scripts modifiés pour la visualisation. Sur backward step $96\times48$, $400$ pas et \texttt{GUARD\_EVERY=20}, l'audit valide que \texttt{inactiveLogged} suit bien \texttt{inactiveSlotsRequested}. Le temps augmente avec la capacité inactive : $3.06$ s pour $8000$ slots, $3.45$ s pour $50000$ slots et $3.80$ s pour $100000$ slots. Le coût est donc réel, même s'il reste modéré dans cette configuration.

Le jalon 0313 ajoute un fast path de pool inactif en queue de tableau. Lorsqu'une grande réserve inactive est préallouée en fin d'état particulaire, il est inutile de scanner toute la capacité pour trouver des slots disponibles. Le code scanne d'abord une fenêtre bornée dans la queue inactive ; si elle suffit, l'ancien scan complet est évité, sinon un fallback exact conserve la robustesse. Ce fast path concerne le hard inlet full-face, l'inlet segmenté et le population guard lors des splits. Il ne modifie ni la collision SRC, ni le thermostat, ni la conservation masse/impulsion/énergie du resampling.

Le jalon 0314 traite l'autre source de coût : les summaries et dumps. Deux nouvelles clés \texttt{params.kv} contrôlent le rôle à écrire ou à résumer :
\begin{equation}
  \texttt{summaryRoleFilter} \in \{\texttt{all},\texttt{fluid}\}, \qquad
  \texttt{dumpRoleFilter} \in \{\texttt{all},\texttt{fluid}\}.
\end{equation}
Avec \texttt{fluid}, le code produit des états compacts destinés à la visualisation et au post-traitement : seuls les slots fluides sont téléchargés et écrits. Ce mode réduit le coût des dumps et évite de charger inutilement des centaines de milliers de slots inactifs dans MATLAB. Il n'est pas restart-compatible au sens strict, car le réservoir \texttt{Inactive} n'est plus présent dans le dump ; les dumps complets restent disponibles avec \texttt{dumpRoleFilter=all}.

\subsection{Paramètres additionnels 0304--0314}

Les inventaires CSV accompagnant cette version du rapport ont été mis à jour pour inclure les variables d'environnement et clés \texttt{params.kv} ajoutées après 0303. Les paramètres les plus importants sont listés ci-dessous.

\begin{itemize}[nosep]
  \item \path{MPCD_CUDA_RESAMPLING_ADAPTIVE_FLAG_0304} et \path{..._EVERY} : diagnostic post-SRC des cellules low-N ou vides sur grille physique non shiftée.
  \item \path{MPCD_CUDA_RESAMPLING_GEOMETRY_DIAG_0305_HIGH_U} : seuil de vitesse cellulaire utilisé par la classification géométrique 0305.
  \item \path{MPCD_CUDA_RESAMPLING_OUTLIER_0306_U_THRESHOLD} : seuil d'outlier de vitesse utilisé dans les worst-cells 0306.
  \item \path{MPCD_CUDA_RESAMPLING_SPLIT_SAFETY_0307} : active la prévention des cascades de split ; le mode nominal combine \path{MPCD_CUDA_RESAMPLING_SPLIT_PREFER_MAX_MASS_DONOR_0307=1}, \path{MPCD_CUDA_RESAMPLING_SPLIT_DONOR_MIN_MASS_0307=0.5} et \path{MPCD_CUDA_RESAMPLING_SPLIT_NEW_PARTICLE_MIN_MASS_0307=0.25}.
  \item \path{MPCD_CUDA_RESAMPLING_SOLID_ADJACENT_SPLIT_MODE_0307} : mode de prudence pour les cellules solid-adjacent. Le défaut nominal reste le mode normal, car le plancher de masse suffit à supprimer la pathologie sans désactiver la correction près des solides.
  \item \path{MPCD_CUDA_INACTIVE_TAIL_POOL_0313} : active le fast path de recherche de slots inactifs en queue de tableau ; les paramètres \path{..._MIN_SCAN_0313}, \path{..._MAX_SCAN_0313} et \path{..._SCAN_MULT_0313} règlent la fenêtre de scan.
  \item \path{summaryRoleFilter} et \path{dumpRoleFilter} : clés \texttt{params.kv} 0314 contrôlant le mode complet ou compact des summaries et dumps.
\end{itemize}

\subsection{Position par rapport à Q6 et aux futurs développements}

Le resampling CUDA et Q6 traitent deux problèmes distincts. Le premier modifie la représentation particulaire afin de maintenir un support statistique local ; le second corrige la composante compressive du champ de vitesse barycentrique sans changer la population, les masses ou les types. Depuis les jalons 0490--0491, ces deux briques disposent d'une extension multi-espèces et peuvent être activées indépendamment. La section suivante remplace donc l'ancien statut, dans lequel Q6 multi-espèces était encore présenté comme un chantier futur.


\section{Traitement multi-espèces}

Les cycles 0490 et 0491 ont introduit deux extensions orthogonales. Le cycle 0490a--0490p rend le resampling conservatif par espèce et résident CUDA dans son domaine de qualification. Le cycle 0491a--0491h conserve une projection Q6 barycentrique unique, puis distribue sa correction finale entre les espèces selon des poids conservatifs dépendant du 
\texttt{type}. Cette séparation préserve les quatre chemins de simulation :

\begin{table}[htbp]
\centering
\small
\caption{Chemins multi-espèces conservés par l'architecture 0490--0491.}
\begin{tabular}{@{}p{0.24\textwidth}ccp{0.47\textwidth}@{}}
\toprule
Chemin & Resampling & Q6 & Contrat principal \\
\midrule
SRC & non & non & Le \texttt{type} est transporté sans conversion ; la collision peut coupler les espèces présentes dans une cellule. \\
SRC+resampling & oui & non & Les opérations de population et la fermeture de masse sont conservatrices par espèce. \\
SRC+Q6 & non & oui & La projection reste barycentrique ; l'application particulaire peut être commune ou pondérée par espèce. \\
SRC+resampling+Q6 & oui & oui & Les deux briques partagent l'état CUDA, avec Q6 avant thermostat et resampling, et avec invalidation explicite des dépôts. \\
\bottomrule
\end{tabular}
\end{table}

Le multi-espèces considéré ici est un modèle de transport et de couplage cinématique. Il n'introduit pas encore une équation d'état propre à chaque phase, une diffusion Maxwell--Stefan explicite, une réaction chimique, une tension superficielle ou une règle de collision dépendant de la paire d'espèces. Ces compléments pourront réutiliser les dépôts et les diagnostics décrits ci-dessous, mais ils ne doivent pas être confondus avec le contrat actuellement validé.

\subsection{État commun, registre d'espèces et dépôts cellulaires}

\subsubsection{Identité physique et état algorithmique}

Le champ particulaire \texttt{role} décrit l'état algorithmique d'un slot : \texttt{Fluid} pour une particule active, \texttt{Inactive} pour un slot libre du pool et \texttt{Latent} pour une particule préallouée mais encore hors dynamique. Le champ \texttt{type} décrit au contraire l'identité physique transportée. Une opération de resampling peut modifier le \texttt{role}, mais ni le resampling, ni Q6, ni le thermostat ne doivent convertir implicitement une espèce en une autre.

Le registre est activé par
\begin{equation}
  \texttt{speciesRegistryEnable=true},
  \qquad \texttt{speciesCount}=N_s,
\end{equation}
et par des déclarations de la forme
\begin{equation}
\begin{aligned}
  \texttt{speciesK} ={}& \texttt{type name phaseFamily}\\
  &\texttt{q6StrengthDeclared massClosureStrengthDeclared}\\
  &[\texttt{referenceCellMassDeclared}].
\end{aligned}
\end{equation}
La famille vaut actuellement \texttt{liquid} ou \texttt{gas}. Le coefficient \texttt{q6StrengthDeclared} intervient dans la distribution Q6 décrite plus loin ; \texttt{massClosureStrengthDeclared} et la masse cellulaire de référence interviennent dans la fermeture du resampling. La clé \texttt{speciesRequireRegisteredTypes} impose que tout slot physique actif ou latent appartienne au registre. Les slots \texttt{Inactive}, qui ne portent pas de masse fluide active, sont exclus de cette exigence.

\subsubsection{Dépôt cellule--espèce partagé}

Pour une cellule physique non décalée $c$ et une espèce enregistrée $s$, le dépôt reconstruit
\begin{align}
  N_{c,s} &= \sum_{i\in c}\mathbf{1}_{\tau_i=s},\\
  M_{c,s} &= \sum_{i\in c}m_i\mathbf{1}_{\tau_i=s},\\
  \bm{P}_{c,s} &= \sum_{i\in c}m_i\bm{v}_i\mathbf{1}_{\tau_i=s}.
\end{align}
Les grandeurs barycentriques sont
\begin{equation}
  M_c=\sum_sM_{c,s},\qquad
  \bm{P}_c=\sum_s\bm{P}_{c,s},\qquad
  \bm{u}_c=\frac{\bm{P}_c}{M_c},
\end{equation}
et les fractions massiques cellulaires sont
\begin{equation}
  Y_{c,s}=\frac{M_{c,s}}{M_c},
  \qquad \sum_sY_{c,s}=1
\end{equation}
dans toute cellule humide. Pour la politique de resampling, on utilise aussi l'occupation normalisée
\begin{equation}
  \theta_{c,s}=\frac{M_{c,s}}{M_s^{\mathrm{ref}}},
\end{equation}
notation distincte du poids Q6 $w_{c,s}$ introduit plus loin.

Le dépôt CPU déterministe 0490b constitue la référence. Le workspace CUDA 0490h contient notamment les types enregistrés, les coefficients Q6, les masses de référence, les familles de phase, l'activation du resampling par espèce, puis les tableaux denses $N_{c,s}$, $M_{c,s}$, $P_{x,c,s}$, $P_{y,c,s}$, $Y_{c,s}$ et $\theta_{c,s}$. Les allocations sont conservées et réutilisées ; les tableaux denses ne sont téléchargés vers l'hôte que dans les smokes d'équivalence. Cette brique, initialement créée pour le resampling, devient avec 0491 un service partagé accessible même lorsque \texttt{resamplingEnable=false}.

\subsection{Resampling multi-espèces : formulation et implémentation 0490}

\subsubsection{Invariants des opérations discrètes}

Le resampling modifie le nombre de représentants et éventuellement leur masse, mais doit conserver l'identité physique. Ses invariants locaux et globaux sont :
\begin{equation}
  \Delta M_s=0,
  \qquad \Delta\bm{P}_s=\bm{0}
  \qquad \text{pour chaque espèce }s
\end{equation}
dans un domaine fermé, à l'exception des flux explicitement imposés par des ouvertures.

Un split hérite du type du parent. Dans sa forme élémentaire,
\begin{equation}
  (m,\bm{v},s)
  \longrightarrow
  \left(\frac{m}{2},\bm{v},s\right)
  +\left(\frac{m}{2},\bm{v},s\right),
\end{equation}
ce qui conserve exactement masse, quantité de mouvement et énergie cinétique du parent avant déplacement géométrique des filles. Un merge n'est autorisé qu'entre deux particules de même type ; la vitesse du représentant fusionné est choisie barycentriquement,
\begin{equation}
  m'=m_1+m_2,
  \qquad
  \bm{v}'=\frac{m_1\bm{v}_1+m_2\bm{v}_2}{m_1+m_2},
\end{equation}
puis la restauration d'énergie relative existante est appliquée lorsque ce chemin est activé.

Le garde de population conserve la bande globale $N_{\min}/N_{\mathrm{target}}/N_{\max}$, mais répartit la cible entre les espèces présentes. Une méthode déterministe du plus fort reste attribue les représentants, avec au moins un porteur par espèce présente lorsque $N_{\mathrm{target}}$ le permet. Le split est dirigé vers l'espèce la plus déficitaire ; le merge vers une espèce excédentaire possédant au moins deux représentants. Les opérations entre types différents sont interdites, même lorsque leur fusion serait numériquement commode.

\subsubsection{Refill et transferts conservatifs par type}

Le remplissage d'une cellule temporairement vide utilise une mémoire résidente $M^{\mathrm{last}}_{c,s}$. La population recréée est distribuée selon la dernière composition admissible, puis une correction globale par espèce restaure masse et impulsion. Le refill est refusé lorsque la population cible est inférieure au nombre d'espèces à représenter ou lorsqu'une espèce disparue ne possède plus aucun porteur actif permettant une compensation conservative exacte.

Les transferts donneur--receveur sont également planifiés par type. Pour une espèce $s$, une cellule receveuse $r$ ne peut consommer que l'excès disponible chez un donneur $d$ de la même espèce. Le plan contient donc explicitement
\begin{equation}
  (r,d,s,\Delta M_{r\leftarrow d,s}),
\end{equation}
et le matérialiseur rejette tout candidat dont le \texttt{type} ne correspond pas. Une cellule liquide pure ne peut pas être alimentée par un donneur gazeux sous prétexte qu'il est plus proche ou plus riche.

\subsubsection{Fermeture de masse dépendante de la composition}

Pour une cellule non vide, on définit
\begin{equation}
  W_c=\sum_s\frac{M_{c,s}}{M_s^{\mathrm{ref}}},
  \qquad
  M_c^{\star}=\frac{M_c}{W_c}.
\end{equation}
Le coefficient de fermeture est une moyenne d'occupation des coefficients déclarés :
\begin{equation}
  \beta_c=\operatorname{clamp}\!\left[
  \beta_{\mathrm{global}}
  \frac{\sum_s\theta_{c,s}\beta_s}
       {\sum_s\theta_{c,s}},0,1\right],
\end{equation}
et la cible cellulaire devient
\begin{equation}
  M_c^{\mathrm{eff}}
  =M_c+\beta_c\left(M_c^{\star}-M_c\right).
\end{equation}
Une espèce liquide avec $\beta_s\simeq1$ participe fortement à la fermeture ; une espèce gazeuse avec $\beta_s\simeq0$ reste plus libre ; une cellule mixte reçoit un comportement intermédiaire.

La première réalisation 0490d appliquait un facteur commun dans chaque cellule et préservait ainsi exactement la composition locale. Une succession de facteurs différents dans des cellules de compositions variables peut toutefois conserver la masse totale tout en faisant dériver les masses globales par espèce. Le chemin résident ultérieur utilise donc une matrice d'échelles cellule--espèce, équilibrée alternativement sur les contraintes cellulaires et sur les sommes globales par espèce. La dernière normalisation impose $M_s$ comme invariant dur ; un décalage barycentrique commun restaure ensuite la vitesse cellulaire lorsque les facteurs propres aux espèces ont modifié le moment local.

\subsubsection{Chaîne CUDA résidente 0490h--0490p}

Les jalons d'implémentation sont :
\begin{itemize}[nosep]
  \item \textbf{0490a--0490b} : registre d'espèces et dépôt CPU cellule--espèce de référence ;
  \item \textbf{0490c--0490d} : split/merge conservatifs par type et fermeture de masse dépendante de la composition ;
  \item \textbf{0490h} : dépôt $N_{c,s}$, $M_{c,s}$ et $\bm{P}_{c,s}$ résident CUDA ;
  \item \textbf{0490i} : fermeture de masse sur l'état CUDA partagé ;
  \item \textbf{0490j} : sélection d'espèce du garde de population ;
  \item \textbf{0490k} : plan donneur--receveur CUDA ordonné par receveur, espèce puis donneur compatible ;
  \item \textbf{0490l} : validation stricte sans plan ni vecteur d'opérations piloté par le CPU ;
  \item \textbf{0490m} : application directe du plan device, matérialisation et mutation de l'état partagé ;
  \item \textbf{0490n} : dépôts post-opération et pool \texttt{Fluid}/\texttt{Latent}/\texttt{Inactive} résidents ;
  \item \textbf{0490p} : politique wet/poor/rich/target calculée et consommée sur device.
\end{itemize}

La correction 0490n-fix2 distingue les sous-remplissages d'une entrée individuelle, possibles à cause de l'indivisibilité des particules, des sous-remplissages agrégés par couple $(d,s)$. Les premiers sont diagnostiques ; les seconds restent fatals, car un excès provenant d'une autre espèce ou d'un autre donneur ne doit pas masquer un vrai déficit.

Les paramètres structurants sont regroupés par fonction :
\begin{itemize}[nosep]
  \item fermeture : \path{speciesResamplingMassClosureEnable} et son backend CUDA ;
  \item population : \path{speciesResamplingPopulationGuardEnable} et son backend CUDA ;
  \item transferts : \path{speciesResamplingTransferEnable} et son backend CUDA ;
  \item production résidente : fast path, dépôts, pool et maintenance stricte.
\end{itemize}
Les diagnostics principaux sont :
\begin{itemize}[nosep]
  \item \path{species_runtime_0490a.csv} et \path{species_cell_runtime_0490b.csv} ;
  \item \path{species_cell_cuda_equivalence_0490h.csv} ;
  \item \path{cuda_species_mass_closure_0490i.csv} ;
  \item \path{cuda_species_transfer_plan_0490k.csv} ;
  \item \path{cuda_species_resident_fast_path_0490m.csv} ;
  \item \path{cuda_species_resident_maintenance_0490n.csv}.
\end{itemize}

\subsubsection{Statut physique et recommandation d'usage}

Le cycle 0490 établit que le resampling peut être rendu conservatif par espèce et exécuté sans retour complet vers le CPU. Il ne démontre pas qu'il est nécessaire ni bénéfique dans toutes les simulations. Les campagnes ultérieures montrent un coût important lorsque de nombreux couples cellule--espèce sont réparés et une sensibilité possible de l'autodiffusion à la fragmentation et aux transferts, alors que les grandeurs bulk de type Taylor--Green, température et composition restent généralement proches de la référence.

Le resampling multi-espèces ne doit donc pas être le réglage nominal du fluide. Il doit rester une correction \emph{opt-in}, déclenchée pour un défaut de support identifié et comparée à un cas sans resampling. Les critères de décision doivent inclure le nombre de cellules réellement réparées, le coût par pas, la distribution des masses particulaires, la conservation par espèce, l'autodiffusion et l'effet sur les observables physiques ciblées. Les cellules vides, les cellules de coupure géométrique et les ouvertures doivent en outre être traitées selon leur origine physique ; une paroi ou une injection ne peut pas être remplacée par une création locale de masse déguisée.

\subsection{Projection Q6 multi-espèces : formulation 0491}

\subsubsection{Une seule contrainte d'incompressibilité barycentrique}

Pour un mélange de densités $\rho_s$ et de vitesses $\bm{u}_s$, la vitesse barycentrique est
\begin{equation}
  \bm{u}=\frac{\sum_s\rho_s\bm{u}_s}{\rho},
  \qquad \rho=\sum_s\rho_s.
\end{equation}
Imposer séparément $\nabla\cdot\bm{u}_s=0$ à chaque espèce introduirait plusieurs contraintes elliptiques sans disposer, dans le modèle courant, de pressions indépendantes correspondantes. Une telle formulation surcontraindrait les vitesses relatives, pourrait supprimer des flux diffusifs admissibles, deviendrait singulière dans les cellules pures ou traces et multiplierait le coût du solveur.

Le contrat 0491 est donc :
\begin{quote}
\textbf{une seule projection Q6 du mouvement barycentrique, suivie d'une distribution par espèce de la correction particulaire finale.}
\end{quote}
Le dépôt barycentrique, le solveur elliptique, les opérateurs divergence/gradient, les conditions aux limites, la correction globale de quantité de mouvement et la mesure de divergence restent ceux du Q6 résident de référence. Seul le kernel qui applique la correction finale aux particules devient sensible au \texttt{type}.

\subsubsection{Poids d'espèce et identité barycentrique}

Le Q6 existant produit, pour chaque cellule humide, une correction finale
\begin{equation}
  \Delta\bm{u}^{Q6}_c,
  \qquad
  \bm{u}'_c=\bm{u}_c+\Delta\bm{u}^{Q6}_c.
\end{equation}
Chaque espèce déclare un coefficient
\begin{equation}
  \alpha_s=\texttt{q6StrengthDeclared}_s\ge0,
\end{equation}
et un paramètre global $\eta\in[0,1]$, fourni par \texttt{speciesQ6Sensitivity}, règle la sensibilité effective aux différences entre espèces. Dans une cellule humide,
\begin{equation}
  \bar\alpha_c=\sum_sY_{c,s}\alpha_s,
\end{equation}
puis
\begin{equation}
  w_{c,s}=(1-\eta)+\eta\frac{\alpha_s}{\bar\alpha_c},
  \qquad
  \Delta\bm{u}_{c,s}=w_{c,s}\Delta\bm{u}^{Q6}_c.
\end{equation}
La propriété centrale est
\begin{align}
  \sum_sY_{c,s}w_{c,s}
  &=(1-\eta)\sum_sY_{c,s}
  +\frac{\eta}{\bar\alpha_c}\sum_sY_{c,s}\alpha_s\\
  &=1,
\end{align}
d'où l'identité de quantité de mouvement
\begin{equation}
  \sum_sM_{c,s}\Delta\bm{u}_{c,s}
  =M_c\Delta\bm{u}^{Q6}_c.
\end{equation}
La recomposition barycentrique après application est donc exactement celle du Q6 commun, même lorsque les vitesses par espèce diffèrent.

Pour chaque particule active $i$ de type $s(i)$ dans la cellule $c(i)$,
\begin{equation}
  \bm{v}'_i
  =\bm{v}_i+w_{c(i),s(i)}\Delta\bm{u}^{Q6}_{c(i)}.
\end{equation}
La variation de vitesse relative entre deux espèces est
\begin{equation}
  \Delta(\bm{u}_{c,s}-\bm{u}_{c,r})
  =(w_{c,s}-w_{c,r})\Delta\bm{u}^{Q6}_c.
\end{equation}
Elle est nulle en mode commun et constitue l'effet nouveau du mode pondéré.

\subsubsection{Modes de référence et cas limites}

Les cas de non-régression sont intégrés à la formule :
\begin{itemize}[nosep]
  \item \texttt{speciesQ6Mode=common} ou $\eta=0$ donne $w_{c,s}=1$ ;
  \item une seule espèce donne toujours $w=1$ ;
  \item une cellule pure donne $w=1$, quelle que soit la valeur de $\alpha_s$ ;
  \item des coefficients $\alpha_s$ égaux donnent $w=1$ pour toutes les espèces ;
  \item une espèce absente, $M_{c,s}=0$, ne reçoit aucune correction particulaire.
\end{itemize}

Lorsque $\bar\alpha_c<\varepsilon$, le mode permissif revient à $w=1$ avec compteur ; le mode strict déclenche une erreur fatale. Un type non enregistré est fatal avec \texttt{speciesRequireRegisteredTypes=true}. Une espèce trace peut recevoir un poids élevé si son $\alpha_s$ est grand ; les poids minimal, maximal et RMS doivent donc être suivis. Un clamp indépendant est interdit, car il casserait $\sum_sY_{c,s}w_{c,s}=1$. Tout limiteur futur devra redistribuer les poids par une méthode conservative de type active-set.

Q6 ne modifie ni $m_i$, ni $\bm{x}_i$, ni \texttt{role}, ni \texttt{type}, ni la population. Le mode pondéré peut en revanche redistribuer l'énergie cinétique entre espèces. Pour une cellule,
\begin{equation}
  \Delta K_c
  =\sum_s\bm{P}_{c,s}\cdot
     \left(w_{c,s}\Delta\bm{u}^{Q6}_c\right)
  +\frac{1}{2}\sum_sM_{c,s}w_{c,s}^2
     \left\|\Delta\bm{u}^{Q6}_c\right\|^2.
\end{equation}
Cette variation n'est pas un défaut de conservation barycentrique, mais elle doit être mesurée par espèce. Le thermostat cell-relative est appliqué après Q6 ; il doit préserver le mouvement barycentrique corrigé sans annuler silencieusement l'effet relatif recherché.

\subsection{Implémentation CUDA résidente de Q6 multi-espèces}

\subsubsection{Réutilisation du dépôt 0490h}

Le module 0491 consomme le même état CUDA partagé que Q6 résident et le resampling :
\begin{equation}
  (x,y,v_x,v_y,m,\texttt{role},\texttt{type}).
\end{equation}
Le dépôt 0490h fournit $M_{c,s}$ et $Y_{c,s}$ avec ou sans resampling. Les coefficients $\alpha_s$ et la table déterministe \texttt{type}$\rightarrow$\texttt{speciesIndex} sont chargés une fois. Un workspace persistant contient $\bar\alpha_c$, les poids $w_{c,s}$ et un résumé diagnostique compact. Les poids et la correction Q6 sont consommés directement sur device ; aucun tableau cellule--espèce complet, aucune table de poids et aucun état particulaire complet ne doivent être téléchargés pour décider de l'application.

Le calcul reste séquencé ainsi :
\begin{equation}
\begin{aligned}
  \text{streaming}
  &\rightarrow \text{collision SRC}\\
  &\rightarrow \text{Q6 barycentrique et distribution par type}\\
  &\rightarrow \text{thermostat/mean-flow}\\
  &\rightarrow \text{resampling éventuel}.
\end{aligned}
\end{equation}
Le dépôt précédant Q6 reste valide pour $N_{c,s}$, $M_{c,s}$ et $Y_{c,s}$ après l'application, mais sa quantité de mouvement devient obsolète. Le thermostat modifie encore $\bm{P}_{c,s}$. Un split, un merge, un refill, un transfert ou une fermeture de masse invalide tout ou partie du dépôt. Le chemin combiné reconstruit donc un dépôt autoritaire au point post-thermostat attendu par le resampling plutôt que de réutiliser aveuglément un snapshot périmé.

\subsubsection{Jalons 0491a--0491h}

L'extension a été construite par portes successives :
\begin{itemize}[nosep]
  \item \textbf{0491a} : contrat mathématique, paramètres et référence CPU déterministe ;
  \item \textbf{0491b} : dépôt partagé et calcul shadow CUDA des $\bar\alpha_c$, $w_{c,s}$ et résidus ;
  \item \textbf{0491c} : application CUDA opt-in dans le chemin SRC+Q6 ;
  \item \textbf{0491d} : matrice explicite des quatre chemins et ordre Q6--thermostat--resampling ;
  \item \textbf{0491e} : audit résident strict, sans tableau cellule--espèce hôte, sans H2D des poids, sans fallback CPU ni téléchargement complet ;
  \item \textbf{0491f} : interaction avec thermostat, mean-flow et bilan énergétique par espèce ;
  \item \textbf{0491g} : intégration aux familles de conditions limites et géométries déjà supportées par le Q6 barycentrique ;
  \item \textbf{0491h} : qualification matricielle, interface persistante, espèce trace, chemin combiné et mesure de performance.
\end{itemize}

Le point important est que 0491 ne duplique pas le solveur elliptique. Le kernel sensible aux espèces remplace uniquement l'application commune de $\Delta\bm{u}^{Q6}_c$. Cette localisation limite le risque de régression sur la divergence, les masques de faces, Darcy/$\chi$, les ouvertures et la correction globale de moment. Les intégrations ultérieures 0431 et 0493 ont étendu le routage résident aux familles périodique, murs, entrée/sortie pleine face ou segmentée, solides pris en charge et Darcy/$\chi$, sans changer la formule de pondération Q6.

\subsubsection{Paramètres et diagnostics}

Les paramètres utilisateur principaux sont :
\begin{itemize}[nosep]
  \item activation : \path{speciesQ6Enable}, \path{speciesQ6Mode}, \path{speciesQ6Sensitivity} ;
  \item backend : \path{speciesQ6CudaEnable}, \path{speciesQ6CudaResidentStrict} ;
  \item configuration CUDA : \path{speciesQ6ThreadsPerBlock} ;
  \item robustesse : \path{speciesQ6AlphaEpsilon}, \path{speciesQ6FallbackMode} ;
  \item comparaison : \path{speciesQ6ComparisonTolerance} ;
  \item diagnostics : \path{speciesQ6DiagnosticsEnable}, \path{speciesQ6DiagnosticsFilename}.
\end{itemize}

Trois forces doivent être distinguées : \texttt{q6ProjectionStrength} règle l'amplitude barycentrique globale ; \texttt{q6StrengthDeclared} règle la participation relative d'une espèce ; \texttt{speciesQ6Sensitivity} détermine à quel degré cette différence est appliquée. L'activation de \texttt{speciesQ6Enable} n'active pas Q6 lui-même et ne doit pas activer implicitement le resampling.

Pour chaque cellule humide, les résidus durs sont
\begin{align}
  r^{P_x}_c&=\sum_sM_{c,s}\Delta u_{x,c,s}
             -M_c\Delta u^{Q6}_{x,c},\\
  r^{P_y}_c&=\sum_sM_{c,s}\Delta u_{y,c,s}
             -M_c\Delta u^{Q6}_{y,c}.
\end{align}
Le fichier \path{cuda_species_q6_0491.csv} enregistre notamment :
\begin{itemize}[nosep]
  \item le mode, $\eta$ et le nombre d'espèces ;
  \item les cellules pures, mixtes, vides ou en repli ;
  \item les extrema et le RMS des poids, ainsi que les types non enregistrés ;
  \item les résidus barycentriques, les variations de moment et d'énergie par espèce et le RMS des vitesses relatives ;
  \item les erreurs CPU/CUDA, les octets D2H/H2D, les indicateurs résident/fallback et les temps par sous-étape.
\end{itemize}

Les invariants attendus sont : masse et population de chaque espèce inchangées par Q6, aucun changement de \texttt{type} ou de \texttt{role}, correction barycentrique identique au Q6 commun et divergence après projection identique à la tolérance configurée. Trois équivalences sont suivies séparément : mono-espèce, coefficients égaux et recomposition barycentrique avec coefficients différents.

\subsection{Validation, limites et lecture actuelle}

La validation logicielle 0491h-fix1d couvre la matrice des quatre chemins, l'équivalence Q6 historique/common, le mode weighted, une interface persistante, une espèce trace de fraction $10^{-7}$ et un run fermé combiné. La campagne courte se termine par
\begin{equation}
  \texttt{PASS matrix=4, closed\_steps=100,}
  \quad \texttt{interface\_retention=1}.
\end{equation}
Elle a également identifié une incompatibilité entre l'ancien reconditionnement de masse 0296 et la fermeture multi-espèces 0490d/0490i ; ce reconditionnement est explicitement supprimé dans ce chemin afin d'éviter deux autorités concurrentes sur les masses.

Cette qualification courte ne vaut pas qualification physique longue universelle. Une campagne 0491h à 10000 pas a conservé environ $99.1\%$ de l'interface selon son critère, mais n'a pas satisfait l'ensemble de la matrice longue. Elle doit donc être lue comme un résultat de robustesse encourageant, non comme une preuve de conservation d'une interface immiscible : aucune tension superficielle ni fermeture interfaciale n'est encore présente.

Les audits bulk 0493k donnent une image plus favorable pour Q6 lui-même. Sur le mélange binaire, les accords de viscosité Taylor--Green, de courbe TG, de température et de composition passent dans les tolérances retenues ; la divergence avec Q6 reste proche de $31\%$ de celle du SRC non projeté. En revanche, l'autodiffusion dans les chemins avec resampling est restée non résolue ou incompatible selon les fenêtres et versions d'analyse. Ce contraste est cohérent avec la nature des briques : Q6 corrige les vitesses sans modifier la représentation, tandis que le resampling fragmente, fusionne ou déplace des porteurs et peut donc agir directement sur les statistiques de diffusion.

L'état actuel conduit aux recommandations suivantes :
\begin{itemize}
  \item conserver le mode \texttt{common} comme référence permanente de non-régression ;
  \item utiliser le mode \texttt{weighted} de manière opt-in, avec registre strict, diagnostics des poids et bilan énergétique par espèce ;
  \item conserver une seule projection barycentrique et ne pas imposer $\nabla\cdot\bm{u}_s=0$ séparément ;
  \item traiter Q6 et resampling comme deux axes indépendants dans tous les runners et toutes les comparaisons ;
  \item ne pas activer le resampling par défaut tant que son bénéfice sur l'observable ciblée ne compense pas son coût et son effet possible sur la diffusion ;
  \item calibrer d'abord le fluide SRC effectif, puis fixer les régimes $Re$, $Ma$, $Sc$ et $\lambda/a$ avant d'interpréter une faible population comme une pathologie algorithmique ;
  \item réserver les développements suivants aux manques physiques identifiés : traînée inter-espèces, diffusion explicite, pression ou compressibilité de phase, tension superficielle, mouillage et échange de quantité de mouvement interfacial.
\end{itemize}

Q6 multi-espèces constitue ainsi la base la plus saine pour les compléments multiphasiques : il conserve le solveur et la contrainte barycentrique déjà validés, tout en exposant un mécanisme contrôlé de participation des espèces. Le resampling multi-espèces reste une infrastructure utile pour les défauts de support avérés, mais son domaine d'emploi doit être défini par une analyse bénéfice--coût et par des diagnostics de transport, et non par une activation systématique.


\section{Addendum CUDA-TOPO : Darcy--Brinkman et champ \texorpdfstring{$\chi$}{chi}}

\subsection{Objectif du chantier topologique}

Le chantier CUDA-TOPO a été ouvert après la stabilisation du chemin SRC classic CUDA résident et du resampling post-SRC. Son objectif immédiat n'était pas encore de fournir une optimisation topologique complète, mais de vérifier si le solveur pouvait accepter un champ géométrique externe, noté $\chi$, et l'utiliser pour imposer une résistance volumique de type Darcy--Brinkman. L'intérêt attendu est double. D'une part, un automate extérieur d'optimisation peut modifier $\chi$ sans modifier le code source du solveur. D'autre part, le même mécanisme doit pouvoir représenter des géométries non analytiques, telles que des coudes, des canaux diagonaux, des obstacles elliptiques ou des profils NACA.

Dans cette étape, $\chi=1$ désigne le fluide libre et $\chi=0$ la zone pénalisée. La convention retenue est donc opposée à une variable de porosité solide : le champ $\chi$ est un indicateur de liberté fluide. La voie validée à ce stade est une voie SRC classic CUDA-VIZ avec pénalisation Darcy--Brinkman. Elle n'active pas Q6, le resampling ni le viriel dans les scripts de validation topologique, mais l'architecture conserve la possibilité de réintroduire ces modules dans un chantier ultérieur.

\subsection{Formulation Darcy--Brinkman utilisée}

La résistance locale est calculée à partir du champ $\chi$ par une interpolation de type Borrvall--Petersson, dans l'esprit des formulations de pénalisation de fluide utilisées en optimisation topologique et des modèles de Brinkman pour milieux poreux \cite{Brinkman1949,BorrvallPetersson2003} :
\begin{equation}
  \alpha(\chi)=\alpha_{\min}+\left(\alpha_{\max}-\alpha_{\min}\right)\frac{q(1-\chi)}{q+\chi}.
\end{equation}
Cette forme donne $\alpha\simeq \alpha_{\min}$ dans le fluide libre et $\alpha\simeq \alpha_{\max}$ dans la zone solide/poreuse. Le paramètre $q$ règle la raideur de l'interpolation près de l'interface. La pénalisation agit ensuite sur la vitesse moyenne cellulaire, sous la forme d'un amortissement exponentiel vers une vitesse solide prescrite $u_s$ :
\begin{equation}
  u^{\star}=u_s+\exp(-\alpha\Delta t)(u-u_s),
\end{equation}
ou, de manière équivalente,
\begin{equation}
  \Delta u=-\lambda(\alpha)(u-u_s),\qquad
  \lambda(\alpha)=1-\exp(-\alpha\Delta t).
\end{equation}
Pour $\Delta t=5\times10^{-4}$, on obtient par exemple $\lambda\simeq0.393$ pour $\alpha=1000$, $\lambda\simeq0.918$ pour $\alpha=5000$, et $\lambda\simeq0.993$ pour $\alpha=10000$. Ces valeurs expliquent pourquoi des pénalisations modérées peuvent annuler correctement la vitesse dans la zone $\chi\simeq0$ sans reproduire pour autant une frontière solide équivalente.

La force diagnostique associée est une force volumique effective de pénalisation. Elle ne doit pas être lue comme une force aérodynamique pariétale exacte : elle mesure l'impulsion retirée au fluide par le frein Darcy dans les cellules pénalisées. Les observables ajoutées dans le chantier sont donc appelées \emph{proxies} de traînée et de portance.

\subsection{Implémentation CUDA-VIZ et contrôles interactifs}

Le premier jalon du chantier ajoute un chemin CUDA-VIZ pour appliquer la pénalisation Darcy--Brinkman dans le pas SRC classic résident. Les champs livevis disponibles incluent notamment $\chi$, $\alpha$ et la puissance dissipée Darcy. Le fichier \texttt{darcy\_cost\_0343.csv} consigne les grandeurs globales associées à la pénalisation : puissance Darcy, vitesse moyenne, fuite solide et diagnostics de coût.

Un second jalon rend le contrôle livevis persistant par le fichier racine \texttt{livevis\_control.kv}. Ce fichier permet de changer pendant l'exécution le champ visualisé, le gain, les bornes de clipping, le nombre de passes de lissage et la colormap. Ce contrôle s'est révélé important pour distinguer les effets physiques de pénalisation des simples effets d'échelle de couleur.

Le troisième jalon transforme $\chi$ en véritable entrée du solveur. Deux modes coexistent : un mode analytique pour les validations rapides et un mode \texttt{file} où le champ $\chi$ est lu une fois côté CPU, envoyé une fois sur GPU, puis utilisé en résident. Le pré-calcul de $\chi$, $\alpha$ et $\lambda$ évite de recalculer l'interpolation à chaque pas. La comparaison cercle analytique/cercle fichier a montré des résultats quasi identiques, ce qui valide le chemin d'entrée externe.

\subsection{Générateurs de géométrie et cas de validation}

Plusieurs générateurs de champ $\chi$ ont été ajoutés progressivement. Les premiers cas ont porté sur des canaux courbes et un canal diagonal, afin de vérifier que le code n'était plus limité aux formes analytiques simples. Les cas \texttt{bend\_pipe} et \texttt{diagonal\_channel} ont validé la lecture de $\chi$ externe, la visualisation directe de la géométrie et l'effet qualitatif de la pénalisation sur le flux.

La famille suivante de générateurs a introduit des obstacles de canal : cylindre, ellipse et profils NACA à quatre chiffres. Le profil NACA0012 a servi de cas de référence car sa symétrie impose une portance nulle à incidence nulle et une portance antisymétrique lorsque l'incidence change de signe. Les scripts de sweep créent les champs $\chi$, lancent les cas Darcy--Brinkman, appliquent les statistiques de fenêtre finale et construisent des polaires proxy.

\subsection{Observables de force et statistiques de fenêtre}

Les observables de benchmark sont volontairement optionnelles et cadencées pour limiter le coût des atomiques GPU. Le fichier \texttt{topo\_benchmark\_0348.csv} contient les grandeurs de force cellulaire projetées sur une direction de traînée et une direction de portance. Le post-traitement \texttt{0348b} extrait ensuite des moyennes, écarts-types, minima, maxima et dernières valeurs sur une fenêtre finale. Cette séparation évite de confondre le transitoire initial et le régime établi.

Les grandeurs principales sont
\begin{itemize}[leftmargin=*]
  \item \texttt{dragProxy}, projection de la force Darcy sur la direction de l'écoulement ;
  \item \texttt{liftProxy}, projection transverse ;
  \item \texttt{darcyPower}, puissance dissipée par la pénalisation ;
  \item \texttt{solidLeakRms}, fuite de vitesse dans la zone pénalisée ;
  \item les rapports \texttt{absLiftOverAbsDrag} et \texttt{liftOverDragProxy}.
\end{itemize}
Ces grandeurs sont cohérentes entre cas, mais elles restent des indicateurs de pénalisation volumique et non des coefficients $C_L$ ou $C_D$ calibrés.

\subsection{Calibration Poiseuille et échelle de Reynolds effective}

Afin de comparer les sweeps NACA à une échelle hydrodynamique, un script de calibration Poiseuille a été ajouté. Pour la grille cible NACA $N_x=600$, $N_y=160$, $\gamma=10$, $\Delta t=5\times10^{-4}$ et $k_BT=0.01$, le fit Poiseuille donne
\begin{equation}
  \nu_{\mathrm{eff}}\simeq 2.328\times 10^{-3}.
\end{equation}
Le cas $k_BT=0.01$ est retenu comme calibration préférée car le fit présente un coefficient $R^2\simeq0.977$, nettement plus propre que le cas $k_BT=0.1$. Pour une corde $c=0.22$ et une vitesse $U_0=1$, le Reynolds effectif vaut alors
\begin{equation}
  Re_{\mathrm{eff}}=\frac{U_0 c}{\nu_{\mathrm{eff}}}\simeq94.5.
\end{equation}
Cette valeur est faible devant les références expérimentales NACA classiques à haut Reynolds. Les comparaisons doivent donc être interprétées comme des comparaisons de forme normalisée, pas comme des validations quantitatives de coefficients aérodynamiques.

\subsection{Résultats NACA : comparaison normalisée aux polaires de référence}

Le sweep NACA0012 calibré à $Re_{\mathrm{eff}}\simeq94.5$ montre une réponse de portance proxy très structurée. La portance est quasi nulle à incidence nulle, change de signe lorsque l'angle d'attaque change de signe et croît presque linéairement sur l'intervalle étudié. Après normalisation par la valeur à $10^\circ$, la forme de la courbe de portance est proche de la courbe de référence NACA/Ladson utilisée comme repère qualitatif.

La traînée proxy est plus problématique. Elle est bien paire et minimale autour de $0^\circ$, mais sa croissance avec $|\alpha|$ reste trop faible. Dans le sweep de référence, la traînée normalisée augmente d'environ $28\%$ entre $0^\circ$ et $14^\circ$, alors que la référence NACA normalisée utilisée dans la comparaison augmente d'un facteur supérieur à deux. La conséquence est une finesse proxy trop favorable aux grands angles d'attaque, car la portance augmente raisonnablement alors que la traînée ne croît pas assez.

Un sweep de sensibilité en $\alpha_{\max}$ et $U_0$ a ensuite montré que l'augmentation de $\alpha_{\max}$ améliore l'imperméabilité locale et réduit la fuite solide, mais ne rend pas la courbe de traînée suffisamment convexe. L'augmentation de $U_0$ accroît davantage la courbure de traînée, mais au prix d'un bruit plus important et sans rétablir une vraie polar d'obstacle solide. Cette observation suggère que la limite ne provient pas uniquement d'une pénalisation trop faible, mais de la nature même du modèle Darcy pur.

\subsection{Comparaison Von Karman : solide immergé et Darcy pur}

Un cas autonome Von Karman a été ajouté pour comparer proprement le cylindre solide défini dans le code et un cylindre représenté uniquement par le champ Darcy. Les scripts ont été réalignés sur le chemin CUDA périodique/cercle validé du script portable : même chemin \texttt{periodic\_circle}, même thermostat non partagé, même angle de rotation effectif et mêmes paramètres de livevis. Après cet alignement, le mode \texttt{immersedSolid} du script autonome reproduit le comportement du script portable.

Le test critique compare ensuite trois modèles : \texttt{immersed}, \texttt{darcy} et, pour diagnostic seulement, \texttt{both}. Le mode \texttt{immersed} active la frontière solide codée et exclut les particules du cylindre dans l'état initial. Le mode \texttt{darcy} désactive la frontière solide et laisse le cylindre rempli de particules ; seule la pénalisation volumique freine la vitesse moyenne dans la zone $\chi\simeq0$.

Les diagnostics de densité de sillage sur $5000$ pas donnent le tableau synthétique suivant :
\begin{center}
\begin{tabular}{lcccc}
\toprule
Modèle & $\langle N_{\mathrm{wake}}/N_{\mathrm{up}}\rangle$ & Dernier $N_{\mathrm{wake}}/N_{\mathrm{up}}$ & $\langle U_{x,\mathrm{wake}}/U_{x,\mathrm{up}}\rangle$ & Dernier $U_{x,\mathrm{wake}}/U_{x,\mathrm{up}}$\\
\midrule
Solide immergé & 0.959 & 0.969 & 0.766 & 0.807\\
Darcy pur & 1.000 & 1.000 & 0.999 & 0.997\\
\bottomrule
\end{tabular}
\end{center}
Le résultat est très instructif. Le solide immergé crée un sillage réel : légère déplétion de population et déficit de vitesse d'environ $20$ à $25\%$ dans la boîte de sillage proche. Le Darcy pur, malgré une vitesse faible dans la zone pénalisée elle-même, ne crée pratiquement ni déficit de population ni déficit de vitesse en aval. Le fluide traverse donc un frein volumique poreux plus qu'il ne contourne une frontière solide.

Cette observation invalide l'hypothèse initiale selon laquelle les mauvais proxies de traînée seraient dus à un appauvrissement excessif du sillage Darcy. C'est au contraire le solide immergé qui crée le sillage plus marqué. La limite du Darcy pur est plutôt son incapacité à imposer une condition de non-pénétration et une production de vorticité pariétale équivalentes à celles d'une paroi solide.



\subsection{Complément 0418--0426 : de Darcy pur vers une frontière pilotée par \texorpdfstring{$\chi$}{chi}}

Les premiers tests ont montré qu'une pénalisation Darcy pure arrête bien la vitesse moyenne dans la zone $\chi\simeq0$, mais ne crée pas le même sillage qu'une frontière solide codée. La suite du chantier a donc consisté à transformer le champ $\chi$ en support géométrique plus riche, sans revenir à une liste de formes analytiques figées. Trois ajouts sont importants.

Le premier est la désactivation initiale des particules dans la zone solide, contrôlée par \texttt{darcyInitialDeactivateBelowChi}. Lorsque ce seuil vaut par exemple $0.05$, les particules placées initialement dans les cellules très solides sont basculées hors du rôle fluide. Le champ $\chi$ agit donc dès l'état initial comme un masque géométrique, et non seulement comme une force volumique appliquée après coup. Ce point est essentiel pour les géométries de type backward step, bend pipe ou obstacle, car il évite de commencer le calcul avec une masse fluide artificielle dans le solide.

Le deuxième ajout est le mode \texttt{mean} de la pénalisation. Le code dépose masse et quantité de mouvement sur la grille, reconstruit la vitesse moyenne cellulaire, puis applique le kick Darcy--Brinkman sur cette vitesse moyenne avant de redistribuer la correction aux particules. Dans ce mode, la force n'est pas une correction particule par particule indépendante : elle est cohérente avec le moment cellulaire local. Le mode \texttt{mean\_outward\_bath} ajoute ensuite un bain orienté par la normale locale. La normale est reconstruite à partir de $\nabla\chi$ ; avec la convention $\chi=0$ solide et $\chi=1$ fluide, elle pointe du solide vers le fluide. Le bain corrige la composante relative des vitesses de manière à favoriser une réémission hors du solide et à limiter les fuites résiduelles au voisinage de l'interface. Dans les validations backward step, ce bain n'est pas apparu comme un coût dominant ; au contraire, il stabilise localement la dynamique et réduit les corrections parasites.

Le troisième ajout est \texttt{chiVP}, c'est-à-dire une contribution de particules virtuelles de collision dérivée de $\chi$. Il ne s'agit pas de créer, streamer ou dumper des particules virtuelles persistantes. Le mécanisme ajoute seulement une masse et une quantité de mouvement effectives dans la collision SRC des cellules d'interface. Si $M$ et $P$ sont la masse et le moment réels déposés dans une cellule et si $M_{\rm vp}$ désigne la masse virtuelle issue de $\chi$, le centre de masse de collision est remplacé par
\begin{equation}
  u_{\rm cm}^{\rm eff} = \frac{P + M_{\rm vp} u_s}{M+M_{\rm vp}}.
\end{equation}
Dans les scripts, \texttt{darcyChiCollisionVpGamma=-1} signifie que l'occupation virtuelle de référence est héritée de la valeur MPCD nominale, \texttt{darcyChiCollisionVpLayers} fixe l'épaisseur de bande, \texttt{darcyChiCollisionVpThreshold} définit le seuil d'interface et \texttt{darcyChiCollisionVpStrength} règle l'intensité relative. Sur le cas Von Karman filtré, la valeur \texttt{0.25} a donné le meilleur compromis global entre proche sillage, sillage aval et zone proche paroi. Cette valeur a donc été retenue pour le backward step Darcy/chiVP.

\subsection{Correction de performance 0426 : fastflags CUDA résidents}

Une ablation backward step $960\times480$ a isolé un problème d'aiguillage important. Les premiers runs Darcy segmentés semblaient environ trois à quatre fois plus lents que le solide rectangle full-face. Les tests successifs ont écarté les hypothèses naturelles : désactiver \texttt{chiVP}, enlever le bain \texttt{outward}, mettre $\alpha=0$ ou désactiver \texttt{darcyBrinkmanEnable} ne ramenait pas le temps vers le chemin solide. L'activation du profil interne a montré que le temps était absorbé par \texttt{src\_collision}, \texttt{force\_stream} et un bloc géométrique générique, et non par les kernels Darcy proprement dits.

La cause était finalement l'absence, dans les scripts Darcy de démonstration, du même groupe de drapeaux CUDA résidents que dans les scripts solides optimisés. Le bloc 0426 active par défaut, entre autres, la rotation sur device, le diagnostic thermostat rapide, le dépôt streaming fusionné, le contrôle paresseux des kernels, la suppression de synchronisations finales ou de setup inutiles, la suppression du téléchargement du workspace et l'évitement du remplissage hôte des cell-id lorsque le chemin résident le permet. Il est encapsulé dans les scripts par \texttt{export\_src\_cuda\_resident\_fastflags\_0426} et peut être désactivé par \texttt{MPCD\_DARCY\_FASTFLAGS\_ENABLE=0} pour diagnostic.

Le résultat de performance est net. Sur $500$ pas du backward step $960\times480$, le cas \texttt{darcyOff} segmenté passe d'environ $109\,\si{\second}$ profilées à environ $28\,\si{\second}$ après activation des fastflags. Le vrai cas physique \texttt{mean\_outward\_bath+chiVP}, avec $\alpha_{\max}=8\times10^5$, donne ensuite environ $28\,\si{\second}$ profilées et $31\,\si{\second}$ de temps mural pour $500$ pas, soit un coût voisin du \texttt{solid\_fullface} profilé. La conclusion est donc forte : le chemin Darcy/chiVP n'est pas intrinsèquement coûteux ; le ralentissement provenait d'une configuration incomplète du backend résident.

\subsection{Conséquence pour les formes solides arbitraires}

Après cette correction, le champ $\chi$ devient une interface pratique pour imposer des formes solides quelconques dans le chemin SRC classic CUDA résident. Le solveur n'a plus besoin d'un \texttt{caseType} ou d'une forme analytique dédiée pour chaque géométrie : un générateur externe, un automate d'optimisation ou un script de prétraitement peut écrire $\chi$, puis le même pipeline applique désactivation initiale, pénalisation moyenne, bain orienté et collision virtuelle d'interface. Cette souplesse est particulièrement importante pour les coudes, marches, obstacles composites et profils issus d'une boucle d'optimisation.

Il reste néanmoins utile de garder la distinction conceptuelle suivante. Le mode Darcy pur représente un milieu poreux ou une résistance volumique. Le mode augmenté \texttt{mean\_outward\_bath+chiVP} fournit une approximation beaucoup plus solid-aware de la frontière, et devient la voie opérationnelle pour les géométries arbitraires pilotées par $\chi$. Une condition de paroi strictement équivalente à un bounce-back analytique reste à qualifier par des métriques de sillage, de production de vorticité et de fuite normale, mais le verrou principal de performance est levé.

\subsection{Statut validé et limite identifiée}

Le point d'étape Darcy--Brinkman doit maintenant être formulé en distinguant clairement trois niveaux :
\begin{itemize}[leftmargin=*]
  \item le chemin CUDA-VIZ applique correctement une pénalisation Darcy--Brinkman issue d'un champ $\chi$ analytique ou externe ;
  \item les champs $\chi$, $\alpha$, vitesse, vorticité et puissance Darcy sont visualisables en temps réel, avec contrôle persistant par \texttt{livevis\_control.kv} ;
  \item les générateurs de géométrie et les sweeps NACA démontrent que $\chi$ peut piloter des formes complexes sans modifier le solveur ;
  \item les observables de force Darcy restent cohérentes comme proxies de pénalisation volumique ;
  \item Darcy pur reste un modèle poreux et ne reproduit pas, à lui seul, une frontière solide impenetrable ni le même sillage qu'un solide immergé codé ;
  \item la chaîne augmentée \texttt{mean\_outward\_bath+chiVP}, combinée aux fastflags 0426, fournit désormais une voie efficace et flexible pour traiter des formes solides quelconques à partir de $\chi$.
\end{itemize}

Cette distinction est importante pour l'optimisation topologique. Le champ $\chi$ est validé comme entrée géométrique externe et comme champ de pénalisation poreuse. Il n'est pas encore validé comme représentation complète d'un solide équivalent. Pour atteindre cet objectif, un futur chantier devra construire une condition solide dérivée de $\chi$ : détection d'interface, normale locale, réflexion ou bounce-back, éventuellement particules virtuelles de paroi dérivées de $\nabla\chi$. On peut alors distinguer deux usages futurs :
\begin{itemize}[leftmargin=*]
  \item \texttt{CYLINDER\_MODEL=darcy}, pour représenter un milieu poreux ou une résistance volumique continue ;
  \item \texttt{CYLINDER\_MODEL=chi\_solid}, à développer, pour représenter un obstacle solide général piloté par $\chi$.
\end{itemize}

Le chantier Darcy--Brinkman est donc validé comme brique topologique de pénalisation et, depuis 0426, comme interface pratique pour généraliser le traitement des formes solides à partir de $\chi$ sans dépendre de formes prédéfinies. La validation fine des conditions pariétales équivalentes -- fuite normale, production de vorticité, longueur de sillage et efforts -- reste à poursuivre, mais l'obstacle principal de performance du chemin Darcy/chiVP a été levé.

\section{Étalonnage du fluide SRC effectif et référence de co-simulation (0493w1)}
\label{sec:etalonnage-src-effectif}

\subsection{Motivation : caractériser le fluide avant de corriger la simulation}

Les développements consacrés aux cellules pauvres, au resampling, aux parois et aux obstacles ont montré qu'une situation apparemment pathologique peut avoir deux origines très différentes. Elle peut provenir d'un défaut numérique réel de support particulaire, mais elle peut aussi être la conséquence normale d'un fluide SRC dont les coefficients de transport et les échelles acoustiques sont mal adaptés au domaine imposé. Une zone de sillage peu repeuplée, une faible diffusion transverse ou une accumulation à proximité d'une ouverture ne suffisent donc pas, à elles seules, à justifier une correction de population.

La première calibration appliquée au cas Darcy segmenté historique a précisément révélé un régime extrême. Pour
\begin{equation}
  a=\frac{1}{128},\qquad
  \Delta t=10^{-3},\qquad
  k_BT=10^{-3},\qquad
  \gamma=20,
\end{equation}
avec une rotation SRC de $90^\circ$, décalage aléatoire de grille et thermostat \texttt{cell\_relative\_rescale} appliqué à chaque pas, les estimations obtenues étaient
\begin{equation}
  \nu\simeq5.09\times10^{-3},\qquad
  D_{\mathrm{self}}\simeq7.64\times10^{-7},\qquad
  Sc\simeq6.66\times10^{3},\qquad
  \frac{\lambda}{a}\simeq5.1\times10^{-3}.
\end{equation}
Pour $U=0.15$ et $L=0.24$, le Reynolds n'était que $Re\simeq7.1$. L'identification acoustique de ce premier essai était invalide du point de vue de la quantité de mouvement ; seules les bornes idéales isotherme et adiabatique 2D pouvaient être utilisées comme proxies, conduisant à $Ma\simeq3.35$--$4.74$. L'échelle combinée $c_sL/\nu=Re/Ma$ n'était alors que de l'ordre de $1.5$--$2.1$. Le cas réunissait donc faible Reynolds, très forte compressibilité à la vitesse imposée, autodiffusion extrêmement faible et libre parcours presque nul devant la cellule. Une partie des besoins de réparation observés dans ce régime doit par conséquent être requalifiée avant d'être érigée en exigence générale du resampling.

La règle méthodologique désormais retenue est la suivante :
\begin{quote}
Toute configuration SRC destinée à un calcul physique doit être caractérisée comme un fluide effectif avant l'activation de Q6/Q9, du resampling ou d'une fermeture quasi incompressible. Les corrections ultérieures doivent être évaluées relativement à cette référence, et non relativement à des proxies théoriques ou à une vitesse imposée arbitrairement.
\end{quote}

\subsection{Configuration qui définit le fluide effectif}

Le fluide n'est pas défini par la seule grille. La signature minimale de la calibration est
\begin{equation}
\Theta_{\mathrm{SRC}}=
\left(
L_x,L_y,N_x,N_y,\gamma,\Delta t,k_BT,m,\alpha,
\mathcal{S}_{\pm},\mathcal{G},\mathcal{T}
\right),
\end{equation}
où $\mathcal{S}_{\pm}$ décrit le choix aléatoire du signe de rotation, $\mathcal{G}$ le décalage aléatoire de grille et $\mathcal{T}$ la politique de thermostat, y compris sa cadence, sa température cible et sa population minimale. La taille locale de cellule
\begin{equation}
  a_x=\frac{L_x}{N_x},\qquad a_y=\frac{L_y}{N_y}
\end{equation}
entre directement dans les coefficients de transport. À paramètres locaux identiques, une augmentation simultanée de $L$ et de $N$ conserve $a$ et constitue un contrôle de taille de boîte ; à domaine fixé, une augmentation de $N$ modifie $a$ et définit un autre fluide SRC.

Le thermostat appartient explicitement à cette signature. Un SRC brut et un SRC thermostaté à chaque pas ne sont pas deux réalisations d'un même fluide : ils peuvent présenter des viscosités, diffusions et réponses acoustiques différentes. Le calibrateur classe donc la configuration comme \texttt{raw\_src} ou \texttt{thermostatted\_src}. Q6 et le resampling sont volontairement désactivés afin de mesurer le fluide sous-jacent sur lequel ces modules agiront ensuite.

\subsection{Runner canonique de cas unique}

Le runner de référence est
\begin{center}
  \texttt{scripts/run\_0493w1\_src\_fluid\_calibrator.sh}.
\end{center}
Il prend ses paramètres par variables d'environnement, produit les états périodiques nécessaires, lance les expériences et appelle l'analyseur \texttt{analyze\_0493w1\_src\_fluid\_calibrator.py}. Son contrat courant est résumé dans le tableau suivant.

\begin{center}
\small
\begin{tabular}{p{0.34\textwidth}p{0.58\textwidth}}
\toprule
Entrées & Rôle physique ou numérique \\
\midrule
\texttt{Lx, Ly, NX, NY} & Domaine périodique de calibration, grille et taille des cellules de collision. \\
\texttt{GAMMA, DT, KBT, PARTICLE\_MASS} & Occupation moyenne, intervalle de collision, énergie thermique et masse particulaire. \\
\texttt{ROTATION\_ANGLE, RANDOM\_ROTATION\_SIGN, GRID\_SHIFT\_ENABLE} & Règle de collision SRC et restauration de l'invariance galiléenne. \\
\texttt{THERMOSTAT\_*} & Activation, mode, cadence, cible $k_BT$ et population minimale du thermostat de production. \\
\texttt{TG\_MODE\_X/Y, TG\_STEPS, TG\_DUMP\_COUNT} & Longueur d'onde, durée et échantillonnage de la mesure de viscosité. \\
\texttt{SOUND\_MODE\_X, SOUND\_REPLICATES, SOUND\_STEPS, SOUND\_DUMP\_COUNT} & Longueur d'onde et ensemble de réalisations pour la mesure acoustique. \\
\texttt{MSD\_GRID\_MODE, MSD\_STEPS, MSD\_DUMP\_COUNT} & Domaine et durée de la mesure d'autodiffusion. \\
\texttt{CHARACTERISTIC\_U, CHARACTERISTIC\_L} & Échelles de la future simulation pour calculer $Re$, $Ma$, $Pe$ et les vitesses compatibles avec des cibles données. \\
\bottomrule
\end{tabular}
\end{center}

Une calibration standard avec \texttt{SOUND\_REPLICATES=4} effectue six simulations SRC indépendantes : un Taylor--Green, quatre modes longitudinaux et un MSD. Le préflight, activé par \texttt{PREFLIGHT\_ONLY=1}, n'exécute aucune simulation et vérifie notamment la géométrie, la résolution des longueurs d'onde, le nombre d'intervalles de régression, le nombre d'échantillons par période acoustique proxy, le volume estimé des dumps et le nombre de \emph{particle-steps}. Les options \texttt{MAX\_DUMP\_GB} et \texttt{ALLOW\_LARGE\_DUMPS} évitent une campagne volumineuse involontaire. \texttt{ANALYZE\_ONLY=1} permet de reconstruire les résultats à partir de dumps existants.

Le runner caractérise un fluide homogène dans une boîte périodique ; il ne reproduit pas l'obstacle ou la paroi de la simulation finale. Deux usages sont distingués :
\begin{itemize}[leftmargin=*]
  \item le mode \emph{cell-equivalent}, qui conserve $a$, $\Delta t$, $k_BT$ et les règles SRC dans une boîte périodique compacte, afin de mesurer à faible coût les propriétés locales du fluide ;
  \item le mode \emph{domaine complet}, qui conserve le domaine et la grille de la simulation cible pour contrôler les effets de taille finie, de statistique et de longueur d'onde. L'option \texttt{MSD\_GRID\_MODE=full} impose alors le domaine complet pour l'autodiffusion.
\end{itemize}

\subsection{Mesure de la viscosité par décroissance Taylor--Green}

Le générateur initialise le mode transverse
\begin{equation}
\begin{aligned}
  u_x(x,y,0)&=A_0\sin(k_xx)\cos(k_yy),\\
  u_y(x,y,0)&=-A_0\cos(k_xx)\sin(k_yy),
\end{aligned}
\qquad
k_x=\frac{2\pi n_x}{L_x},\quad
k_y=\frac{2\pi n_y}{L_y}.
\end{equation}
L'amplitude est reconstruite par projection massique des vitesses particulaires sur cette base. Dans le régime hydrodynamique linéaire,
\begin{equation}
  A(t)=A_0\exp\left[-\nu(k_x^2+k_y^2)t\right],
\end{equation}
d'où
\begin{equation}
  \nu=-\frac{1}{k_x^2+k_y^2}\frac{\mathrm d}{\mathrm dt}\ln|A(t)|.
\end{equation}

L'analyseur construit plusieurs fenêtres de fit, écarte les points situés sous $2.5\,\%$ de l'amplitude initiale afin de ne pas ajuster le plancher stochastique, impose au moins sept points et une décroissance logarithmique suffisante, puis choisit la fenêtre de meilleure linéarité. Il publie la valeur choisie, $R^2$, le nombre de points et la dispersion des valeurs de $\nu$ sur les fenêtres acceptables. Le statut est \texttt{PASS} pour une valeur positive avec $R^2\geq0.98$, \texttt{REVIEW} pour $0.93\leq R^2<0.98$ et \texttt{INVALID} dans les autres cas.

Une réalisation TG unique peut toutefois rester dominée par le bruit lorsque le mode devient faible. Pour une référence précise, la campagne 0493w1 a montré qu'il faut moyenner les amplitudes signées de plusieurs graines à temps identiques, puis effectuer une seule régression sur le signal d'ensemble. La moyenne des viscosités ajustées séparément est moins robuste, car chaque trajectoire peut sélectionner une fenêtre différente après son entrée dans le bruit.

\subsection{Mesure acoustique par régression hydrodynamique cumulative}

L'état acoustique est initialisé sans vitesse cohérente,
\begin{equation}
  \rho(x,0)=\rho_0\left[1+\varepsilon\cos(kx)\right],
  \qquad u_x(x,0)=0,
\end{equation}
avec $\varepsilon=0.08$ dans la campagne fix3. À chaque dump, l'analyseur calcule les modes complexes normalisés de densité et de vitesse longitudinale. Les quatre réalisations thermiques indépendantes sont synchronisées et moyennées avant le fit.

Avec la convention de Fourier $\exp(-ikx)$, le modèle linéaire utilisé est
\begin{equation}
  \frac{\mathrm d\widehat\rho}{\mathrm dt}=-ik\widehat u,
\end{equation}
\begin{equation}
  \frac{\mathrm d\widehat u}{\mathrm dt}
  =-ikc_s^2\widehat\rho-\nu_Lk^2\widehat u.
\end{equation}
Plutôt que de différencier des données bruitées ou d'exiger une période acoustique complète, l'analyseur intègre ces deux équations depuis l'état initial par la règle du trapèze :
\begin{equation}
  \widehat\rho(t)-\widehat\rho(0)
  =-ik\int_0^t\widehat u(\tau)\,\mathrm d\tau,
\end{equation}
\begin{equation}
  \widehat u(t)-\widehat u(0)
  =-ikc_s^2\int_0^t\widehat\rho(\tau)\,\mathrm d\tau
  -\nu_Lk^2\int_0^t\widehat u(\tau)\,\mathrm d\tau.
\end{equation}
Une régression complexe à deux paramètres avec contraintes de non-négativité fournit $c_s^2$ et $\nu_L$. Les résidus relatifs des balances de continuité et de quantité de mouvement, le conditionnement, les contraintes actives et un jackknife par retrait d'une réalisation quantifient la qualité. Une valeur acoustique \texttt{INVALID} ou \texttt{REVIEW} ne doit pas servir à publier un Mach mesuré ; le runner conserve alors seulement les proxies idéaux
\begin{equation}
  c_{s,\mathrm{iso}}=\sqrt{\frac{k_BT}{m}},
  \qquad
  c_{s,\mathrm{ad},2D}=\sqrt{\frac{2k_BT}{m}}.
\end{equation}
Dans les six fluides de la campagne physique, la vitesse mesurée est proche de la branche isotherme, ce qui est cohérent avec le thermostat cellulaire appliqué à chaque pas.

\subsection{Mesure de l'autodiffusion par MSD périodique}

Le troisième état est un fluide homogène à l'équilibre. Un sous-ensemble persistant de particules est suivi entre les dumps. Les incréments périodiques sont déroulés, puis le déplacement du centre de masse de l'échantillon est retiré afin de ne pas confondre diffusion et dérive globale. En deux dimensions,
\begin{equation}
  \left\langle|\Delta\bm r(t)|^2\right\rangle=4D_{\mathrm{self}}t.
\end{equation}
L'analyseur ajuste les portions tardives commençant à $20$, $30$, $40$ ou $50\,\%$ de la durée totale et retient la fenêtre de meilleur $R^2$. Il publie également l'anisotropie $MSD_x/MSD_y$, le plus grand incrément relatif à la boîte et le paramètre non gaussien 2D
\begin{equation}
  \alpha_2=\frac{\langle r^4\rangle}{2\langle r^2\rangle^2}-1.
\end{equation}
Le statut de diffusion est \texttt{PASS} pour $R^2\geq0.995$, \texttt{REVIEW} pour $0.98\leq R^2<0.995$ et \texttt{INVALID} sinon.

\subsection{Grandeurs dérivées et lecture physique}

À partir des trois mesures, le runner calcule notamment
\begin{equation}
  Sc=\frac{\nu}{D_{\mathrm{self}}},
  \qquad
  Re=\frac{UL}{\nu},
  \qquad
  Ma=\frac{U}{c_s},
  \qquad
  Pe_{\mathrm{mass}}=\frac{UL}{D_{\mathrm{self}}}.
\end{equation}
Le libre parcours indiqué est un proxy de déplacement balistique thermique pendant un intervalle de collision,
\begin{equation}
  \lambda_{\mathrm{mean}}=\sqrt{\frac{\pi k_BT}{2m}}\,\Delta t,
  \qquad
  \frac{\lambda}{a}=\frac{\lambda_{\mathrm{mean}}}{\sqrt{a_xa_y}},
\end{equation}
et non un libre parcours moléculaire mesuré. Le runner fournit aussi
\begin{equation}
  \frac{Re}{Ma}=\frac{c_sL}{\nu},
\end{equation}
qui mesure directement la capacité du fluide à atteindre simultanément un Reynolds élevé et un Mach faible, ainsi que la vitesse compatible avec $Ma=0.3$, le Reynolds correspondant, la vitesse imposant $Re=50$ et le Mach résultant.

\subsection{Pré-balayage physique de six configurations}

Le pré-balayage 0493w1 a conservé $\gamma=20$, $m=1$, une rotation de $90^\circ$, le signe aléatoire, le décalage de grille et le thermostat \texttt{cell\_relative\_rescale} à chaque pas. Il a varié $a$, $\Delta t$ et $k_BT$. Les valeurs du tableau suivant proviennent d'une calibration complète par cas et ont toutes obtenu les trois statuts \texttt{PASS}. Elles constituent un balayage de sélection ; la viscosité du candidat final est ensuite affinée par un ensemble multi-graines.

\begin{center}
\scriptsize
\begin{tabular}{lrrrrrrrr}
\toprule
Cas & $a$ & $\Delta t$ & $k_BT$ & $\lambda/a$ & $\nu$ & $c_s$ & $D_{\mathrm{self}}$ & $Sc$ \\
\midrule
\texttt{a192\_dt002\_k030} & $1/192$ & .002 & .030 & .083 & $1.025\,10^{-3}$ & .17054 & $4.226\,10^{-5}$ & 24.25 \\
\texttt{a192\_dt002\_k125} & $1/192$ & .002 & .125 & .170 & $1.252\,10^{-3}$ & .35058 & $1.650\,10^{-4}$ & 7.59 \\
\texttt{a256\_dt002\_k030} & $1/256$ & .002 & .030 & .111 & $5.385\,10^{-4}$ & .17229 & $3.892\,10^{-5}$ & 13.84 \\
\texttt{a256\_dt002\_k125} & $1/256$ & .002 & .125 & .227 & $4.891\,10^{-4}$ & .35460 & $1.626\,10^{-4}$ & 3.01 \\
\texttt{a256\_dt004\_k125} & $1/256$ & .004 & .125 & .454 & $3.113\,10^{-4}$ & .35474 & $3.061\,10^{-4}$ & 1.02 \\
\texttt{a384\_dt002\_k125} & $1/384$ & .002 & .125 & .340 & $2.974\,10^{-4}$ & .35425 & $1.616\,10^{-4}$ & 1.84 \\
\bottomrule
\end{tabular}
\end{center}

Pour $U=0.15$ et $L=0.24$, les indicateurs de sélection étaient :

\begin{center}
\small
\begin{tabular}{lrrr}
\toprule
Cas & $Re(U=0.15)$ & $Ma(U=0.15)$ & $Re(Ma=0.3)$ \\
\midrule
\texttt{a192\_dt002\_k030} & 35.13 & .880 & 11.98 \\
\texttt{a192\_dt002\_k125} & 28.75 & .428 & 20.16 \\
\texttt{a256\_dt002\_k030} & 66.85 & .871 & 23.03 \\
\texttt{a256\_dt002\_k125} & 73.60 & .423 & 52.20 \\
\texttt{a256\_dt004\_k125} & 115.64 & .423 & 82.05 \\
\texttt{a384\_dt002\_k125} & 121.05 & .423 & 85.76 \\
\bottomrule
\end{tabular}
\end{center}

Les deux fluides à $k_BT=0.03$ restent trop compressibles à $U=0.15$. Les variantes \texttt{a256\_dt004\_k125} et \texttt{a384\_dt002\_k125} permettent un Reynolds plus élevé à Mach donné, mais leur $Sc$ proche de l'unité ou de deux et leur $\lambda/a$ élevé les rapprochent d'un régime cinétique/gazeux. Le cas \texttt{a256\_dt002\_k125} a été retenu comme compromis : libre parcours appréciable sans être dominant, $Sc$ de quelques unités, vitesse du son élevée et résolution de l'obstacle améliorée. Les valeurs du balayage ne doivent toutefois pas être prises comme valeurs finales de viscosité, car elles ne reposent que sur une réalisation TG.

\subsection{Contrôle de taille de boîte et ensemble multi-graines}

Le candidat a été recalibré avec la même taille de cellule $a=1/256$ dans deux boîtes :
\begin{equation}
\begin{aligned}
  64^2 &: L_x=L_y=0.25,\quad (n_x,n_y)_{TG}=(1,1),\quad n_{s}=2,\\
  128^2 &: L_x=L_y=0.50,\quad (n_x,n_y)_{TG}=(2,2),\quad n_{s}=4.
\end{aligned}
\end{equation}
Les longueurs d'onde physiques sont ainsi identiques : $64$ cellules pour Taylor--Green et $32$ cellules pour le son. Les durées et cadences ont été appariées : $2000$ pas et $50$ dumps demandés pour TG, $1200$ pas et $120$ dumps demandés pour chacune des quatre réalisations acoustiques, $3000$ pas et $60$ dumps demandés pour le MSD complet. Quatre graines de base ont été utilisées : $493201$, $593201$, $693201$ et $793201$.

Les viscosités ajustées séparément étaient beaucoup plus dispersées en $64^2$ qu'en $128^2$ :
\begin{equation}
  \nu_{64,\mathrm{ind}}=(5.880\pm1.186)\times10^{-4},
  \qquad
  \nu_{128,\mathrm{ind}}=(6.051\pm0.320)\times10^{-4},
\end{equation}
où les incertitudes sont les écarts-types entre graines. Cette différence de dispersion provient du meilleur rapport signal/bruit de la projection TG dans la grande boîte, qui contient quatre fois plus de particules.

La méthode de référence consiste à moyenner les amplitudes TG signées avant la régression. Sur la fenêtre commune $0.08\leq t\leq1.36$, on obtient
\begin{equation}
  \nu_{64,\mathrm{ens}}=5.8526\times10^{-4},
  \qquad R^2=0.99870,
\end{equation}
avec une sensibilité jackknife $3.71\times10^{-5}$, et
\begin{equation}
  \nu_{128,\mathrm{ens}}=5.9751\times10^{-4},
  \qquad R^2=0.99972,
\end{equation}
avec une sensibilité jackknife $1.29\times10^{-5}$. L'écart entre les deux boîtes n'est que $2.09\,\%$ et reste inférieur à l'incertitude de l'ensemble $64^2$. Aucun effet significatif de taille de boîte n'est donc détecté à taille de cellule et nombre d'onde physique fixés. La boîte $128^2$ est néanmoins recommandée pour une référence précise, car elle réduit fortement la variance de l'estimateur TG.

Sur les quatre calibrations $128^2$, les autres propriétés donnent
\begin{equation}
  c_s=0.35459\pm0.00050,
  \qquad
  D_{\mathrm{self}}=(1.6588\pm0.0268)\times10^{-4},
\end{equation}
les incertitudes étant les écarts-types inter-graines. En combinant $D_{\mathrm{self}}$ avec la viscosité du signal d'ensemble,
\begin{equation}
  \boxed{\nu=5.98\times10^{-4}},\qquad
  \boxed{c_s=0.35459},\qquad
  \boxed{D_{\mathrm{self}}=1.659\times10^{-4}},\qquad
  \boxed{Sc\simeq3.60}.
\end{equation}

Pour $L=0.24$ et $U=0.1064$,
\begin{equation}
  Ma=0.3001,\qquad Re=42.74,
\end{equation}
et
\begin{equation}
  \frac{c_sL}{\nu}=142.4.
\end{equation}
Atteindre $Re=50$ avec ce fluide demande
\begin{equation}
  U\simeq0.1245,\qquad Ma\simeq0.351.
\end{equation}
Le compromis physiquement documenté est donc $Re\simeq43$ à $Ma\simeq0.30$, ou $Re=50$ au prix d'un Mach voisin de $0.35$. Ces valeurs remplacent les premières estimations fondées sur une seule trajectoire TG.

\subsection{Recommandations pour les domaines et modules à traiter}

\paragraph{Écoulements périodiques et benchmarks bulk.}
La calibration du SRC sous-jacent doit être effectuée avant toute comparaison Q6/Q9. Taylor--Green, onde longitudinale et MSD fournissent la référence à laquelle les variantes projetées doivent être comparées. Le resampling doit rester désactivé lorsqu'aucun défaut de support n'est observé. Une boîte $128^2$ et plusieurs graines sont préférables lorsque le Reynolds doit être fixé avec précision.

\paragraph{Obstacles, sillages et pénalisation Darcy/$\chi$.}
La vitesse d'entrée doit d'abord être choisie à partir de $c_s$ et de $\nu$, et non imposée indépendamment. Les sillages du cas Darcy historique doivent être rejoués avec le fluide calibré \texttt{a256\_dt002\_k125} avant de conclure à un défaut de repeuplement. Il faut distinguer la pénalisation poreuse volumique d'une frontière solide et vérifier séparément fuite normale, production de vorticité, déficit de vitesse et bilan de masse. Une absence de sillage à $Re\simeq7$ et $Ma>3$ ne constitue pas un test pertinent d'un mécanisme de resampling destiné à un écoulement faiblement compressible.

\paragraph{Parois et cellules coupées.}
Les particules virtuelles peuvent compléter les moments de collision dans les fractions solides, mais elles ne doivent pas être comptées comme masse physique. Dans une cellule coupée, la cible de masse doit être liée à la fraction accessible et aux conditions de flux, non à une occupation bulk uniforme. Les cellules totalement solides ne doivent pas conserver une population fluide persistante. Toute réparation de support doit préserver masse, quantité de mouvement et énergie thermique relatives dans la partie fluide.

\paragraph{Ouvertures et domaines segmentés.}
Les distributions injectées doivent être compatibles avec le $k_BT$, la masse, le débit et le Mach calibrés. Avant tout split ou transfert, il faut auditer la suppression aux sorties, la réinjection, le bilan de masse par face et la capacité du transport thermique à alimenter les zones transverses. Un dépeuplement créé par une condition limite incohérente ne doit pas être compensé silencieusement par le resampling.

\paragraph{Resampling.}
Le resampling est un contrôle de représentation, pas une équation d'état ni une fermeture incompressible. Il ne doit être activé qu'après démonstration qu'un support statistique insuffisant subsiste dans un régime physique correctement calibré. Les critères doivent être géométriques et statistiques, dépendre de la population effective et de la fraction fluide, rester paramétrables, et conserver exactement masse, moment, énergie thermique et espèces. Le split seul peut supprimer les cellules pauvres mais provoquer une cascade de masses faibles et un coût élevé ; les résultats 0493o doivent donc être requalifiés sur les nouveaux régimes avant promotion en chemin de production.

\paragraph{Fermeture quasi incompressible Q6/Q9 et viriel.}
Une projection ne doit pas servir à masquer un Mach intrinsèquement trop élevé. La première action consiste à choisir $U$, $a$, $\Delta t$ et $k_BT$ afin d'obtenir un domaine $Re$--$Ma$ accessible. Q6/Q9 doit ensuite être évalué comme correction du fluide calibré : réduction des modes compressifs, conservation du transport transverse et absence de modification indue de $\nu$, $c_s$ ou $D_{\mathrm{self}}$. La fermeture virielle doit rester un modèle EOS explicite, distinct du maintien du support particulaire.

\paragraph{Mélanges multi-espèces.}
La calibration doit être répétée pour toute modification de masses, températures, occupations, règles de collision ou thermostat par espèce. Les grandeurs barycentriques ne suffisent pas nécessairement : les diffusivités propres et croisées, le Schmidt de chaque composant et la réponse acoustique du mélange peuvent devenir déterminants. Q6 barycentrique et resampling conservatif par type doivent être qualifiés relativement à cette référence multi-espèces.

\subsection{Contrat cible pour une référence de co-simulation}

Le runner 0493w1 est actuellement un outil de cas unique fonctionnel. À terme, il doit devenir une interface stable du code : l'utilisateur fournit la configuration SRC et les échelles de son domaine, puis reçoit un artefact de référence utilisable par le solveur, un outil externe ou une boucle de co-simulation. Le contrat cible devrait comprendre :
\begin{itemize}[leftmargin=*]
  \item un manifeste exhaustif des paramètres qui définissent le fluide, du domaine de calibration, des modes, cadences et graines ;
  \item l'identifiant du commit, le hash du binaire et un hash canonique de la configuration ;
  \item les statuts indépendants \texttt{viscosityStatus}, \texttt{soundStatus} et \texttt{diffusionStatus}, accompagnés des résidus, $R^2$ et incertitudes ;
  \item un mode ensemble qui moyenne automatiquement les signaux TG et acoustiques avant régression ;
  \item deux profils explicites, \texttt{cell\_equivalent} et \texttt{full\_domain}, sans ambiguïté sur la taille de cellule et les longueurs d'onde ;
  \item des recommandations machine-readable : $U(Ma_{\mathrm{cible}})$, $Re(Ma_{\mathrm{cible}})$, $U(Re_{\mathrm{cible}})$, $Ma(Re_{\mathrm{cible}})$ et avertissements de régime cinétique ;
  \item une politique d'archivage non destructive empêchant l'écrasement silencieux d'une référence existante.
\end{itemize}

Les sorties canoniques actuelles sont
\begin{verbatim}
analysis/fluid_calibration_0493w1.json
analysis/fluid_calibration_0493w1.csv
analysis/README_0493w1_RESULTS.md
analysis/tg_decay_0493w1.csv
analysis/sound_mode_0493w1.csv
analysis/msd_0493w1.csv
\end{verbatim}
Le JSON constitue l'interface principale pour la co-simulation ; les trois séries CSV permettent de recalculer les fits, de construire un ensemble multi-graines et d'auditer les diagnostics sans conserver les dumps particulaires volumineux.

Un appel représentatif du candidat de référence s'écrit :
\begin{verbatim}
LIVE_PROGRESS=1 ALLOW_LARGE_DUMPS=1 \
RUN_ROOT=runs/calibrations/a256_dt002_k125_size128 \
Lx=0.5 Ly=0.5 NX=128 NY=128 GAMMA=20 \
DT=0.002 KBT=0.125 PARTICLE_MASS=1 \
ROTATION_ANGLE=1.5707963267948966 \
RANDOM_ROTATION_SIGN=true GRID_SHIFT_ENABLE=true \
THERMOSTAT_ENABLE=true THERMOSTAT_MODE=cell_relative_rescale \
THERMOSTAT_EVERY=1 THERMOSTAT_TARGET_KBT=0.125 \
THERMOSTAT_MIN_PARTICLES=3 \
TG_MODE_X=2 TG_MODE_Y=2 TG_STEPS=2000 TG_DUMP_COUNT=50 \
SOUND_MODE_X=4 SOUND_REPLICATES=4 \
SOUND_STEPS=1200 SOUND_DUMP_COUNT=120 \
MSD_GRID_MODE=full MSD_STEPS=3000 MSD_DUMP_COUNT=60 \
CHARACTERISTIC_U=0.1064 CHARACTERISTIC_L=0.24 \
bash scripts/run_0493w1_src_fluid_calibrator.sh
\end{verbatim}

Cette étape d'étalonnage ne rend pas inutiles les développements de support, de projection ou de fermeture. Elle fournit le référentiel qui permet enfin de décider lesquels sont nécessaires, dans quelles zones, et avec quel rapport bénéfice/coût. Les validations futures doivent donc être organisées en deux niveaux : d'abord le fluide SRC calibré et ses conditions limites, puis l'effet incrémental de Q6/Q9, du resampling, du viriel ou des traitements $\chi$/paroi.


\section{Conclusion}


Le dernier développement GPU ajoute un jalon distinct. La branche \texttt{SRC\_GPU} dispose maintenant d'un chemin \textbf{SRC/MPCD classic full CUDA résident}. Ici, \emph{SRC classic} signifie explicitement advection/streaming, décalage aléatoire de grille, collision/rotation SRC et thermostat. Les validations consolidées couvrent les familles périodique, wall-simple, solide rectangle, solide circulaire, piston/mobile wall, inlet/outlet full-face, inlet/outlet segmenté, et cercle avec inlet/outlet pour le cas Von Karman. Les ouvertures CUDA disposent en outre de régimes explicites \texttt{neumann}, \texttt{equilibrium\_flux} et \texttt{forced\_flux}, utilisables en full-face ou segmenté dans le périmètre SRC classic-only. La fermeture liquide -- Q6, resampling et viriel -- reste séparée et préservée pour les chantiers futurs, en particulier Q6 CUDA et resampling CUDA complet.

Les développements post-0143 ajoutent à cette hiérarchie un autre résultat important. Les conditions limites segmentées permettent désormais plusieurs inlets/outlets sur une même face solide, ce qui généralise fortement les cas de cuve, injection localisée ou injection multi-type marquée. La suppression de \texttt{method} rend les ablations propres : Q6 et resampling sont des briques indépendantes, activées par \texttt{projectionEnable} et \texttt{resamplingEnable}. Enfin, la réponse de capacité fermée fournit une voie continue entre régime quasi-incompressible et régime faiblement compressible : une surmasse globale affaiblit Q6, augmente le viriel effectif, modifie la cible de remap pour conserver la compression, puis projette \(P_{\rm kin}+P_{\rm vir}\) en charge pariétale. La validation statique de boîte uniformément surcompressée montre que la pression pariétale est uniforme par face et que la force nette est nulle, ce qui constitue un jalon majeur pour un futur couplage à des parois déformables ou rompables.



Le dernier cycle CUDA ajoute maintenant une réponse partielle mais validée à cette difficulté. Le resampling post-SRC CUDA ne modifie pas la collision SRC ; il reconditionne la représentation particulaire après que le pas physique a produit l'état fluide. Les modules 0295--0299 établissent une chaîne survey, reconditionnement de masse, split/merge local, restauration $M/P/K_{\mathrm{rel}}$ et filtrage boundary-aware. Les campagnes 0300--0302 montrent que le backward step, qui perd jusqu'à $191$ cellules humides vides à $U_{\mathrm{in}}=0.60$ en SRC classic, n'en conserve plus aucune avec le réglage nominal \texttt{12:20:32}, alors que les erreurs locales de conservation restent au niveau du roundoff. Ce résultat donne une première validation physique du resampling CUDA comme contrôle de support particulaire, distinct de Q6 et complémentaire du chemin SRC classic résident.


Le cycle 0490a--0490p généralise maintenant cette chaîne aux mélanges. La conservation est imposée par \texttt{type}, depuis le split/merge local jusqu'au refill, au plan donneur--receveur et à la fermeture de masse. Le chemin CUDA résident élimine le plan CPU, les vecteurs d'opérations hôte, les copies complètes de rollback, les scans de particules pour les dépôts et le pool, puis le dernier miroir CPU wet/poor/rich. La validation à 10000 pas et l'audit \texttt{remaining\_cpu\_scope=none} constituent un jalon d'architecture : le resampling multi-espèces peut désormais rester sur l'état CUDA partagé dans le domaine périodique qualifié. La prochaine difficulté n'est plus la résidence du mécanisme, mais son élargissement physique aux frontières, aux solides et à une éventuelle projection Q6 dépendante de la phase.

Le projet SRC/MPCD quasi-incompressible aboutit maintenant à une hiérarchie plus claire des méthodes. La redistribution dure et la fermeture virielle historique ont joué un rôle exploratoire, mais ne constituent plus le socle principal de validation. La génération Q6/Q9 a fourni une formulation plus robuste de l'incompressibilité : collision SRC classique, projection de vitesse, correction basse fréquence du flux de masse, thermostat corrigé, diagnostics de conservation et traitement explicite des parois par particules virtuelles.

La difficulté décisive n'était toutefois pas seulement la divergence du champ de vitesse ou du flux de masse. Les cas de poches vides, de remplissage et d'inlet/outlet ont montré que la méthode doit garantir un support particulaire local suffisant. Une cellule pauvre ne doit pas être traitée comme un puits de pression ou une anomalie eulérienne à corriger par force ; elle doit être traitée comme un défaut de représentation particulaire locale.

La formulation SRC/MPCD pondérée avec resampling fournit cette réponse. Dans sa forme stabilisée, elle combine :
\begin{equation}
  \text{rôles particulaires}
  + \text{pool préalloué}
  + \text{garde }N_{\min}/N_{\mathrm{target}}/N_{\max}
  + \text{masses pondérées}
  + \text{remap }M,P,E
  + \text{Q6}.
\end{equation}
Le point essentiel est l'ordre logique : on maintient d'abord la population active des cellules mouillées, puis on impose la masse et le moment par remap pondéré. La correction de masse seule est insuffisante, car elle peut masquer un dépeuplement local par des particules trop lourdes.

Le portage C++/OpenMP constitue une étape majeure : il démontre que la méthode peut être organisée dans une architecture compatible avec une exécution parallèle, des états persistants \texttt{.smpcd}, des scripts MATLAB de préparation et d'analyse, et des scripts \texttt{bash} de lancement reproductibles. Les validations récentes indiquent que Taylor--Green, Poiseuille \texttt{wallVP}, le cylindre périodique, la marche, l'obstacle ouvert et le remplissage d'un domaine initialement inactif se comportent de façon cohérente. Le cas Poiseuille montre que la physique de canal est préservée ; le cas injection/fill montre que la couche wet/dry et le pool actif/inactif sont opérationnels ; le cas obstacle ouvert a joué un rôle utile en révélant la nécessité du garde de population maintenant réintroduit.

Le statut actuel peut donc être résumé ainsi : le socle Q6/Q9 et les solides fixes sont validés comme base, le resampling pondéré est devenu la voie principale pour les cellules pauvres et domaines évolutifs, et le portage OpenMP post-0140 doit maintenant être revalidé systématiquement. Les priorités immédiates sont de confirmer la bande de population $N_c\in[14,26]$ sur tous les cas mouillés, de balayer la cadence de remap de masse, puis de reprendre les validations de sillage. Les optimisations de Q6 et l'extension aux cylindres inlet/outlet doivent être traitées dans des chantiers séparés, afin de ne pas compromettre la validation physique du mécanisme principal.


Le complément Darcy/chiVP modifie également le statut des géométries complexes. Le champ $\chi$ peut maintenant être traité comme une donnée géométrique externe de même rang pratique que l'état particulaire initial : il décrit la forme, désactive les particules initiales dans le solide, pilote le frein Darcy moyen, oriente le bain de réémission et fournit les moments virtuels de collision d'interface. La correction 0426 des fastflags CUDA résidents montre que cette généralité n'impose pas de pénalité majeure de performance lorsque les scripts sélectionnent le backend résident optimisé. Cette observation est importante pour les futurs chantiers d'optimisation topologique, car elle rend réaliste une boucle externe qui modifie $\chi$ sans recompilation du solveur.

La conclusion méthodologique est forte : rendre SRC/MPCD quasi-incompressible ne consiste pas seulement à projeter un champ de vitesse. La réussite principale est de coupler cette projection à un contrôle du support particulaire. C'est cette combinaison qui ouvre la voie à des écoulements ouverts, des fronts de remplissage, des solides immergés et, à terme, des géométries complexes avec un coût parallèle maîtrisable.




\begin{thebibliography}{99}

\bibitem{MalevanetsKapral1999}
A. Malevanets and R. Kapral,
``Mesoscopic model for solvent dynamics,''
\emph{The Journal of Chemical Physics}, vol. 110, no. 17, pp. 8605--8613, 1999.
doi: \href{https://doi.org/10.1063/1.478857}{10.1063/1.478857}.

\bibitem{GompperIhleKrollWinkler2009}
G. Gompper, T. Ihle, D. M. Kroll, and R. G. Winkler,
``Multi-Particle Collision Dynamics: A Particle-Based Mesoscale Simulation Approach to the Hydrodynamics of Complex Fluids,''
in \emph{Advanced Computer Simulation Approaches for Soft Matter Sciences III},
Advances in Polymer Science, vol. 221, Springer, Berlin, Heidelberg, pp. 1--87, 2009.
See also arXiv:0808.2157.

\bibitem{IhleTuzelKroll2006}
T. Ihle, E. Tüzel, and D. M. Kroll,
``Consistent particle-based algorithm with a non-ideal equation of state,''
\emph{Europhysics Letters}, vol. 73, no. 5, pp. 664--670, 2006.

\bibitem{ZantopStark2021a}
A. W. Zantop and H. Stark,
``Multi-particle collision dynamics with a non-ideal equation of state. I,''
\emph{The Journal of Chemical Physics}, vol. 154, 024105, 2021.

\bibitem{ZantopStark2021b}
A. W. Zantop and H. Stark,
``Multi-particle collision dynamics with a non-ideal equation of state. II,''
\emph{The Journal of Chemical Physics}, vol. 155, 134904, 2021.

\bibitem{Chorin1968}
A. J. Chorin,
``Numerical solution of the Navier--Stokes equations,''
\emph{Mathematics of Computation}, vol. 22, no. 104, pp. 745--762, 1968.

\bibitem{SolenthalerPajarola2009}
B. Solenthaler and R. Pajarola,
``Predictive-corrective incompressible SPH,''
\emph{ACM Transactions on Graphics}, vol. 28, no. 3, article 40, 2009.

\bibitem{DeCourcy2024}
J. J. De Courcy, T. C. S. Rendall, L. Constantin, B. Titurus, and J. E. Cooper,
``Incompressible $\delta$-SPH via artificial compressibility,''
\emph{Computer Methods in Applied Mechanics and Engineering}, vol. 420, article 116700, 2024.

\bibitem{Yang2021}
K. Yang, T. Aoki, and co-authors,
``Weakly compressible Navier--Stokes solver based on evolving pressure projection method for two-phase flow simulations,''
\emph{Journal of Computational Physics}, vol. 431, article 110113, 2021.

\bibitem{WinklerHuang2009}
R. G. Winkler and C.-C. Huang,
``Stress tensors of multiparticle collision dynamics fluids,''
\emph{The Journal of Chemical Physics}, vol. 130, 074907, 2009.

\bibitem{ImperioPaddingBriels2011}
A. Imperio, J. T. Padding, and W. J. Briels,
``Force calculation on walls and embedded particles in multi-particle collision dynamics simulations,''
\emph{Physical Review E}, vol. 83, 046704, 2011.

\bibitem{ValizadehMonaghan2015}
A. Valizadeh and J. J. Monaghan,
``A study of solid wall models for weakly compressible SPH,''
\emph{Journal of Computational Physics}, vol. 300, pp. 5--19, 2015.

\bibitem{Islam2024}
M. R. I. Islam,
``SPH-based framework for modelling fluid--structure interaction problems with finite deformation and fracturing,''
\emph{Ocean Engineering}, vol. 294, article 116722, 2024.

\bibitem{Singh2025}
V. Singh, C. Peng, and M. R. I. Islam,
``Three-dimensional SPH modeling of brittle fracture under hydrodynamic loading,''
\emph{Computer Methods in Applied Mechanics and Engineering}, 2025.

\end{thebibliography}

\end{document}
